%% Factorization Swiss-Cheese: the primary treatment
%% Layers 0--5 of the six-layer hierarchy

%% Macro safety: providecommand only
\providecommand{\fSet}{\mathsf{fSet}}
\providecommand{\IndCoh}{\operatorname{IndCoh}}
\providecommand{\QCoh}{\operatorname{QCoh}}
\providecommand{\Dmod}{D\text{-}\mathsf{mod}}
\providecommand{\Mbar}{\overline{\cM}}
\providecommand{\colim}{\operatorname{colim}}
\providecommand{\varprojlim}{\underleftarrow{\lim}}
\providecommand{\Obs}{\mathsf{Obs}}
\providecommand{\id}{\mathrm{id}}
\providecommand{\cG}{\mathcal{G}}
\providecommand{\cB}{\mathcal{B}}
\providecommand{\bD}{\mathbb{D}}
\providecommand{\op}{\mathrm{op}}
% Vol III two-stage factorisation (cross-volume macros)
\providecommand{\PhiFA}{\Phi^{\mathrm{FA}}}
\providecommand{\SpCh}{\mathrm{Sp}_{\mathrm{Ch}}}
\providecommand{\HolFA}{\mathrm{HolFA}}
\providecommand{\EdHolFA}{E_d\text{-}\HolFA}
\providecommand{\EnHolFA}{E_n\text{-}\HolFA}
\providecommand{\intSigma}{\textstyle\int_{\Sigma}}
\providecommand{\hCS}{\mathrm{hCS}}

\chapter{Factorization Swiss-Cheese Algebras}
\label{ch:factorization-swiss-cheese}

%%%%%%%%%%%%%%%%%%%%%%%%%%%%%%%%%%%%%%%%%%%%%%%%%%%%%%%%%%%%%%%%%%%%%%%%%%%%%
\section{The factorization substrate}
\label{sec:factorization-substrate}
%%%%%%%%%%%%%%%%%%%%%%%%%%%%%%%%%%%%%%%%%%%%%%%%%%%%%%%%%%%%%%%%%%%%%%%%%%%%%

A chiral algebra is not an algebra over an operad. It is a
factorizable $D$-module on the Ran space. The operad
$\SCchtop$ is what remains after formal completion at collision
points: it is the local shadow, not the governing structure.
Formal completion is a functor
\[
 \operatorname{Loc}_x\colon
 \Fact_{\mathrm{BD}}(X)
 \longrightarrow
 \operatorname{Alg}_{C_\ast(\FM_\bullet(T_xX))}
\]
from global factorization data to the collision model at~$x$; it is
not an inverse construction. At genus~$0$ on the affine chart
$\mathbb{P}^1\smallsetminus\{\infty\}\simeq\C$, it identifies the
translation-invariant, unramified local subcategory with its
formal-disc model. Extending an arbitrary local object across
$\infty$ is a separate descent problem. At genus~$g \geq 1$ the
forgetful nature of $\operatorname{Loc}_x$ is unavoidable: the periods
of~$\Sigma_g$, the family monodromy, and the factorization pattern on
$\Ran(\Sigma_g)$ carry information invisible to the local operad. On
the scalar bar lane this global excess is measured by the anomaly
$\dfib^{\,2} = \kappa(\cA) \cdot \omega_g$.

All results in this chapter apply to logarithmic $\SCchtop$-algebras (Definition~\ref{def:log-SC-algebra}). Beyond the affine and $\cW$-algebra lineages, the question of which chiral algebras admit a logarithmic $\SCchtop$-algebra structure remains open.

The Swiss-cheese structure is constructed directly from
factorization data; the local operadic model follows as a
consequence. The treatment is parallel: Beilinson--Drinfeld
(algebraic, $D$-module) and Costello--Gwilliam (differential-geometric,
prefactorization algebra) reach the same object through complementary
routes, and their equivalence is proved at the end of the chapter.

\begin{definition}[Ran space and factorization pattern]
\label{def:ran-space}
\index{Ran space|textbf}
Let $X$ be a smooth curve over a field~$k$ of characteristic~$0$.
\begin{enumerate}[label=\textup{(\roman*)}]
\item The \emph{Ran space} of~$X$ is the (pre)sheaf
\begin{equation}\label{eq:ran-colim}
 \Ran(X)
 \;=\;
 \colim_{S \in \fSet^{\mathrm{surj}}} X^S,
\end{equation}
where the colimit is taken over the category $\fSet^{\mathrm{surj}}$ of
non-empty finite sets with surjections. Points of $\Ran(X)$ are finite
non-empty subsets $\{x_1,\ldots,x_n\} \subset X$.

\item For a partition $I = I_1 \sqcup I_2$ of a finite set, the
\emph{disjointness locus}
\[
 U_{I_1,I_2}
 \;:=\;
 \bigl\{(S_1,S_2) \in \Ran(X) \times \Ran(X) :
 S_1 \cap S_2 = \varnothing\bigr\}
\]
embeds into $\Ran(X)$ via union. The \emph{factorization structure} is
encoded by the functoriality of the assignment $S \mapsto F(S)$ with
respect to partitions of the index set.
\end{enumerate}
\end{definition}

\begin{definition}[Factorizable $D$-module]
\label{def:factorizable-dmod}
\index{factorizable D-module@factorizable $D$-module|textbf}
A \emph{factorizable $D$-module} on $\Ran(X)$ is a $D$-module~$F$
on $\Ran(X)$ equipped with \emph{factorization isomorphisms}: for
every decomposition of a finite set $I = I_1 \sqcup I_2$ into
non-empty disjoint subsets, an isomorphism
\begin{equation}\label{eq:fact-iso}
 F\big|_{U_{I_1,I_2}}
 \;\xrightarrow{\;\sim\;}
 F\big|_{U_{I_1}} \boxtimes F\big|_{U_{I_2}}
\end{equation}
on the disjointness locus, satisfying the associativity and
commutativity constraints inherited from the operadic composition for iterated partitions
$I = I_1 \sqcup \cdots \sqcup I_r$. In the language of
\cite{BD04}, this is a factorization algebra in the category of
$D$-modules on~$X$.
\end{definition}

\begin{definition}[Chiral operations from factorization]
\label{def:chiral-ops-fact}
\index{chiral operations!from factorization}
Let $F$ be a factorizable $D$-module on $\Ran(X)$. The
\emph{chiral $n$-ary operations} are \emph{not} maps between
fibers at points. They are diagonal cospecialisation morphisms. For
the collision of the $i$th and $j$th variables, write
$j_{ij}\colon X^n\smallsetminus\Delta_{ij}\hookrightarrow X^n$ and
$\Delta_{ij}\colon X^{n-1}\hookrightarrow X^n$. The operation is a
$D$-module morphism
\begin{equation}\label{eq:chiral-op-diagonal}
 \mu_{ij}\colon
 j_{ij,*}j_{ij}^{*}\bigl(F^{\boxtimes n}\bigr)
 \;\longrightarrow\;
 \Delta_{ij,*}\bigl(F^{\boxtimes(n-1)}\bigr)
\end{equation}
from the complement of the diagonal to the diagonal stratum
$\Delta_{ij}=\{z_i=z_j\}\subset X^n$. The composition law arises from
iterating such cospecialisations along nested diagonals together with
the factorization isomorphisms on disjointness loci. Operadic
composition composes operations between \emph{fibers}; factorization
composition composes \emph{$D$-module morphisms on configuration
spaces}.
\end{definition}

%%%%%%%%%%%%%%%%%%%%%%%%%%%%%%%%%%%%%%%%%%%%%%%%%%%%%%%%%%%%%%%%%%%%%%%%%%%%%
%% Verdier duality on Ran(X): global definition + well-definedness
%%%%%%%%%%%%%%%%%%%%%%%%%%%%%%%%%%%%%%%%%%%%%%%%%%%%%%%%%%%%%%%%%%%%%%%%%%%%%

\begin{definition}[Verdier duality on $\Ran(X)$]
\label{def:ran-verdier-duality}
\index{Verdier duality!on Ran space|textbf}
Write $\Ran_{\leq n}(X) \subset \Ran(X)$ for the closed sub-ind-scheme of
configurations of cardinality $\leq n$, so that
$\Ran(X) = \colim_n \Ran_{\leq n}(X)$ as ind-schemes of ind-finite type
(\cite{BD04}, 3.4.1; \cite{FG11}, 2.2). For each $n$, Verdier duality
on $D$-modules of ind-finite type gives a contravariant auto-equivalence
$\bD_{\leq n}\colon \Dmod(\Ran_{\leq n}(X))^{\op}
\xrightarrow{\sim} \Dmod(\Ran_{\leq n}(X))$.
The \emph{Verdier duality functor on $\Ran(X)$} is the colimit
\begin{equation}\label{eq:ran-verdier}
 \bD_{\Ran}\;:=\;
 \colim_n\, \bD_{\leq n}
 \colon \Dmod(\Ran(X))^{\op}
 \xrightarrow{\;\sim\;}
 \Dmod(\Ran(X)),
\end{equation}
taken along the system of closed embeddings
$i_{n,n+1}\colon \Ran_{\leq n}(X) \hookrightarrow \Ran_{\leq n+1}(X)$
and the compatibility $\bD_{\leq n+1}\,(i_{n,n+1})_{!} \simeq
(i_{n,n+1})_{\ast}\,\bD_{\leq n}$ (proper base change for the closed
immersion $i_{n,n+1}$, which is proper of ind-finite type).
\end{definition}

\begin{proposition}[Well-definedness of $\bD_{\Ran}$;
\ClaimStatusProvedHere]
\label{prop:ran-verdier-well-defined}
The functor $\bD_{\Ran}$ of Definition~\ref{def:ran-verdier-duality} is
well-defined on $\Dmod(\Ran(X))$ and satisfies:
\begin{enumerate}[label=\textup{(\roman*)}]
\item \emph{Stratumwise compatibility.} For each $n$, the restriction
 $\bD_{\Ran}|_{\Ran_{\leq n}} = \bD_{\leq n}$ agrees with classical
 Verdier duality on the ind-finite-type stratum.
\item \emph{Involutivity.} $\bD_{\Ran}^{2} \simeq \id$ as auto-equivalences.
\item \emph{Chiral-tensor compatibility.} On the disjointness locus
 $U_{I_1,I_2} \subset \Ran(X) \times \Ran(X)$,
 $\bD_{\Ran}(F \otimes^{\mathrm{ch}} G)|_{U_{I_1,I_2}} \simeq
 \bD_{\Ran}(F) \otimes^{\mathrm{ch}} \bD_{\Ran}(G)|_{U_{I_1,I_2}}$
 up to a Tate twist of weight $2 \dim X = 2$ on each $n$-point stratum
 (normalised so that the dualising sheaf of $X^n$ is $\omega_{X^n}[n]$).
\item \emph{Factorization preservation.} If $F$ is a factorizable
 $D$-module in the sense of Definition~\ref{def:factorizable-dmod},
 then so is $\bD_{\Ran}(F)$; the factorization isomorphism for
 $\bD_{\Ran}(F)$ is the Verdier dual of the factorization isomorphism
 for $F$.
\end{enumerate}
\end{proposition}

\begin{proof}
Well-definedness reduces to checking that the colimit
in~\eqref{eq:ran-verdier} exists in the $(\infty,2)$-category
$\Dmod(\Ran(X))$ of Francis--Gaitsgory (\cite{GR17}, IV.5.3). For each
closed immersion $i_{n,n+1}$, the $!$-pullback
$i_{n,n+1}^{!}\colon \Dmod(\Ran_{\leq n+1}) \to \Dmod(\Ran_{\leq n})$
admits a left adjoint $(i_{n,n+1})_{!}$ (the $*$-pushforward on
$D$-modules equals the $!$-pushforward for closed embeddings), and
Verdier duality intertwines these via
$\bD_{\leq n+1}\,(i_{n,n+1})_{!} \simeq (i_{n,n+1})_{\ast}\,\bD_{\leq n}$
(\cite{BD04}, 7.3.6). Passing to the colimit, $\bD_{\Ran}$ is the
unique extension making the squares commute; this is
\cite{GR17}, Volume~I, Chapter~4, Proposition~2.3.5 (Verdier duality
commutes with colimits in ind-schemes of ind-finite type).

(i) is the construction. (ii) follows from involutivity on each
stratum and compatibility with the transition maps. (iii) is the
chiral-tensor analogue of the standard external-tensor formula
$\bD(F \boxtimes G) \simeq \bD(F) \boxtimes \bD(G)[2\dim X]$
(\cite{KS90}, 3.4.3) restricted to the disjointness locus, where
$\otimes^{\mathrm{ch}}$ reduces to $\boxtimes$ by
\cite{BD04}, 3.4.10. (iv) follows by applying $\bD_{\Ran}$ to the
factorization isomorphism~\eqref{eq:fact-iso} and using (iii).
\end{proof}

\begin{remark}[Consequences for Theorem~A]
\label{rem:ran-verdier-thmA}
Proposition~\ref{prop:ran-verdier-well-defined} makes rigorous the
intertwining $\bD_{\Ran}(\barB(\cA)) \simeq \barB(\cA^!)$ invoked in the
introduction and in the Vol~I Theorem~A bar-cobar adjunction: both
sides are factorizable $D$-modules on $\Ran(X)$, and the isomorphism is
the chiral-tensor-compatible Verdier dual at each Ran-stratum, glued by
the colimit structure. In particular, the decorator
\texttt{test\_verdier\_intertwining\_iv} of the independent
verification protocol targets a well-defined global statement, not a
stratum-wise assertion whose colimit compatibility is unchecked.
\end{remark}

\begin{remark}[The local shadow]
\label{rem:local-shadow-extraction}
\index{local shadow|textbf}
The topological operad $\SCchtop$ is extracted from the factorization
structure by three steps:
\begin{enumerate}[label=\textup{(\arabic*)}]
\item Restrict to a formal neighbourhood of a point $x \in X$
 (where $X \simeq \C$ locally).
\item Replace the $D$-module structure by its formal completion, which
 is a cochain complex (a vector space with a differential) rather
 than a quasicoherent sheaf with connection.
\item Identify the formal neighbourhood of the small diagonal in
 $X^n$ with $\FM_n(\C)$, the Fulton--MacPherson compactification.
\end{enumerate}
This extraction loses global information: monodromy around cycles,
the extension data of the $D$-module connection, and period integrals
over homology classes. At genus~$0$ on the affine chart of
$\mathbb{P}^1$ the loss vanishes only for the translation-invariant
local subcategory that is already known to extend across the chosen
point at infinity. At genus~$g \geq 1$ the loss has a scalar
projection measured in the bar construction by
$\kappa(\cA)\cdot\omega_g$; this scalar is not visible in the
formal-disc operad.

Thus the local model $\SCchtop$
(Definition~\ref{def:SC-operations}) is the genus-$0$ shadow of
the factorization structure. The homotopy-Koszulity theorem
(Theorem~\ref{thm:homotopy-Koszul}) is a theorem about this local
model; its promotion to an all-genus statement requires the
factorization-algebraic input of this chapter.
\end{remark}


%%%%%%%%%%%%%%%%%%%%%%%%%%%%%%%%%%%%%%%%%%%%%%%%%%%%%%%%%%%%%%%%%%%%%%%%%%%%%
\section{Worked examples}
\label{sec:worked-examples-factorization}
%%%%%%%%%%%%%%%%%%%%%%%%%%%%%%%%%%%%%%%%%%%%%%%%%%%%%%%%%%%%%%%%%%%%%%%%%%%%%

Before developing the general theory, we compute the factorization
Swiss-cheese structure for the three canonical examples: the free boson
(Heisenberg), the affine Kac--Moody algebra, and the Virasoro
algebra. Each computation is carried out in explicit coordinates,
with all signs and structure maps written out. The general
definitions and constructions follow in
Sections~\ref{sec:BD-factorization-SC}--\ref{sec:factorization-koszul-duality}.

%% =====================================
\subsection{The Heisenberg algebra (free boson)}
\label{subsec:heisenberg-SC}
%% =====================================

\begin{example}[Heisenberg as BD factorization Swiss-cheese algebra]
\label{ex:heisenberg-SC}
\index{Heisenberg algebra!BD factorization Swiss-cheese}
\ClaimStatusProvedHere{}
Let $X$ be a smooth curve. The \emph{Heisenberg chiral algebra} on~$X$
is the symmetric algebra $D$-module
$\cA = \Sym(\omega_X)$ on~$X$, where $\omega_X = \Omega^1_X$ is the
canonical bundle. The chiral product is the $!$-extension of
the exterior tensor product $\omega_X \boxtimes \omega_X$ to the
diagonal, with residue determined by the tautological pairing
$\omega_X \otimes \omega_X^{-1} \to \cO_X$. In coordinates: for
$a = f(z)\,dz$ and $b = g(z)\,dz$ sections of $\omega_X$, the OPE is
\begin{equation}\label{eq:heisenberg-OPE}
 a(z) \cdot b(w)
 \;\sim\;
 \frac{f(z)\,g(w)}{z - w}\, dz \otimes dw
 \;+\; \text{regular},
\end{equation}
with propagator $K(z,w) = dz/(z-w)$ on $\C$.

\smallskip\noindent
\textbf{Factorizable $D$-module on $\Ran(X)$.}
The $I$-th component of the factorizable $D$-module
(Construction~\ref{constr:fact-dmod-from-chiral}) is
\[
 F_{\mathrm{Heis}}^I
 \;=\;
 j_{!*}\bigl(\omega_X^{\boxtimes I}\bigr)
 \big|_{\Delta_{\mathrm{arr}}},
\]
where $j\colon U_I \hookrightarrow X^I$ is the diagonal complement.
For $I = \{1,2\}$, the space $F_{\mathrm{Heis}}^{\{1,2\}}$ consists
of sections of $\omega_X \boxtimes \omega_X$ on $X \times X$
with at most simple poles along $z_1 = z_2$: concretely, expressions
of the form
\[
 \frac{f(z_1,z_2)}{z_1 - z_2}\, dz_1 \otimes dz_2
 \;+\;
 g(z_1,z_2)\, dz_1 \otimes dz_2,
\]
where $f, g$ are regular. The factorization isomorphism on the
disjointness locus $\{z_1 \neq z_2\}$ is the tautological
splitting
$\omega_X \boxtimes \omega_X|_{z_1 \neq z_2} =
\omega_X|_{z_1} \otimes \omega_X|_{z_2}$.

\smallskip\noindent
\textbf{Swiss-cheese structure.}
The BD factorization Swiss-cheese algebra is:
\begin{itemize}
\item \emph{Closed sector}: $F_{\mathsf{cl}} = F_{\mathrm{Heis}}$
 on $\Ran(X)$.
\item \emph{Open sector}: $F_{\mathsf{op}} = \Sym(V)$ viewed as an
 $\Eone$-algebra on $\Ran(\R)$ via ordered multiplication, where
 $V = H^0(D, \omega_X)$ for a formal disc $D \subset X$.
\item \emph{Mixed sector}: the chiral module structure
 $F_{\mathrm{Heis}} \otimes F_{\mathsf{op}} \to F_{\mathsf{op}}$,
 given by the contour-integral formula
 $a(z) \cdot v = \oint_{|z-w|=\varepsilon}
 K(z,w)\, a(z)\, v\, dz$ for $v \in F_{\mathsf{op}}$.
\end{itemize}
The CG prefactorization algebra (Definition~\ref{def:CG-prefact})
associated to this data is the \emph{free boson}: the prefactorization
algebra $\Obs$ on $\C \times \R$ with propagator
$P(z_1,t_1;\,z_2,t_2) = \frac{dz_1}{z_1 - z_2}\cdot\theta(t_1 - t_2)$,
matching Definition~\ref{def:BV-origin} with $\fg = \C$ (abelian)
and $S_{\mathrm{int}} = 0$.
\end{example}

\begin{example}[Heisenberg bar complex at arity $3$]
\label{ex:heisenberg-bar-3}
\index{Heisenberg algebra!bar complex at arity 3}
\ClaimStatusProvedHere{}
Let $\cA = \Sym(\omega_X)$ as above, with a generator
$J = J(z)\,dz \in \omega_X$. The augmentation ideal
$\bar\cA = \cA_+ = \bigoplus_{n \geq 1}\Sym^n(\omega_X)$, and the
desuspended generators are $s^{-1}J$, $s^{-1}J^2$, etc. Working
at tree level (genus~$0$, $X = \C$), the bar complex
$\barB(\cA) = T^c(s^{-1}\bar\cA)$ has the following arity-$3$
elements, restricting to the case where all tensor factors are
the generator~$J$:
\begin{align}
 \xi_1 &= [s^{-1}J \,|\, s^{-1}J \,|\, s^{-1}J],
 \label{eq:heis-bar-el-1}\\
 \xi_2 &= [s^{-1}J^2 \,|\, s^{-1}J],
 \label{eq:heis-bar-el-2}\\
 \xi_3 &= [s^{-1}J \,|\, s^{-1}J^2].
 \label{eq:heis-bar-el-3}
\end{align}
Here $\xi_1$ is an element of bar degree~$3$ (three tensor factors in
the bar construction), while $\xi_2$ and $\xi_3$ have bar degree~$2$
with one factor of higher polynomial degree.

\smallskip\noindent
\textbf{The bar differential $d_{\barB}$.}
The bar differential is the coderivation determined by the chiral
binary product $\mu_2\colon s^{-1}\bar\cA \otimes s^{-1}\bar\cA
\to s^{-1}\bar\cA$. For the Heisenberg algebra, $\mu_2$ acts as:
\begin{itemize}
\item $\mu_2(s^{-1}J \otimes s^{-1}J) = 0$, since the Heisenberg OPE
 $J(z)\,J(w) \sim 1/(z-w)$ has singular part a \emph{scalar}
 (the identity), which lies in $\cA^0 = k$ (the augmentation), not
 in $\bar\cA$. So the projected binary product $\bar\mu_2$ on
 $\bar\cA$ vanishes.
\end{itemize}
Consequently, $d_{\barB}(\xi_1) = 0$: the bar differential on a
degree-$3$ element is
\begin{align*}
 d_{\barB}(\xi_1)
 &= [s^{-1}\bar\mu_2(J,J) \,|\, s^{-1}J]
 + (-1)^{|s^{-1}J|}\,
 [s^{-1}J \,|\, s^{-1}\bar\mu_2(J,J)] \\
 &= [0 \,|\, s^{-1}J] + (-1)^{0} [s^{-1}J \,|\, 0]
 \;=\; 0.
\end{align*}
Here $|s^{-1}J| = |J| - 1 = 0$ (the generator $J$ has cohomological
degree~$1$ as a section of~$\omega_X$, and desuspension subtracts~$1$).

More generally, all higher arity bar elements in the Heisenberg bar
complex are $d_{\barB}$-closed, because the projected binary product
$\bar\mu_2$ vanishes identically. The Heisenberg algebra is
\emph{intrinsically formal}: $\barB(\cA)$ has trivial differential,
and the bar complex is quasi-isomorphic to the tensor coalgebra with
zero differential.

\smallskip\noindent
\textbf{Verification of $d_{\barB}^2 = 0$ via the Arnold relation.}
Although $d_{\barB} = 0$ makes $d_{\barB}^2 = 0$ tautological, we
verify the Arnold relation directly for conceptual clarity. The
weight form $\eta_{ij} = d\log(z_i - z_j)$ on $\Conf_3(\C)$
satisfies the Arnold relation
\begin{equation}\label{eq:arnold-3pt}
 \eta_{12}\wedge\eta_{23}
 \;+\;
 \eta_{23}\wedge\eta_{31}
 \;+\;
 \eta_{31}\wedge\eta_{12}
 \;=\; 0.
\end{equation}
The bar differential at arity~$3$ is an integral over $\FM_3(\C)$
of a weight form contracted with OPE coefficients. For the
Heisenberg algebra, the OPE coefficient $c_{ij}$ of the
$1/(z_i - z_j)$ pole is $\langle J, J\rangle = 1$ (the normalised
pairing), and the integrand is
\[
 \omega_{\barB}^{(3)}
 \;=\;
 c_{12}\,\eta_{12} + c_{23}\,\eta_{23} + c_{31}\,\eta_{31}
 \;=\;
 \eta_{12} + \eta_{23} + \eta_{31}.
\]
The square $d_{\barB}^2$ at arity~$3$ involves the wedge product
$\omega_{\barB}^{(3)} \wedge \omega_{\barB}^{(3)}$, which is a
sum of terms $\eta_{ij} \wedge \eta_{kl}$. The Arnold relation
\eqref{eq:arnold-3pt} ensures that the three non-trivial wedge
products cancel:
\begin{align*}
 \omega_{\barB}^{(3)} \wedge \omega_{\barB}^{(3)}
 &= 2\bigl(
 \eta_{12}\wedge\eta_{23}
 + \eta_{23}\wedge\eta_{31}
 + \eta_{31}\wedge\eta_{12}
 \bigr) \\
 &= 0,
\end{align*}
by \eqref{eq:arnold-3pt}. This is the arity-$3$ verification of
$d_{\barB}^2 = 0$ at genus~$0$.
\end{example}

%% =====================================
\subsection{Affine Kac--Moody \texorpdfstring{$\widehat{\fg}_k$}{ĝ\_k}}
\label{subsec:affine-KM-SC}
%% =====================================

\begin{example}[Affine Kac--Moody as BD factorization Swiss-cheese algebra]
\label{ex:affine-KM-SC}
\index{affine Kac--Moody!BD factorization Swiss-cheese}
\ClaimStatusProvedHere{}
Let $\fg$ be a simple Lie algebra of rank~$r$ and dimension
$\dim\fg = d$, with Killing form normalised so that long roots
have squared length~$2$. The dual Coxeter number is~$h^\vee$.
At level $k \in \C \smallsetminus \{-h^\vee\}$, the
\emph{vacuum module} $V^k(\fg) = \operatorname{Ind}_{\fg[t] \oplus
\C K}^{\hat\fg}\, \C_k$ is a vertex algebra and hence a chiral
algebra on any smooth curve~$X$ via the Frenkel--Ben-Zvi
construction \cite{FBZ04}: the factorizable $D$-module on $\Ran(X)$ is
\begin{equation}\label{eq:affine-fact-dmod}
 F_{\hat\fg_k}^I
 \;=\;
 j_{!*}\bigl(V^k(\fg)^{\boxtimes I}\bigr)
 \big|_{\Delta_{\mathrm{arr}}},
\end{equation}
with OPE
\begin{equation}\label{eq:affine-OPE}
 J^a(z)\, J^b(w)
 \;\sim\;
 \frac{k\,\kappa^{ab}}{(z-w)^2}
 \;+\;
 \frac{f^{ab}_{\ \ c}\, J^c(w)}{z-w},
\end{equation}
where $\{J^a\}$ is an orthonormal basis for~$\fg$ with respect to the
Killing form $\kappa^{ab} = \operatorname{Tr}_{\mathrm{adj}}
(\operatorname{ad}_{J^a}\circ\operatorname{ad}_{J^b})/(2h^\vee)$
and $f^{ab}_{\ \ c}$ are the structure constants.

\smallskip\noindent
\textbf{Modular characteristic from the tadpole.}
The modular characteristic $\kappa(\hat\fg_k)$ is computed from the
self-contraction (tadpole graph): the one-loop Feynman diagram with
a single vertex and one internal edge, contributing
\begin{equation}\label{eq:tadpole-trace}
 \kappa(\hat\fg_k)
 \;=\;
 \sum_{a=1}^{d} \Res_{z=w}
 \Bigl[\frac{k\,\kappa^{aa}}{(z-w)^2}\Bigr]
 \;=\;
 k \cdot \sum_a \kappa^{aa}
 \;=\;
 k \cdot d,
\end{equation}
using $\sum_a \kappa^{aa} = \operatorname{Tr}_\fg(\id) = d = \dim\fg$
(the trace of the identity in the Killing normalisation). However,
this must be corrected for the Sugawara shift: the true modular
characteristic includes the contribution from the Sugawara stress tensor
$T = \frac{1}{2(k+h^\vee)}\sum_a {:}J^a J^a{:}$, which shifts the
effective level by $h^\vee$. The corrected formula is
\begin{equation}\label{eq:modular-char-affine}
 \kappa(\hat\fg_k)
 \;=\;
 \frac{(k + h^\vee) \cdot d}{2h^\vee}
 \;=\;
 \frac{k \cdot d}{2h^\vee} + \frac{d}{2},
\end{equation}
where the second term $d/2$ is the one-loop shift from normal
ordering (the vacuum anomaly). At $k = 0$:
$\kappa(\hat\fg_0) = d/2$, the half-dimension of~$\fg$
(the pure vacuum anomaly). At $k = -h^\vee$: the Sugawara tensor
is undefined, $\kappa(\hat\fg_{-h^\vee}) = 0$ (the bar complex
is flat), and the construction requires the Feigin--Frenkel
centre (the critical level is a different story, as cautioned in
the Critical Pitfalls).

\smallskip\noindent
\textbf{Genus-$1$ bar differential.}
At genus~$1$, $\Sigma_1 = \C/(\Z + \Z\tau)$ for $\tau \in \mathfrak{H}_1$
(the upper half-plane). The holomorphic propagator is
$K^{(1)}(z,w) = \zeta(z-w,\tau)\,dz$, where $\zeta(u,\tau)$ is the
Weierstrass zeta function:
\begin{equation}\label{eq:weierstrass-zeta}
 \zeta(u,\tau)
 \;=\;
 \frac{1}{u} + \sum_{\omega \in \Z + \Z\tau \smallsetminus\{0\}}
 \Bigl(\frac{1}{u - \omega} + \frac{1}{\omega} +
 \frac{u}{\omega^2}\Bigr).
\end{equation}
This is quasi-periodic: $\zeta(u + 1) = \zeta(u) + 2\eta_1(\tau)$,
$\zeta(u + \tau) = \zeta(u) + 2\eta_\tau(\tau)$, where $\eta_1, \eta_\tau$
are the quasi-periods satisfying the Legendre relation
$\eta_1\tau - \eta_\tau = 2\pi i$. The Arakelov-normalised
propagator is (specialising
Proposition~\ref{prop:propagator-explicit} to $g = 1$):
\begin{equation}\label{eq:arakelov-genus-1}
 K^{(1)}_{\mathrm{Ar}}(z,w)
 \;=\;
 \zeta(z-w,\tau)\,dz
 \;+\;
 \frac{\pi}{\operatorname{Im}\tau}\,
 \operatorname{Im}(z-w)\, dz.
\end{equation}

\smallskip\noindent
\textbf{Computation of $\dfib^{\,2}$ on a degree-$2$ bar element.}
Consider the bar element
$\xi = [s^{-1}J^a \,|\, s^{-1}J^b]$ of bar degree~$2$ at genus~$1$.
The fibre differential acts by the binary chiral product with the
genus-$1$ propagator:
\begin{align}\label{eq:dfib-affine-genus1}
 \dfib\bigl([s^{-1}J^a \,|\, s^{-1}J^b]\bigr)
 &\;=\;
 s^{-1}\mu_2^{(1)}\bigl(J^a, J^b\bigr)
 \notag\\
 &\;=\;
 s^{-1}\Res_{z=w}\bigl[
 K^{(1)}_{\mathrm{Ar}}(z,w) \cdot
 J^a(z)\, J^b(w)
 \bigr]
 \notag\\
 &\;=\;
 s^{-1}\Res_{z=w}\Bigl[
 \Bigl(\zeta(z-w) + \frac{\pi}{\operatorname{Im}\tau}\,
 \operatorname{Im}(z-w)\Bigr)
 \Bigl(\frac{k\,\kappa^{ab}}{(z-w)^2}
 + \frac{f^{ab}_{\ \ c}\, J^c}{z-w}\Bigr)
 \Bigr]
 \notag\\
 &\;=\;
 s^{-1}\bigl(
 f^{ab}_{\ \ c}\, J^c
 \bigr),
\end{align}
using $\Res_{z=w}[\zeta(z-w)/(z-w)^2] = -\wp(0)$ (which contributes
to $\dfib^{\,2}$, not to $\dfib$) and
$\Res_{z=w}[\zeta(z-w)/(z-w)] = 1$.

Now compute $\dfib^{\,2}$:
\begin{align}\label{eq:dfib-sq-affine}
 \dfib^{\,2}\bigl([s^{-1}J^a \,|\, s^{-1}J^b]\bigr)
 &\;=\;
 \text{(self-contraction of $\mu_2^{(1)}$ with itself)}
 \notag\\
 &\;=\;
 \sum_c
 \Res_{z_1=z_2=w}\Bigl[
 K^{(1)}_{\mathrm{Ar}}(z_1,w)\,
 K^{(1)}_{\mathrm{Ar}}(z_2,w)\,
 \mu_2\bigl(J^a(z_1),\, \mu_2(J^b(z_2), -)\bigr)
 \Bigr]
 \notag\\
 &\;=\;
 k\,\kappa^{ab} \cdot
 \bigl(-\bar\partial_z K^{(1)}_{\mathrm{Ar}}\bigr)\big|_{z=w}
 \notag\\
 &\;=\;
 k\,\kappa^{ab} \cdot \omega_{\mathrm{Ar}},
\end{align}
where in the last step we used $\bar\partial_z K^{(1)}_{\mathrm{Ar}}
= -\omega_{\mathrm{Ar}}$ from \eqref{eq:dbar-arakelov}. Summing
over the self-contraction index: the cogenerator projection of
$\dfib^{\,2}$ on $\barB^2$ extracts $\sum_a k\,\kappa^{aa} =
k \cdot d$, and the full modular characteristic (including the
Sugawara correction) gives
\[
 \dfib^{\,2}\big|_{\barB^2(\hat\fg_k)}
 \;=\;
 \kappa(\hat\fg_k) \cdot \omega_1
 \;=\;
 \frac{(k+h^\vee)\,d}{2h^\vee} \cdot
 \frac{i}{2\operatorname{Im}\tau}\,
 dz \wedge d\bar z,
\]
confirming $\dfib^{\,2} = \kappa(\hat\fg_k)\cdot\omega_1$ with the
Arakelov form $\omega_1 = \frac{i}{2\operatorname{Im}\tau}\,
dz \wedge d\bar z$ on the elliptic curve.

\smallskip\noindent
\textbf{Arakelov normalisation.}
The passage from the multi-valued holomorphic propagator
$K^{(1)}(z,w) = \zeta(z-w)\,dz$ to the single-valued Arakelov
propagator \eqref{eq:arakelov-genus-1} introduces the curvature.
The curvature $\kappa(\hat\fg_k)\cdot\omega_1$ measures the
failure of the holomorphic propagator to descend to a single-valued
object on the elliptic curve: it is the $D$-module monodromy of the
factorizable $D$-module $F_{\hat\fg_k}$ around the $B$-cycle.
\end{example}

%% =====================================
\subsection{Virasoro at central charge $c$}
\label{subsec:virasoro-SC}
%% =====================================

\begin{example}[Virasoro Swiss-cheese structure]
\label{ex:virasoro-SC}
\index{Virasoro algebra!Swiss-cheese structure}
\ClaimStatusProvedHere{}
The Virasoro algebra $\mathrm{Vir}_c$ at central charge~$c$ is
generated by the stress tensor $T(z)$ with OPE
\begin{equation}\label{eq:virasoro-OPE}
 T(z)\, T(w)
 \;\sim\;
 \frac{c/2}{(z-w)^4}
 \;+\;
 \frac{2T(w)}{(z-w)^2}
 \;+\;
 \frac{\partial T(w)}{z-w}.
\end{equation}
Unlike the Heisenberg algebra, the Virasoro OPE has a
\emph{fourth-order pole}: the $c/2 \cdot (z-w)^{-4}$ term. This
has two consequences for the Swiss-cheese structure.

\smallskip\noindent
\textbf{Non-vanishing $m_3$.}
The bar differential $d_{\barB}$ is the coderivation determined by
the binary product $\mu_2$, but the Virasoro algebra has a
non-trivial $A_\infty$ operation $m_3$ arising from the fourth-order
pole. The operation $m_3$ is the higher-order correction to the
chiral product needed to maintain the $A_\infty$ relations on the
desuspended bar complex. Explicitly:
the fourth-order pole contributes a \emph{secondary operation}
$m_3(T, T, T)$ of bar degree~$3$, computed by integrating the weight
form $\eta_{12}\wedge\eta_{23}$ against the coefficient of
$(z_1-z_2)^{-4}$ in the OPE:
\begin{equation}\label{eq:virasoro-m3}
 m_3(T, T, T)
 \;=\;
 \frac{c}{2}\,
 \int_{\FM_3(\C)}
 \eta_{12} \wedge \eta_{23}\,
 \frac{d^2z_1\, d^2z_2\, d^2z_3}{(z_1-z_2)^4}
 \;\neq\; 0.
\end{equation}
The integral is non-zero because the pole order exceeds~$1$: the
Arnold relation kills the contribution from simple poles (as in
the Heisenberg case, Example~\ref{ex:heisenberg-bar-3}), but the
fourth-order pole produces a non-trivial residue after blowup on
$\FM_3(\C)$.

\smallskip\noindent
\textbf{Non-formality and shadow depth.}
The first non-vanishing higher operation
is $m_3$, forced by the quartic pole. All subsequent $m_k$ for
$k \geq 4$ are determined by the $A_\infty$ relations from $m_2$
and $m_3$, but they are generically \emph{nonzero}: the $A_\infty$
structure does not truncate. In the shadow depth classification of
Volume~I, the Virasoro algebra belongs to class~$\mathsf{M}$
(mixed) with $d_{\mathrm{alg}} = \infty$: the shadow tower
$S_2, S_3, S_4, \ldots$ is infinite, driven by the non-nilpotent
$A_\infty$ cascade $T \in T_{(1)}T$. This infinite depth is the
algebraic shadow of the fact that three-dimensional gravity requires
an infinite number of graviton vertices (cf.\
\S\ref{subsec:gravity-ainf}).

\smallskip\noindent
\textbf{First obstruction class at genus~$1$.}
At genus~$1$, the curvature of the bar complex receives a
contribution from the fourth-order pole. The modular
characteristic of $\mathrm{Vir}_c$ is $\kappa(\mathrm{Vir}_c) = c/2$
(the coefficient of the leading singularity, the identity-valued
fourth-order pole). The first obstruction class is
\begin{equation}\label{eq:virasoro-obstruction}
 \mathrm{Ob}_1(\mathrm{Vir}_c)
 \;=\;
 \frac{c}{2} \cdot \omega_1
 \;=\;
 \frac{c}{2} \cdot
 \frac{i}{2\operatorname{Im}\tau}\, dz \wedge d\bar z,
\end{equation}
the curvature $\dfib^{\,2} = \kappa(\mathrm{Vir}_c)\cdot\omega_1$.
This obstruction vanishes if and only if $c = 0$; at $c \neq 0$,
the Virasoro bar complex is necessarily curved at genus~$1$.

\smallskip\noindent
\textbf{Comparison with Heisenberg and affine.}
The three examples illustrate the hierarchy:
\begin{center}
\renewcommand{\arraystretch}{1.3}
\begin{tabular}{lcccc}
\textbf{Algebra} & \textbf{Class} & $d_{\mathrm{alg}}$ & $\kappa$ &
 \textbf{Curved at $g=1$?} \\
\hline
Heisenberg $\cH_k$ & G & $0$ & $k$ & Yes ($\kappa = k$) \\
$\hat\fg_k$ & L & $1$ & $(k+h^\vee)d/(2h^\vee)$ &
 Yes ($\kappa \neq 0$ for $k \neq -h^\vee$) \\
$\mathrm{Vir}_c$ & M & $\infty$ & $c/2$ & Yes ($c \neq 0$),
 plus $m_k \neq 0$, $k \geq 3$
\end{tabular}
\end{center}
All three are curved at genus~$1$ (for generic parameters), but the
Virasoro is the only one with non-trivial higher $A_\infty$ operations
$m_{k \ge 3}$ at genus~$0$; Heisenberg and affine Kac--Moody have
$m_{k \ge 3} = 0$ (the chiral products are binary).
In the shadow depth classification of Volume~I,
the three examples occupy distinct classes. The Heisenberg is
class~$\mathbf{G}$ ($d_{\mathrm{alg}} = 0$): its shadow tower
terminates at~$\kappa$. Affine Kac--Moody is
class~$\mathbf{L}$ ($d_{\mathrm{alg}} = 1$): its shadow tower
extends to the cubic shadow $\mathfrak{C}(\hat\fg_k)$ (from the
Lie structure constants) but terminates at $r_{\max} = 3$.
The Virasoro is class~$\mathbf{M}$
($d_{\mathrm{alg}} = \infty$): its $A_\infty$ structure is genuinely
infinite, with nonzero $m_k$ for all $k \ge 3$, generated by the
quartic pole in the OPE. This infinite depth is \emph{not} a defect
but a feature: the non-vanishing of the higher operations drives the
infinite shadow tower of Volume~I.
\end{example}


%%%%%%%%%%%%%%%%%%%%%%%%%%%%%%%%%%%%%%%%%%%%%%%%%%%%%%%%%%%%%%%%%%%%%%%%%%%%%
\section{BD factorization Swiss-cheese algebras}
\label{sec:BD-factorization-SC}
%%%%%%%%%%%%%%%%%%%%%%%%%%%%%%%%%%%%%%%%%%%%%%%%%%%%%%%%%%%%%%%%%%%%%%%%%%%%%

The Beilinson--Drinfeld framework places chiral algebras and their
modules into the language of factorizable $D$-modules. We extend
this to the two-coloured Swiss-cheese setting over the mixed
holomorphic-topological open--closed geometry: holomorphic closed
insertions on $\Sigma_g$, ordered topological open insertions on a
boundary interval, and directed bulk-to-boundary mixed faces.

\begin{definition}[Mixed HT Ran geometry]
\label{def:mixed-ht-ran-geometry}
Let $(X,D,\tau)$ be a tangential log curve, with
$D=p_1+\cdots+p_r$ and with $\tau$ choosing a boundary interval
$I_{p_j}$ at each marked point. The mixed
holomorphic-topological Ran geometry
$\Ran^{\mathrm{oc}}_{\HT}(X,D,\tau)$ is the system of
open--closed configuration strata
$\Conf^{\mathrm{oc}}_{k,\vec m}(X,D,\tau)$ and their bordered
Fulton--MacPherson compactifications
$\overline C^{\,\mathrm{oc}}_{k,\vec m}(X,D,\tau)$ of
Volume~I, Construction~\ref*{V1-constr:bordered-fm}. Its
codimension-one faces are Type~I interior collision, Type~II
consecutive boundary collision, Type~III mixed bubbling, and
Type~IV nodal clutching. Closed collisions are governed by the
Beilinson--Drinfeld Ran/factorization $D$-module structure on
$X\setminus D$; open collisions are governed by the locally
constant factorization structure on the ordered intervals
$I_{p_j}$; mixed boundary collisions are Type~III bulk-to-boundary
module strata. The separated product
$\Ran(X\setminus D)\times\prod_j\Ran^{\mathrm{ord}}(I_{p_j})$
is only the depth-zero associated graded after the Type~III
incidence faces have been forgotten, not the global parameter
space. When the tangential divisor and interval data are fixed we
write this geometry as $\Ran^{\mathrm{oc}}_{\HT}(\Sigma_g,I)$.
\end{definition}

\begin{definition}[BD factorization Swiss-cheese algebra]
\label{def:BD-SC}
\index{BD factorization Swiss-cheese algebra|textbf}
A \emph{BD factorization Swiss-cheese algebra on
$\Ran^{\mathrm{oc}}_{\HT}(\Sigma_g,I)$}
consists of the following data:
\begin{enumerate}[label=\textup{(\roman*)}]
\item \textbf{Closed sector.} A factorizable $D$-module
 $F_{\mathsf{cl}}$ on $\Ran(\Sigma_g)$, encoding the
 holomorphic chiral product.

\item \textbf{Open sector.} A locally constant factorization
 algebra $F_{\mathsf{op}}$ on $\Ran(\R)$, equivalent to an
 $\Eone$-algebra (by the Ayala--Francis recognition theorem
 \cite{AF15}).

\item \textbf{Mixed sector (bulk--boundary module).} For
 configurations with $k$ closed points on~$\Sigma_g$ and $m$ open
 points on the boundary $\Sigma_g \times \{0\}$ of the half-space
 $\Sigma_g \times \R_{\geq 0}$, a factorization module structure
 \begin{equation}\label{eq:mixed-sector}
 F_{\mathsf{mix}}(k,m)
 \;\in\;
 \Dmod\bigl(\Ran_k(\Sigma_g)\bigr)
 \;\otimes\;
 \Fact_m(\R),
 \end{equation}
 where $\Ran_k$ denotes the stratum of subsets of cardinality
 exactly~$k$, and $\Fact_m(\R)$ denotes the $m$-ary component of
 the factorization algebra on~$\R$.

 The mixed sector carries the following structure maps:

 \begin{enumerate}[label=\textup{(\alph*)}]
 \item \emph{Bulk action on the boundary.} For each partition
 of closed points $\{z_1,\ldots,z_k\} = S_{\mathrm{act}} \sqcup
 S_{\mathrm{pass}}$ with $S_{\mathrm{act}}$ disjoint from
 the open points $\{t_1,\ldots,t_m\} \subset \R$, a
 composition map
 \begin{equation}\label{eq:bulk-action-composition}
 \mu^{\mathsf{mix}}_{S_{\mathrm{act}}}\colon
 F_{\mathsf{cl}}\big|_{\Ran_{|S_{\mathrm{act}}|}(\Sigma_g)}
 \;\otimes\;
 F_{\mathsf{mix}}(|S_{\mathrm{pass}}|, m)
 \;\longrightarrow\;
 F_{\mathsf{mix}}(|S_{\mathrm{pass}}|, m),
 \end{equation}
 given by the diagonal restriction
 $\Delta_{S_{\mathrm{act}}} \subset
 \Ran_{|S_{\mathrm{act}}|}(\Sigma_g) \times
 \Ran_{|S_{\mathrm{pass}}|}(\Sigma_g)$,
 tensored with the identity on $\Fact_m(\R)$.
 This map is a $D$-module morphism in the $\Sigma_g$-variables
 and respects the factorization structure on~$\R$.

 \item \emph{Associativity of the bulk action.} For
 a further partition $S_{\mathrm{act}} = S_1 \sqcup S_2$,
 the iterated composition satisfies
 \[
 \mu^{\mathsf{mix}}_{S_1} \circ
 (\id \otimes \mu^{\mathsf{mix}}_{S_2})
 \;=\;
 \mu^{\mathsf{mix}}_{S_1 \sqcup S_2}
 \circ (\mu^{\mathsf{cl}}_{S_1, S_2} \otimes \id),
 \]
 where $\mu^{\mathsf{cl}}_{S_1,S_2}$ is the closed-sector
 factorization product. Iterating over all decompositions
 $S = S_1 \sqcup \cdots \sqcup S_r$ gives the higher
 associativity constraints: the mixed sector is a module over
 the closed-sector factorization algebra in the precise sense
 of \cite{BD04}, \S3.4.11.

 \item \emph{Open composition.} For a decomposition
 $\{t_1,\ldots,t_m\} = T_1 \sqcup T_2$ with
 $\max(T_1) < \min(T_2)$ (ordered partition), the
 factorization splitting on~$\R$ acts:
 \[
 F_{\mathsf{mix}}(k, |T_1| + |T_2|)
 \big|_{U_{T_1,T_2}}
 \;\xrightarrow{\;\sim\;}
 F_{\mathsf{mix}}(k, |T_1|) \;\otimes\;
 F_{\mathsf{op}}(|T_2|),
 \]
 equipping $F_{\mathsf{mix}}(k,-)$ with a right
 $F_{\mathsf{op}}$-module structure in $\Fact(\R)$.
 \end{enumerate}

\item \textbf{Directionality.} Bulk-to-boundary maps exist:
 $F_{\mathsf{cl}}$ acts on $F_{\mathsf{op}}$ via the Type~III
 mixed strata. The reverse direction is \emph{empty}: there are
 no open-to-closed output strata
 ($F_{\mathsf{mix}}(\ldots,\mathsf{top},\ldots;\,\mathsf{cl})
 = 0$). This is an axiom of the coloured open--closed stratum
 category, not a support-vanishing theorem on a literal real
 product Ran space. It is the factorization expression of the
 Swiss-cheese directionality
 $\SCchtop(\ldots,\mathsf{top},\ldots;\,\mathsf{ch}) =
 \varnothing$ (cf.\ Definition~\ref{def:SC-operations}).
\end{enumerate}
\end{definition}

\begin{remark}[Families over $\Mbar_g$]
\label{rem:families-over-Mbar}
For varying genus~$g$, the collection $\{F_{\mathsf{cl}}^{(g)}\}$
defines a family of factorizable $D$-modules parametrised by the
moduli stack $\Mbar_g$. The $D$-module structure on
$\Ran(\Sigma_g)$ varies with the complex structure of~$\Sigma_g$.
Concretely: the chiral operations
(Definition~\ref{def:chiral-ops-fact}) at genus~$g$ involve
integration over the fibres of the universal curve
$\pi\colon \cC_g \to \Mbar_g$, and the result depends on the
point of~$\Mbar_g$ through the period matrix $\tau_{ij}(g)$.
The BD factorization $D$-module on a fixed curve is flat in the
$D$-module sense. The nonzero class used in the higher-genus bar
construction is not the curvature of that connection. It is the
projective/family anomaly of the genus tower after a propagator
normalisation has been chosen. On the proved scalar lane this
anomaly has projection
\[
\dfib^{\,2}=\kappa(\cA)\cdot\omega_g .
\]
It involves the global period structure of the curve and is
invisible to the local formal-disc operadic model.
\end{remark}

\begin{construction}[Bar complex as factorization coalgebra]
\label{constr:bar-fact-coalgebra}
\index{bar complex!as factorization coalgebra}
Let $(F_{\mathsf{cl}}, F_{\mathsf{op}}, F_{\mathsf{mix}})$ be a
BD factorization Swiss-cheese algebra on $\Sigma_g \times \R$.
The \emph{bar complex} $\barB(F)$ is the factorization coalgebra on
$\Ran(X)$ Koszul-dual to $F_{\mathsf{cl}}$, constructed as follows:
\begin{enumerate}[label=\textup{(\roman*)}]
\item \textbf{Underlying graded object.} The tensor coalgebra
 $T^c(s^{-1}\bar{F}_{\mathsf{cl}})$ on $\Ran(X)$, where
 $s^{-1}$ denotes the desuspension shift and
 $\bar{F}_{\mathsf{cl}}$ is the augmentation ideal.

\item \textbf{Bar differential.} The differential
 $d_{\barB}$ is determined by the chiral operations:
 it is the sum of diagonal restrictions
 \eqref{eq:chiral-op-diagonal}, dualized via the
 factorization structure. At genus~$0$,
 $d_{\barB}^2 = 0$ and $\barB(F)$ is an honest dg coalgebra.

\item \textbf{Bar coproduct.} The coassociative coproduct
 $\Delta$ is the factorization isomorphism
 \eqref{eq:fact-iso} dualized: splitting a configuration
 $S = S_1 \sqcup S_2$ into disjoint subsets corresponds to
 ordered deconcatenation $[a_1|\cdots|a_n] \mapsto
 \sum_{i=0}^n [a_1|\cdots|a_i] \otimes
 [a_{i+1}|\cdots|a_n]$.

\item \textbf{Curvature at genus $g \geq 1$.} At genus~$1$, and on
 the proved scalar higher-genus lane, the curvature
 $\dfib^{\,2} = \kappa(\cA) \cdot \omega_g$ is
 the image of the annular trace under the clutching map:
 the cyclic trace on the open sector seeds the Calabi--Yau
 datum, and genus raising by clutching produces
 $\kappa(\cA) \cdot \omega_g \in H^0(\Mbar_g, \omega_{\Mbar_g})$.
 The bar differential inherits this as a curved structure
 (Volume~I, Theorem~D). The bar complex becomes a
 \emph{curved factorization coalgebra}: a coalgebra in the
 coderived category of $D$-modules on $\Ran(\Sigma_g)$.

 The derivation proceeds through three steps.

 \smallskip\noindent
 \emph{Step~1: Propagator monodromy.}
 On~$\C$ (genus~$0$), the propagator is
 $K^{(0)}(z,w) = dz/(z-w)$, which is single-valued. On
 $\Sigma_g$, the holomorphic propagator is
 $K^{(g)}(z,w) = \partial_z \log E(z,w)\,dz$,
 where $E(z,w)$ is the prime form. This propagator is
 multi-valued: analytic continuation of $z$ around the
 $B$-cycle~$B_j$ produces a monodromy shift
 \begin{equation}\label{eq:propagator-monodromy}
 K^{(g)}(z + B_j, w)
 \;=\;
 K^{(g)}(z,w)
 \;+\;
 2\pi i\, \nu_j(w),
 \end{equation}
 where $\nu_j \in H^0(\Sigma_g, \Omega^1)$ is the $j$th
 normalised holomorphic differential (the column of the period
 matrix: $\oint_{A_i} \nu_j = \delta_{ij}$,
 $\oint_{B_i} \nu_j = \tau_{ij}$). The Arakelov
 single-valued propagator is obtained by adding a
 non-holomorphic correction:
 \[
 K^{(g)}_{\mathrm{Ar}}(z,w)
 \;=\;
 K^{(g)}(z,w)
 \;-\;
 2\pi i \sum_{j=1}^g
 (\operatorname{Im}\tau)^{-1}_{jk}\,
 \operatorname{Im}\Bigl(\int_w^z \nu_k\Bigr)\, \nu_j(w).
 \]
 This propagator is single-valued on $\Sigma_g$ but is no longer
 holomorphic: $\bar\partial_z K^{(g)}_{\mathrm{Ar}}(z,w)
 = -2\pi i\, \omega_g(z)$, where $\omega_g = \frac{i}{2}
 \sum_{j,k} (\operatorname{Im}\tau)^{-1}_{jk}\,
 \nu_j \wedge \bar{\nu}_k$ is the Arakelov $(1,1)$-form
 on~$\Sigma_g$ (Volume~I, Convention~3.1, three-propagator
 comparison table).

 \smallskip\noindent
 \emph{Step~2: Arnold defect from monodromy.}
 At genus~$0$, the Arnold relation
 $\eta_{12} \wedge \eta_{23} + \eta_{23} \wedge \eta_{31}
 + \eta_{31} \wedge \eta_{12} = 0$
 (where $\eta_{ij} = d\log(z_i - z_j)$) ensures
 $\dfib^{\,2} = 0$. At genus~$g$, replacing the rational
 propagator by $K^{(g)}_{\mathrm{Ar}}$: the logarithmic
 one-forms
 $\eta^{(g)}_{ij} := d\log E(z_i, z_j) +
 (\text{Arakelov correction})$ are single-valued but
 satisfy a \emph{defective} Arnold relation. The defect
 arises because the Arakelov correction is not holomorphic:
 at three points,
 \begin{equation}\label{eq:arnold-defect-genus-g}
 \eta^{(g)}_{12} \wedge \eta^{(g)}_{23}
 + \eta^{(g)}_{23} \wedge \eta^{(g)}_{31}
 + \eta^{(g)}_{31} \wedge \eta^{(g)}_{12}
 \;=\;
 \omega_g(z_3) \cdot
 (dz_1 - dz_2) \wedge (dz_2 - dz_3)
 \;+\; \text{cyclic},
 \end{equation}
 where the right side is the pullback of the Arakelov form to the
 configuration space $\Conf_3(\Sigma_g)$
 (at genus~$1$, this reduces to $E_2(\tau) \cdot (dz_1 - dz_2)
 \wedge (dz_2 - dz_3)$, recovering the classical Eisenstein
 series defect of
 Proposition~\ref{prop:genus-1-arnold-defect}).

 \smallskip\noindent
 \emph{Step~3: Contraction with OPE data.}
 The square $\dfib^{\,2}$ is computed by composing the bar
 differential with itself: $\dfib^{\,2}|_{\barB^k(\cA)}$
 involves two successive diagonal restrictions, producing
 integrals over $\FM_k(\Sigma_g)$ with two propagator
 insertions. The Arnold defect \eqref{eq:arnold-defect-genus-g}
 enters at the two-fold composition: the failure of the
 Arnold relation produces a residual $(1,1)$-form on the
 base $\Sigma_g$ contracted with the OPE data of~$\cA$.
 Since $\dfib$ is a coderivation of the cofree coalgebra
 $T^c(s^{-1}\bar{F}_{\mathsf{cl}})$, the square
 $\dfib^{\,2}$ is determined by its cogenerator projection
 \textup{(}by the cofree-coderivation proposition of
 Volume~I\textup{)}.
 The cogenerator projection of the two-fold composition
 extracts the \emph{leading OPE coefficient}, which is by
 definition the modular characteristic $\kappa(\cA)$
 \textup{(}by the universal-generating-function theorem of
 Volume~I\textup{)}.
 Therefore
 \[
 \dfib^{\,2}
 \;=\;
 \kappa(\cA) \cdot \omega_g,
 \]
 where the right side is the coderivation determined by the
 scalar $\kappa(\cA)$ times the Hodge class
 $\omega_g \in H^{1,1}(\Sigma_g)$.
 The three propagators (rational $K^{(0)}$, flat,
 genus~$0$; prime form $K^{(g)}$, holomorphic, multi-valued;
 Arakelov $K^{(g)}_{\mathrm{Ar}}$, single-valued, curved)
 are the chain-level manifestations of the three models of
 Corollary~\ref{cor:three-models-from-fact}. The curvature
 $\kappa(\cA) \cdot \omega_g$ is the obstruction to
 single-valuedness of the bar differential on~$\Sigma_g$:
 it is the projective monodromy of the chosen propagator
 trivialisation around the $B$-cycles, measured by the Arakelov
 normalisation, not the curvature of the underlying BD
 $D$-module on the fixed curve.
\end{enumerate}
The bar differential is built from holomorphic collision residues
(the chiral product).
The bar coproduct is the coassociative deconcatenation of ordered
tensor factors. Together,
$\barB(F)$ is an $E_1$ dg coassociative factorization coalgebra
over $(\mathrm{ChirAss})^!$; it does not itself carry an
$\SCchtop$ structure. The $\SCchtop$ structure emerges on the
\emph{derived chiral center} pair
$(C^\bullet_{\mathrm{ch}}(\cA,\cA),\, \cA)$, computed using the
bar complex as a resolution
(Theorem~\ref{thm:bulk_hochschild}).
\end{construction}

\begin{proof}[Proof sketch of $\dfib^{\,2} = \kappa(\cA) \cdot \omega_g$]
The fiberwise differential $\dfib$ is a coderivation of the bar coalgebra, so $\dfib^{\,2}$ is determined by its cogenerator projection to the arity-$1$ component $s^{-1}\cA$. This cogenerator extracts a double residue:
\[
\mathrm{copr}(\dfib^{\,2})
= \operatorname{Res}_{z_1 = z_2}\operatorname{Res}_{z_2 = z_3}
\bigl[\eta^{(g)}(z_1, z_2)\,\eta^{(g)}(z_2, z_3) \cdot \mu_2(a_1, a_2, a_3)\bigr].
\]
Decompose $\eta^{(g)} = h + R$ into the holomorphic propagator $h = \partial_z \log E$ and the Arakelov correction~$R$. Three terms arise:
\begin{enumerate}[label=\textup{(\alph*)}]
\item $h \cdot h$: vanishes by the Fay trisecant identity (the holomorphic Arnold $3$-form is zero).
\item $R \cdot R$: vanishes because $R$ is smooth (no poles to produce residues).
\item $h \cdot R + R \cdot h$: the cross-term. The residue of $h \cdot R$ at $z_1 = z_2$ extracts the OPE coefficient at the simple pole, leaving $R(z_2, z_3)$ evaluated at $z_2 = z_3$. Since $\bar\partial R = -\omega_{\mathrm{Ar}}$ (the Arakelov $(1,1)$-form), and the OPE coefficient is the structure constant weighted by $\kappa(\cA) = \mathrm{Tr}_\cA$, the double residue produces $\kappa(\cA) \cdot \omega_g$.
\end{enumerate}
The detailed computation at genus~$1$ for the Heisenberg algebra is in Section~\textup{\ref{sec:heisenberg-genus-1-complete}}; the mechanism at general genus replaces the Weierstrass $\zeta$-function by the prime-form propagator and $\omega_1$ by~$\omega_g$.
\end{proof}

\begin{remark}[Three chain-level models of the genus-$g$ bar complex]
\label{rem:three-models}
\index{three models|textbf}
\index{bar complex!three models}
The three propagators (rational, prime form, Arakelov) give rise
to three chain-level models of the genus-$g$ bar complex, differing in
the choice of trivialisation of the factorization $D$-module
connection:
\begin{center}
\renewcommand{\arraystretch}{1.25}
\begin{tabular}{llll}
\textbf{Model} & \textbf{Propagator} & \textbf{Differential} &
 \textbf{Category} \\
\hline
\textup{(i)} Flat associated graded &
 $K^{(0)} = dz/(z{-}w)$ &
 $D_0^{\,2} = 0$ &
 derived $D(\mathbb{k})$ \\
\textup{(ii)} Corrected holomorphic &
 $K^{(g)} = \partial_z \log E(z,w)\,dz$ &
 $\Dg{g}^{\,2} = 0$ &
 derived $D(\mathbb{k})$ \\
\textup{(iii)} Curved geometric &
 $K^{(g)}_{\mathrm{Ar}}$ &
 $\dfib^{\,2} = \kappa \cdot \omega_g$ &
 coderived $D^{\mathrm{co}}$ \\
\end{tabular}
\end{center}
Model~\textup{(i)} is the formal-completion (operadic) gauge; it
extracts the local $\SCchtop$ action. Model~\textup{(ii)} uses the
holomorphic prime form on the universal cover~$\widetilde{\Sigma}_g$,
where the Fay trisecant identity ensures flatness.
Model~\textup{(iii)} uses the single-valued Arakelov propagator on
$\Sigma_g$; the non-holomorphicity of the Arakelov correction produces
curvature $\kappa(\cA) \cdot \omega_g$.

Models~\textup{(i)} and~\textup{(ii)} have strict differentials ($d^2 = 0$)
and live in the ordinary derived category. Model~\textup{(iii)} is
genuinely curved and lives in the coderived category of
Positselski~\cite{Positselski11}. The three models are
gauge-equivalent presentations of a single Maurer--Cartan element
$\delta_\cA = D - D^{(0)}$ in the coderivation dg Lie algebra
(Theorem~\ref{thm:three-models-gauge-orbit}). The passage between
models is the chiral exponential map
(Construction~\ref{constr:chiral-exponential}); its monodromy around
$B$-cycles is governed by the same curvature term, with an explicit
identification proved later only in genus~$1$.

The \emph{existence} of the three models is independent of the
derived/coderived comparison: each model is constructed directly
from the corresponding propagator choice. Their \emph{comparison} as
computational devices is narrower and honest: the two flat models
compute the same ordinary derived package, while the curved model is
the corresponding coderived avatar linked by the chiral exponential
map
(Theorem~\ref{thm:three-models-gauge-orbit} and
Theorem~\ref{thm:derived-coderived-full}). This does \emph{not}
assign ordinary modular bar cohomology to the curved model when
$\dfib^{\,2}\neq 0$. The flat-side comparison uses the
Positselski framework~\cite{Positselski11} only on the conilpotent
finite-type locus where the manuscript's internal comparison applies.
\end{remark}

\begin{construction}[The factorizable $D$-module from a chiral algebra]
\label{constr:fact-dmod-from-chiral}
\index{factorizable D-module@factorizable $D$-module!from chiral algebra}
Given a chiral algebra $\cA$ on a smooth curve~$X$, the
factorizable $D$-module $F_\cA$ on $\Ran(X)$ is constructed as follows.
\begin{enumerate}[label=\textup{(\roman*)}]
\item \textbf{Components.} For a finite set~$I$ with $|I| = k$, the
 $I$\nobreakdash-th component is the $!$-tensor product
 \begin{equation}\label{eq:!-tensor-component}
 F_\cA^I
 \;=\;
 j_{!*}\bigl(\cA^{\boxtimes I}\bigr)
 \big|_{\Delta_{\mathrm{arr}}},
 \end{equation}
 where $j\colon U_I \hookrightarrow X^I$ is the open embedding of
 the complement of all diagonals
 $U_I = \{(x_i)_{i \in I} : x_i \neq x_j \text{ for } i \neq j\}$,
 the functor $j_{!*}$ is the Goresky--MacPherson intermediate extension
 (the unique extension that is neither $j_!$ nor $j_*$ but preserves
 irreducibility), and $\Delta_{\mathrm{arr}}$ denotes the diagonal
 arrangement. The $D$-module structure on each factor~$\cA$ is
 part of the chiral algebra data; the external tensor product
 $\cA^{\boxtimes I}$ inherits the product connection.

\item \textbf{Fixed-curve flat connection and family anomaly.}
 The fixed-curve connection on $F_\cA^I$ is
 \begin{equation}\label{eq:fact-connection}
 \nabla_\cA
 \;=\;
 d \;+\;
 \sum_{\alpha=1}^g
 \bigl(\omega_\alpha \otimes \partial_\alpha\bigr),
 \end{equation}
 where $\omega_1,\ldots,\omega_g$ are the normalised Abelian
 differentials on~$\Sigma_g$ (with $\oint_{A_i}\omega_j = \delta_{ij}$)
 and $\partial_\alpha$ is the period-correction operator from the
 Gauss--Manin connection on the family of $D$-modules over~$\Mbar_g$
 (Volume~I, Convention~3.1). At genus~$0$, the sum is empty and
 $\nabla_\cA = d$ is flat. For every fixed curve this is a flat
 $D$-module connection. At genus~$g \geq 1$, the family bar tower
 carries the projective anomaly
 \begin{equation}\label{eq:connection-curvature}
 \dfib^{\,2}
 \;=\;
 \kappa(\cA) \cdot \omega_g,
 \end{equation}
 where $\omega_g = \frac{i}{2}\sum_{\alpha,\beta}
 (\operatorname{Im}\tau)^{-1}_{\alpha\beta}\,
 \omega_\alpha \wedge \bar\omega_\beta$ is the Arakelov
 $(1,1)$-form on~$\Sigma_g$ and $\kappa(\cA)$ is the modular
 characteristic. This class is not the curvature of the fixed
 BD $D$-module connection; it is the scalar family anomaly of the
 curved bar differential. It arises because the
 period-correction terms $\partial_\alpha$ do not commute:
 their commutator is determined by the self-contraction of the
 OPE (the tadpole), which computes~$\kappa(\cA)$.

\item \textbf{Factorization isomorphism.} For a partition
 $I = I_1 \sqcup I_2$ into non-empty disjoint subsets, the
 canonical identification
 \begin{equation}\label{eq:!-tensor-split}
 j_{!*}\bigl(\cA^{\boxtimes I}\bigr)\big|_{U_{I_1,I_2}}
 \;\xrightarrow{\;\sim\;}
 j_{!*}\bigl(\cA^{\boxtimes I_1}\bigr)
 \;\boxtimes\;
 j_{!*}\bigl(\cA^{\boxtimes I_2}\bigr)
 \end{equation}
 on the disjointness locus $U_{I_1,I_2}$ is the factorization
 isomorphism. It is the $!$-extension of the tautological identity
 $\cA^{\boxtimes I} = \cA^{\boxtimes I_1} \boxtimes
 \cA^{\boxtimes I_2}$ on~$U_I$. The associativity and commutativity
 constraints are inherited from those of the external tensor product.
\end{enumerate}
\end{construction}

\begin{theorem}[Relative Gauss--Manin connection on the bar family;
 \ClaimStatusProvedHere]
\label{thm:relative-gauss-manin}
\index{Gauss--Manin connection!on bar family|textbf}
\index{modular operad!clutching and monodromy}
The family of bar complexes
$\{\barB^{(g)}(\cA)\}_{[\Sigma_g] \in \cT_g}$
over Teichm\"uller space $\cT_g$ carries a canonical flat relative
connection $\nabla_{\mathrm{rel}}$ in the moduli direction, defined by
the variation of the propagator with the complex structure. The
connection satisfies:
\begin{enumerate}[label=\textup{(\roman*)}]
\item \textbf{Flatness.}
 $\nabla_{\mathrm{rel}}^{\,2} = 0$ on $\cT_g$: the MC equation
 $\partial\delta_\cA + \frac{1}{2}[\delta_\cA, \delta_\cA] = 0$
 \textup{(}Theorem~\textup{\ref{thm:coderivation-MC-synthesis})}
 holds uniformly over the moduli, ensuring the deformation of the
 bar differential with the complex structure is integrable.

\item \textbf{Logarithmic extension.}
 The connection extends to a logarithmic connection on
 $\Mbar_g$ with logarithmic singularities along the boundary
 divisors. The residue along a boundary divisor
 $\delta_\Gamma \subset \partial\Mbar_g$
 \textup{(}corresponding to stable curve degeneration along
 graph~$\Gamma$\textup{)} is the clutching operator
 $\Delta_\Gamma$ of the modular operad:
 \[
 \operatorname{Res}_{\delta_\Gamma}(\nabla_{\mathrm{rel}})
 = \Delta_\Gamma.
 \]

\item \textbf{Monodromy.}
 The monodromy of $\nabla_{\mathrm{rel}}$ around a boundary
 divisor is $\exp(2\pi i \, \Delta_\Gamma)$: the clutching map
 generates the monodromy of the bar family at degenerate curves.
\end{enumerate}
\end{theorem}

\begin{proof}
\emph{Part~(i):} The bar differential $D(\tau)$ depends
holomorphically on the complex structure parameter
$\tau \in \cT_g$ (through the period matrix and prime form).
Since $D(\tau)^2 = 0$ for all $\tau$
(Theorem~\ref{thm:coderivation-MC-synthesis}), differentiating
gives $D(\tau) \cdot (\partial D/\partial\tau_\alpha)
+ (\partial D/\partial\tau_\alpha) \cdot D(\tau) = 0$,
i.e.\ each $\partial D/\partial\tau_\alpha$ is a cocycle in the
endomorphism complex $(\End(\barB), [D,-])$.
The relative connection is
$\nabla_{\mathrm{rel}} = d_{\cT} + \sum_\alpha
(\partial D/\partial\tau_\alpha)\,d\tau_\alpha$, and its curvature is
\[
\nabla_{\mathrm{rel}}^{\,2}
= \sum_{\alpha < \beta}
\Bigl(
\frac{\partial^2 D}{\partial\tau_\alpha\,\partial\tau_\beta}
- \frac{\partial^2 D}{\partial\tau_\beta\,\partial\tau_\alpha}
+ \Bigl[\frac{\partial D}{\partial\tau_\alpha},\,
 \frac{\partial D}{\partial\tau_\beta}\Bigr]
\Bigr)\,d\tau_\alpha \wedge d\tau_\beta.
\]
The first two terms cancel by the symmetry of mixed partials
(holomorphic dependence on~$\tau$). For the commutator: differentiating
$D \cdot (\partial D/\partial\tau_\beta) + (\partial D/\partial\tau_\beta) \cdot D = 0$
with respect to $\tau_\alpha$ and using
$(\partial D/\partial\tau_\alpha) \cdot (\partial D/\partial\tau_\beta)
+ (\partial D/\partial\tau_\beta) \cdot (\partial D/\partial\tau_\alpha)
= -D \cdot (\partial^2 D/\partial\tau_\alpha\partial\tau_\beta)
 - (\partial^2 D/\partial\tau_\alpha\partial\tau_\beta) \cdot D$
shows that
$[\partial D/\partial\tau_\alpha, \partial D/\partial\tau_\beta]$
is a $[D,-]$-coboundary in $\End(\barB)$. In particular
it vanishes as an endomorphism of $\barB$
(both sides of the Leibniz identity act on the same complex),
so $\nabla_{\mathrm{rel}}^{\,2} = 0$.

\emph{Part~(ii):} Near a boundary divisor $\delta_\Gamma$
parametrising a nodal degeneration $\Sigma_g \leadsto \Sigma_\Gamma$,
the prime form $E(z,w)$ acquires a logarithmic singularity of the
form $E(z,w) \sim t_e \cdot E_{\mathrm{reg}}(z,w)$, where $t_e$
is the smoothing parameter for the node~$e$ and $E_{\mathrm{reg}}$
is regular. The bar differential therefore has a simple pole
$D \sim D_{\mathrm{reg}} + (1/t_e)\Delta_e$ where $\Delta_e$ is the
clutching operator contracting the two half-edges at the node via the
invariant pairing. The residue of $\nabla_{\mathrm{rel}}$ is thus
$\Delta_\Gamma = \sum_{e \in E(\Gamma)} \Delta_e$.

\emph{Part~(iii):} The monodromy around $\{t_e = 0\}$ is
$\exp(2\pi i \, \operatorname{Res}_{t_e=0} \nabla_{\mathrm{rel}})
= \exp(2\pi i \, \Delta_e)$. For the full stable graph~$\Gamma$,
the monodromy is the product over edges.
\end{proof}

\begin{construction}[Non-separating clutching map]
\label{constr:nonsep-clutching}
\index{clutching map!non-separating}
The non-separating clutching map $\Delta_{\mathrm{cyc}}$ raises
genus by~$1$. On a bar element of degree~$2$:
\begin{equation}\label{eq:nonsep-clutching-explicit}
 \Delta_{\mathrm{cyc}}\bigl(
 s^{-1}a_1 \otimes s^{-1}a_2\bigr)
 \;=\;
 \langle a_1, a_2 \rangle_{\mathrm{inv}},
\end{equation}
where $\langle{-},{-}\rangle_{\mathrm{inv}}$ is the invariant bilinear
form on~$\cA$ (the Shapovalov form, normalised so that the vacuum
pairs to~$1$). Geometrically: the two points $z_1, z_2 \in \Sigma_g$
are identified by the clutching map
$\xi\colon \Sigma_g \smallsetminus \{z_1, z_2\} \to \Sigma_{g+1}$,
which creates a new cycle through the identified point. The bar
element $[s^{-1}a_1 | s^{-1}a_2]$ of bar degree~$2$ at genus~$g$
is sent to the scalar $\langle a_1, a_2\rangle_{\mathrm{inv}}$ at
genus~$g+1$. The desuspension shifts cancel:
$|s^{-1}a_1| + |s^{-1}a_2| = |a_1| - 1 + |a_2| - 1$, and the
pairing has degree~$2$ (it is the bilinear form on the desuspended
augmentation ideal), so
$|\Delta_{\mathrm{cyc}}(s^{-1}a_1 \otimes s^{-1}a_2)| =
|a_1| + |a_2| - 2 + 2 = |a_1| + |a_2|$, consistent with the
genus shift raising the total degree by the Euler characteristic
contribution $\chi(\Sigma_{g+1}) - \chi(\Sigma_g) = -2$.

More generally, on a bar element of degree~$k$,
$\Delta_{\mathrm{cyc}}$ sums over all pairs $(i,j)$ with
$1 \leq i < j \leq k$:
\begin{equation}\label{eq:nonsep-clutching-general}
 \Delta_{\mathrm{cyc}}\bigl(
 s^{-1}a_1 \otimes \cdots \otimes s^{-1}a_k\bigr)
 \;=\;
 \sum_{1 \leq i < j \leq k}
 \pm\,
 \langle a_i, a_j \rangle_{\mathrm{inv}}
 \cdot
 s^{-1}a_1 \otimes \cdots
 \widehat{s^{-1}a_i} \cdots
 \widehat{s^{-1}a_j} \cdots
 \otimes s^{-1}a_k,
\end{equation}
where the hat denotes omission and the sign is determined by the
Koszul rule for permuting $s^{-1}a_i$ and $s^{-1}a_j$ past their
neighbours (see Remark~\ref{rem:bicomplex-signs} below). This is the
genus-raising differential $D_1 = D_{\mathrm{nsep}}$ of the modular
bar complex $\Bmod(\cA)$.
\end{construction}

\begin{proposition}[Explicit propagator on $\Sigma_g$]
\label{prop:propagator-explicit}
\index{propagator!on $\Sigma_g$}
\ClaimStatusProvedHere{}
The holomorphic propagator on a genus-$g$ Riemann surface
$\Sigma_g$ with period matrix $\tau = (\tau_{\alpha\beta})
\in \mathfrak{H}_g$ \textup{(}the Siegel upper half-space\textup{)}
is
\begin{equation}\label{eq:hol-propagator-theta}
 K^{(g)}(z,w)
 \;=\;
 \partial_z \log E(z,w)\, dz
 \;=\;
 \partial_z \log
 \frac{\vartheta\!\begin{bmatrix}\delta'\\\delta''\end{bmatrix}
 \!\bigl(\int_w^z \vec\omega\,,\,\tau\bigr)}
 {\cE(z,w)}\; dz,
\end{equation}
where $E(z,w)$ is the prime form,
$\vartheta\!\left[\begin{smallmatrix}\delta'\\\delta''\end{smallmatrix}\right]$
is a theta function with odd characteristic
$(\delta',\delta'')$, $\vec\omega = (\omega_1,\ldots,\omega_g)$
is the vector of normalised holomorphic differentials, and
$\cE(z,w) = \sqrt{dz}\sqrt{dw}$ is the $(-\frac{1}{2}, -\frac{1}{2})$
bidifferential making $E$ a section of
$\omega_X^{-1/2} \boxtimes \omega_X^{-1/2}$ on $X \times X$.

The Arakelov-normalised single-valued propagator is
\begin{equation}\label{eq:arakelov-propagator}
 K^{(g)}_{\mathrm{Ar}}(z,w)
 \;=\;
 \partial_z \log E(z,w)\, dz
 \;+\;
 \pi \sum_{\alpha,\beta=1}^g
 \omega_\alpha(z)\,
 (\operatorname{Im}\tau)^{-1}_{\alpha\beta}\,
 \operatorname{Im}\Bigl(\int_w^z \omega_\beta\Bigr).
\end{equation}
\end{proposition}
\begin{proof}
The prime form $E(z,w)$ is the unique section of
$\omega_X^{-1/2} \boxtimes \omega_X^{-1/2}(\Delta)$ vanishing to
first order on the diagonal $\Delta$, with no other zeros on
$X \times X$ (Fay \cite{Fay73}, Chapter~II). Its logarithmic
derivative $\partial_z \log E(z,w)\,dz$ is therefore a meromorphic
$1$-form in~$z$ with a simple pole at $z = w$ with residue $+1$
and no other poles. On the universal cover of~$\Sigma_g$ it is
single-valued, but on $\Sigma_g$ it has monodromy: analytic
continuation of $z$ around the $B$-cycle~$B_j$ produces
$K^{(g)}(z + B_j, w) = K^{(g)}(z,w) + 2\pi i\, \omega_j(w)$,
as stated in \eqref{eq:propagator-monodromy}. The $A$-cycle
monodromy vanishes by the normalisation $\oint_{A_i}\omega_j = \delta_{ij}$.

To construct a single-valued propagator, add the Arakelov correction:
the function
$\operatorname{Im}\!\bigl(\int_w^z \omega_\beta\bigr)$
shifts by $\operatorname{Im}(\tau_{j\beta})$ when $z$ traverses
$B_j$. The correction term
$\pi\sum_{\alpha,\beta} \omega_\alpha(z)\,
(\operatorname{Im}\tau)^{-1}_{\alpha\beta}\,
\operatorname{Im}\!\bigl(\int_w^z \omega_\beta\bigr)$
therefore shifts by
$\pi \sum_\alpha \omega_\alpha(z)\,
(\operatorname{Im}\tau)^{-1}_{\alpha j}\,
\operatorname{Im}(\tau_{j\beta})\cdot
(\operatorname{Im}\tau)^{-1}_{\beta\gamma}$,
but more directly: under $z \mapsto z + B_j$, the
correction shifts by $2\pi i\, \omega_j(w)$, cancelling the
monodromy of the prime form exactly.

The cost is that $K^{(g)}_{\mathrm{Ar}}$ is no longer holomorphic
in~$z$. Taking $\bar\partial_z$:
\begin{align}\label{eq:dbar-arakelov}
 \bar\partial_z K^{(g)}_{\mathrm{Ar}}(z,w)
 &\;=\;
 \pi \sum_{\alpha,\beta}
 \bar\partial_z\bigl[\omega_\alpha(z)\bigr]\,
 (\operatorname{Im}\tau)^{-1}_{\alpha\beta}\,
 \operatorname{Im}\Bigl(\int_w^z \omega_\beta\Bigr)
 \;+\;
 \pi \sum_{\alpha,\beta}
 \omega_\alpha(z)\,
 (\operatorname{Im}\tau)^{-1}_{\alpha\beta}\,
 \bar\partial_z
 \operatorname{Im}\Bigl(\int_w^z \omega_\beta\Bigr)
 \notag\\[3pt]
 &\;=\;
 0 \;+\;
 \frac{\pi}{2i}
 \sum_{\alpha,\beta}
 \omega_\alpha(z)\,
 (\operatorname{Im}\tau)^{-1}_{\alpha\beta}\,
 \bar\omega_\beta(z)
 \;=\;
 -\omega_{\mathrm{Ar}}(z),
\end{align}
where the first term vanishes because $\omega_\alpha(z)$ is
holomorphic ($\bar\partial_z\omega_\alpha = 0$), and the second
term uses $\bar\partial_z\operatorname{Im}(\int_w^z\omega_\beta) =
\frac{1}{2i}\bar\omega_\beta(z)$. The Arakelov $(1,1)$-form
$\omega_{\mathrm{Ar}} = -\frac{\pi}{2i}\sum_{\alpha,\beta}
(\operatorname{Im}\tau)^{-1}_{\alpha\beta}\,
\omega_\alpha \wedge \bar\omega_\beta$ is the
unique $(1,1)$-form on~$\Sigma_g$ with
$\int_{\Sigma_g}\omega_{\mathrm{Ar}} = 1$ and constant
pointwise norm with respect to the Arakelov metric.

Contracting with OPE data: the square $\dfib^{\,2}$ of the
bar differential on a degree-$k$ bar element involves two
propagator insertions. At the cogenerator projection, only
the self-contraction (tadpole) survives, extracting
$\kappa(\cA)$ \textup{(}Volume~I, Theorem~D at genus~$1$ and on the
scalar higher-genus lane\textup{)}. Therefore
$\dfib^{\,2} = \kappa(\cA) \cdot \omega_g$.
\end{proof}

\begin{remark}[Comparison with Lurie HA \S5.5]
\label{rem:lurie-comparison}
\index{Lurie!factorization on stratified spaces}
The half-space $\Sigma_g \times \R_{\geq 0}$ with boundary
$\Sigma_g \times \{0\}$ is a stratified space in the sense of
Lurie \cite{HA}, Chapter~5.5, with two strata: the open stratum
$\Sigma_g \times \R_{>0}$ (the ``bulk'') and the closed stratum
$\Sigma_g \times \{0\}$ (the ``boundary''). Lurie's
factorization algebras on stratified spaces
(\emph{loc.\ cit.}, Definition~5.5.2.9) assign to each stratum
a factorization algebra, with compatibility data along the
inclusion of strata. Our BD factorization Swiss-cheese algebra
(Definition~\ref{def:BD-SC}) is the product-split refinement of such
an object:
\begin{itemize}
\item The open stratum carries $F_{\mathsf{cl}} \boxtimes
 F_{\mathsf{op}}$ (factorizable $D$-module on the bulk).
\item The closed stratum carries $F_{\mathsf{op}}$ (the boundary
 algebra).
\item The compatibility datum is the mixed sector
 $F_{\mathsf{mix}}$, encoding how bulk observables approaching the
 boundary act on boundary observables.
\item The directionality (Definition~\ref{def:BD-SC}(iv)) is the
 choice of the open--closed exit indexing category: collar disks in
 the bulk may limit to the boundary and give operations with boundary
 output, while the category has no morphism with boundary input and
 bulk output.
\end{itemize}
Lurie's constructibility supplies local constancy on strata and
compatibility along the frontier; it does not by itself identify the
two-coloured Swiss-cheese object with a single-colour $E_3$-algebra,
nor does it reverse the bulk-to-boundary direction. The content of
this chapter beyond Lurie's general framework is the \emph{specific}
product decomposition
$\Sigma_g \times \R$, which splits the factorization structure
into holomorphic ($\Sigma_g$) and topological ($\R$) colours.
This splitting is \emph{not} part of Lurie's stratified
factorization formalism; it is the extra structure that makes
Swiss-cheese two-coloured rather than merely stratified.
\end{remark}


%%%%%%%%%%%%%%%%%%%%%%%%%%%%%%%%%%%%%%%%%%%%%%%%%%%%%%%%%%%%%%%%%%%%%%%%%%%%%
\section{CG factorization Swiss-cheese algebras}
\label{sec:CG-factorization-SC}
%%%%%%%%%%%%%%%%%%%%%%%%%%%%%%%%%%%%%%%%%%%%%%%%%%%%%%%%%%%%%%%%%%%%%%%%%%%%%

The Costello--Gwilliam framework \cite{CG17} approaches the same
structure from the differential-geometric side: prefactorization
algebras on $\Sigma_g \times \R$, quantized via the BV formalism.

\begin{definition}[CG prefactorization algebra on $\Sigma_g \times \R$]
\label{def:CG-prefact}
\index{CG prefactorization algebra|textbf}
A \emph{CG prefactorization algebra of HT type on
$\Sigma_g \times \R$} is a prefactorization algebra $\Obs$ on the
product manifold $\Sigma_g \times \R$ with values in cochain
complexes, satisfying:
\begin{enumerate}[label=\textup{(\roman*)}]
\item \textbf{Holomorphic factorization in $\Sigma_g$.} For
 disjoint open discs $D_1,\ldots,D_k \subset \Sigma_g$ and
 an interval $I \subset \R$, the multiplication map
 \[
 \bigotimes_{i=1}^k \Obs(D_i \times I)
 \;\longrightarrow\;
 \Obs\Bigl(\bigl(\textstyle\bigcup_i D_i\bigr) \times I\Bigr)
 \]
 is holomorphic in the positions of the discs, with meromorphic
 singularities (OPE poles) along collision diagonals.

\item \textbf{Local constancy in $\R$.} For a disc
 $D \subset \Sigma_g$ and inclusions of intervals
 $I \hookrightarrow J$ in~$\R$, the restriction map
 $\Obs(D \times I) \to \Obs(D \times J)$ is a
 quasi-isomorphism. The assignment $I \mapsto \Obs(D \times I)$
 is therefore a locally constant factorization algebra on~$\R$,
 yielding an $\Eone$-algebra structure.

\item \textbf{Mixed structure.} The holomorphic and topological
 factorizations are compatible: for rectangles
 $D_1 \times I_1, \ldots, D_k \times I_k$ contained in a larger
 rectangle $D \times I$, the iterated multiplication factors
 through the product $\FM_k(\Sigma_g) \times \Eone(m)$ of
 holomorphic and topological operation spaces.

\item \textbf{Directionality.} The half-space
 $\Sigma_g \times \R_{\geq 0}$ carries a boundary condition at
 $\Sigma_g \times \{0\}$. Bulk observables approaching the
 boundary act on boundary observables; there is no structure map
 with boundary inputs and bulk output.
\end{enumerate}
\end{definition}

\begin{definition}[BV quantization origin]
\label{def:BV-origin}
\index{BV quantization!factorization origin}
The standard source of CG prefactorization algebras of HT type is
the BV quantization of a classical action functional on
$\Sigma_g \times \R$. Let $\fg$ be a Lie algebra with
non-degenerate invariant pairing $\langle{-},{-}\rangle$, and
let $\cE = \Omega^{0,\bullet}(\Sigma_g) \otimes \fg \otimes
\Omega^\bullet(\R)$ be the space of fields. Write
$\cE = \cE^{\mathrm{hol}} \oplus \cE^{\mathrm{top}}$ for the
decomposition by form degree on $\R$, with $\phi \in \cE$ a
general field. The \emph{classical action functional}
$S = S_{\mathrm{free}} + S_{\mathrm{int}}$ has:

\smallskip\noindent
\emph{Free (kinetic) part.} The kinetic term is
\begin{equation}\label{eq:HT-kinetic}
 S_{\mathrm{free}}(\phi)
 \;=\;
 \int_{\Sigma_g \times \R}
 \bigl\langle \phi,\;
 (\bar\partial_{\Sigma} + \partial_t)\, \phi
 \bigr\rangle\,
 dz \wedge d\bar{z} \wedge dt,
\end{equation}
where $\bar\partial_{\Sigma}$ is the Dolbeault differential on
$\Sigma_g$ and $\partial_t$ is the de~Rham differential on~$\R$.
The total kinetic operator $Q_0 = \bar\partial_{\Sigma} + \partial_t$
splits along the product: $\bar\partial_\Sigma$ acts in the
holomorphic direction, $\partial_t$ in the topological direction.

\smallskip\noindent
\emph{Interaction part.} The cubic interaction is
\begin{equation}\label{eq:HT-interaction}
 S_{\mathrm{int}}(\phi)
 \;=\;
 \frac{1}{2}
 \int_{\Sigma_g \times \R}
 \bigl\langle \phi,\;
 [\phi, \phi]
 \bigr\rangle\,
 dz \wedge d\bar{z} \wedge dt,
\end{equation}
where $[{-},{-}]$ is the Lie bracket on~$\fg$. The full BRST
differential is
\begin{equation}\label{eq:BRST-splitting}
 Q
 \;=\;
 \bar\partial_{\Sigma}
 \;+\;
 \partial_t
 \;+\;
 \{S_{\mathrm{int}},{-}\}_{\mathrm{BV}},
\end{equation}
where $\{{-},{-}\}_{\mathrm{BV}}$ is the BV antibracket. The
identity $Q^2 = 0$ is the classical master equation
$\{S,S\}_{\mathrm{BV}} = 0$.

\smallskip\noindent
\emph{Propagator.} The Green's function for $Q_0$ on
$\Sigma_g \times \R$ factors as
\begin{equation}\label{eq:HT-propagator}
 P(z_1, t_1;\, z_2, t_2)
 \;=\;
 K^{(g)}(z_1, z_2)
 \;\cdot\;
 \theta(t_1 - t_2),
\end{equation}
where $K^{(g)}(z_1,z_2) = \partial_{z_1}\log E(z_1,z_2)\,dz_1$
is the meromorphic (Cauchy/Szeg\H{o}) kernel on $\Sigma_g$
(at genus~$0$: $K^{(0)} = dz_1/(z_1 - z_2)$) and
$\theta(t_1 - t_2)$ is the Heaviside step function on~$\R$,
encoding causal (retarded) ordering in the topological direction.
The factored form $P = K^{(g)} \cdot \theta$ reflects the
HT product structure: holomorphic meromorphy in $\Sigma_g$
and exponential decay in~$\R$.

\smallskip\noindent
\emph{Cosheaf descent.}
Passage from prefactorization to genuine factorization via the Weiss
topology is ensured by the Weiss descent lemma
(Lemma~\ref{lem:product-weiss-descent}).
\end{definition}

\begin{remark}[Comparison with Costello's renormalization framework]
\label{rem:costello-renorm-comparison}
\index{Costello!renormalization comparison}
The propagator $P = K^{(g)} \cdot \theta$ of
\eqref{eq:HT-propagator} differs from the heat kernel regularization
$P_\varepsilon^L = \int_\varepsilon^L K_t\, dt$ of Costello's
renormalization programme \cite{Cos11}, and the difference is
substantive.
Costello regularises the propagator by smearing the Green's
function over a time window $[\varepsilon, L]$, producing a smooth
kernel; the renormalization group flow tracks the
$\varepsilon$-dependence, and the $\varepsilon \to 0$ limit exists
after counterterm subtraction. In the holomorphic-topological
setting, the $\Sigma_g$-component of the propagator is \emph{already}
meromorphic (not distributional): $K^{(g)}(z,w) =
\partial_z\log E(z,w)\,dz$ has a simple pole at $z = w$ and no
further singularities. The $\R$-component $\theta(t_1 - t_2)$ is
distributional (a step function), but the $\Eone$-direction is
purely topological and requires no regularization; it simply
enforces time-ordering.

Concretely: Costello's heat kernel on a Riemannian manifold~$M$ is
the distributional kernel of $e^{-t\Delta_M}$, where $\Delta_M$ is
the Laplacian. On $\C$ with flat metric,
$K_t(z,w) = \frac{1}{4\pi t}e^{-|z-w|^2/4t}$, and
$\lim_{t \to 0^+} K_t(z,w) = \delta(z-w)$.
In the holomorphic limit ($\bar\partial$-complex rather than the
Laplacian), the heat kernel degenerates to the Cauchy kernel:
formally $K^{(0)}(z,w) = \frac{dz}{z-w}$ is the $t = 0$ limit
of the heat kernel, already finite. The HT propagator
\eqref{eq:HT-propagator} is therefore the \emph{holomorphic limit}
of Costello's regularised propagator: no smearing is needed in the
$\C$-direction, and the $\R$-direction uses time-ordering rather than
heat-kernel smearing. The one-loop finiteness condition
(hypothesis~(H1) of the physics bridge theorem,
Theorem~\ref{thm:physics-bridge}) replaces Costello's
renormalization group analysis: it is the statement that the
holomorphic limit is finite without counterterms beyond one loop.
\end{remark}

\begin{construction}[Factorization homology and the genus tower]
\label{constr:fact-homology-genus}
\index{factorization homology!genus tower}
The \emph{factorization homology} of $\Obs$ on $\Sigma_g \times \R$ is
the global sections
\begin{equation}\label{eq:fact-homology}
 \int_{\Sigma_g \times \R} \Obs
 \;:=\;
 \Obs(\Sigma_g \times \R),
\end{equation}
computed as the homotopy colimit of the factorization structure over
all open covers. For varying~$g$, this produces a family
$\{\int_{\Sigma_g \times \R} \Obs\}_{g \geq 0}$ parametrised by the
complex structure of~$\Sigma_g$ (i.e.\ by points of $\Mbar_g$).

The genus-$g$ bar complex of~$\cA$ is recovered as:
\[
 \barB^{(g)}(\cA)
 \;\simeq\;
 \int_{\Sigma_g \times \R} \Obs_\cA,
\]
where $\Obs_\cA$ is the CG prefactorization algebra associated to the
chiral algebra~$\cA$ via BV quantization. The dependence on complex
structure produces the genus tower of
$\S$\ref{subsec:genus-completion} and the curved bar complex of
Theorem~\ref{thm:modular-bar}.
\end{construction}

\begin{remark}[CG Koszul duality]
\label{rem:CG-koszul}
\index{CG Koszul duality|textbf}
The Ayala--Francis framework \cite{AF15} establishes Koszul duality
for locally constant factorization algebras on~$\R^n$: it exchanges
$E_n$-algebras and $E_n$-coalgebras. In the HT product
$\Sigma_g \times \R$, this theorem applies directly only to the
locally constant topological factor. The holomorphic factor is
chiral, not locally constant, and is governed by Beilinson--Drinfeld
factorization duality. Thus the two inputs are:
\begin{itemize}
\item In the $\R$-direction: $\Eone$ Koszul duality (bar-cobar for
 associative algebras).
\item In the $\Sigma_g$-direction: chiral Koszul duality
 (bar-cobar for chiral/factorization algebras, Volume~I
 Theorem~A).
\item Mixed direction: compatibility of these two dualities across
 the bulk-to-boundary module is the mixed factorization convergence
 hypothesis isolated in Theorem~\ref{thm:factorization-SC-koszul};
 it is not a formal consequence of Lurie's locally constant theorem.
\end{itemize}
The bar construction of Construction~\ref{constr:bar-fact-coalgebra}
is the factorization-level object on which these two dualities meet.
\end{remark}


%%%%%%%%%%%%%%%%%%%%%%%%%%%%%%%%%%%%%%%%%%%%%%%%%%%%%%%%%%%%%%%%%%%%%%%%%%%%%
\section{Equivalence of BD and CG descriptions}
\label{sec:BD-CG-equivalence}
%%%%%%%%%%%%%%%%%%%%%%%%%%%%%%%%%%%%%%%%%%%%%%%%%%%%%%%%%%%%%%%%%%%%%%%%%%%%%

The two descriptions (algebraic: $D$-modules on the Ran space;
differential-geometric: prefactorization algebras) produce
equivalent structures under standard comparison hypotheses.

\begin{theorem}[BD--CG comparison and local shadow; \ClaimStatusProvedElsewhere{} for the BD--CG comparison, \ClaimStatusProvedHere{} for the mixed local-shadow functor]
\label{thm:BD-CG-equivalence}
\index{BD--CG equivalence|textbf}
Let $k$ be a field of characteristic~$0$, $X$ a smooth projective
curve over~$k$, and $\cA$ a chiral algebra on~$X$ that is
holomorphically translation-invariant (in the sense of
\textup{\cite{BD04}, \S3.4}). Then:
\begin{enumerate}[label=\textup{(\roman*)}]
\item \textbf{BD--CG comparison.} The category of
 factorizable $D$-modules on $\Ran(X)$ that are $\cA$-modules
 \textup{(}Beilinson--Drinfeld \textup{\cite{BD04}},
 Theorem~\textup{4.5.4; Francis--Gaitsgory \cite{FG12})}
 is equivalent, under the standard analytic comparison hypotheses,
 to the category of holomorphic
 prefactorization algebras on~$X$ with values in ind-coherent
 sheaves, satisfying cosheaf descent in the Weiss topology
 \textup{(}Costello--Gwilliam \textup{\cite{CG17})}.

\item \textbf{Local shadow.} Formal completion at collision points
 defines a functor from the BD/CG factorization category to algebras
 over
 $C_\ast(\FM_\bullet(T_xX))$ in the completed symmetric monoidal dg
 category at~$x$ \textup{(}the closed-colour component of
 a $C_\ast(W(\SCchtop))$-algebra\textup{)}.

\item \textbf{Scope of equivalence.} The local-shadow functor is an
 equivalence only on the formal-disc, translation-invariant
 genus-$0$ subcategory obtained from the affine chart
 $\mathbb P^1\setminus\{\infty\}\simeq \mathbb C$ with extension
 across $\infty$ already supplied. For arbitrary global
 factorizable $D$-modules on $\mathbb P^1$, extension across
 $\infty$ is additional descent data. At genus~$g\ge 1$ the
 functor is not an equivalence: it forgets period data, family
 monodromy, and the scalar anomaly $\kappa(\cA)\cdot\omega_g$.
\end{enumerate}
\end{theorem}

\begin{remark}[Proof inputs for the equivalence]
\label{rem:bdcg-proof-inputs}
The three layers of the proof invoke:
\begin{enumerate}[label=\textup{(\roman*)}]
\item Beilinson--Drinfeld \cite{BD04} and Francis--Gaitsgory \cite{FG12} for the chiral/closed equivalence.
\item Lurie \cite{HA}, Theorem~5.4.5.9, for the associative/open equivalence (locally constant factorization algebras on~$\R$ are $\Eone$-algebras).
\item The mixed local-shadow comparison (Layer~3) is proved here:
the Riemann--Hilbert comparison functor $\Phi_{\mathsf{mix}}$ is
constructed explicitly and verified to intertwine the factorization
structure maps after formal completion at collision points.
\end{enumerate}
The theorem is conditional on~(i) and~(ii), which are standard
foundational results proved in the cited references. The novel content
is~(iii): the mixed-sector local-shadow comparison and its compatibility
with Swiss-cheese directionality. No assertion here makes the local
operadic model equivalent to global Ran-space factorization at
genus~$g\ge 1$.
\end{remark}

\begin{proof}
The proof proceeds through three layers of comparison.

\medskip\noindent
\textbf{Layer 1: Closed (chiral) colour.}
Beilinson--Drinfeld \cite{BD04}, Theorem~4.5.4, establish that
the category of chiral algebras on~$X$ (as factorizable
$D$-modules on $\Ran(X)$) is equivalent to the category of
factorization algebras on~$X$ in the sense of \emph{loc.\ cit.},
\S3.4. Francis--Gaitsgory \cite{FG12} extend this to the
derived/$\infty$-categorical setting. This is a global Ran-space
statement. Passing from it to algebras over
$C_\ast(\FM_\bullet(T_xX))$ is the separate formal-completion
functor at~$x$; Francis--Gaitsgory do not identify arbitrary global
factorization $D$-modules with their local operadic shadows. The
Costello--Gwilliam prefactorization algebra on~$X$, after
passage to the cosheaf topology, is equivalent to the
Francis--Gaitsgory factorization algebra by the comparison
of~\cite{CG17}, \S5.5. This establishes the BD--CG comparison on
the closed colour. Passing from this factorization category to the
operadic local model is a separate formal-completion functor; it is an
equivalence only in the genus-$0$ affine-chart case recorded in the
theorem.

\medskip\noindent
\textbf{Layer 2: Open (associative) colour.}
Lurie \cite{HA}, Theorem~5.4.5.9, identifies locally
constant factorization algebras on~$\R$ with $\Eone$-algebras
(associative algebras up to coherent homotopy) in any presentable
symmetric monoidal $\infty$-category. Both the BD and CG
descriptions of the open colour reduce to $\Eone$-algebras:
the BD description via the locally constant $D$-module on
$\Ran(\R)$, the CG description via the locally constant
prefactorization algebra on~$\R$. This is Lurie's result applied
in the two settings.

\medskip\noindent
\textbf{Layer 3: Mixed (Swiss-cheese) colour.}
The mixed sector involves configurations with both closed points
on~$X$ and open points on~$\R$. We construct an explicit
comparison functor
\[
 \Phi_{\mathsf{mix}}\colon
 \Dmod\bigl(\Ran_k(X)\bigr) \otimes \Fact_m(\R)
 \;\longrightarrow\;
 \Obs^{\mathsf{CG}}_{\mathsf{mix}}(k,m)
\]
as follows. A section $s \in F_{\mathsf{mix}}(k,m)$
(Definition~\ref{def:BD-SC}(iii)) is a $D$-module section on
$\Ran_k(X)$ tensored with a factorization-algebra element on
$\Fact_m(\R)$. The functor $\Phi_{\mathsf{mix}}$ sends $s$ to
the CG observable defined by:
\begin{enumerate}[label=\textup{(\alph*)}]
\item Apply the Riemann--Hilbert correspondence on $\Ran_k(X)$:
 convert the $D$-module section to a flat section of the
 associated local system (the de~Rham complex of the
 $D$-module), producing a differential form on
 $\Conf_k(X)$.
\item On the $\R$-factor, apply the identity: both BD and CG
 descriptions of locally constant factorization algebras on
 $\R$ produce $\Eone$-algebras (Layer~2), and the
 identification is Lurie's equivalence.
\item The combined map sends $s$ to the product
 $\mathrm{dR}(s|_{\Ran_k}) \boxtimes s|_{\Fact_m}$, an
 element of the CG mixed observables on $\Conf_k(X) \times
 \Conf_m(\R)$.
\end{enumerate}
After formal completion at the mixed collision stratum,
$\Phi_{\mathsf{mix}}$ is a quasi-isomorphism on the completed local
mixed complexes by the K\"unneth quasi-isomorphism
\[
 C_\ast\bigl(\FM_k(X) \times \Eone(m)\bigr)
 \;\simeq\;
 C_\ast\bigl(\FM_k(X)\bigr) \otimes C_\ast\bigl(\Eone(m)\bigr),
\]
which holds because $\FM_k(X)$ has finite-dimensional cohomology
over a field of characteristic~$0$ (it is a compact manifold
with corners), and $\Eone(m)$ has finite-dimensional cells.
The completed factorization module structure maps
(Definition~\ref{def:BD-SC}(iii)(a)--(c)) are intertwined by
$\Phi_{\mathsf{mix}}$ with the CG mixed multiplication maps
(Definition~\ref{def:CG-prefact}(iii)), because both are
determined, after completion, by restriction to collision diagonals in
$\FM_k(X)$ tensored with interval inclusions in $\Eone(m)$.
For disjoint $I_1, I_2$, the diagram
\[
\begin{tikzcd}[column sep=3em]
F_{\mathsf{mix}}\big|_{U_{I_1,I_2}} \ar[r, "\Phi_{\mathsf{mix}}"] \ar[d, "\text{fact.\ iso}"'] &
\Obs^{\mathsf{CG}}_{\mathsf{mix}}\big|_{U_1 \times U_2} \ar[d, "\text{CG mult}"] \\
F_{\mathsf{cl}}\big|_{U_{I_1}} \boxtimes F_{\mathsf{op}}\big|_{U_{I_2}} \ar[r, "\Phi_{\mathsf{cl}} \boxtimes \Phi_{\mathsf{op}}"] &
\Obs^{\mathsf{CG}}_{\mathsf{cl}}(U_1) \otimes \Obs^{\mathsf{CG}}_{\mathsf{op}}(U_2)
\end{tikzcd}
\]
commutes because $\Phi_{\mathsf{mix}}$ is defined component-wise on the two factors and both vertical maps are the ``forget the separation'' isomorphism.

The no-open-to-closed constraint is preserved: the vanishing
$F_{\mathsf{mix}}(\ldots,\mathsf{top},\ldots;\,\mathsf{cl}) = 0$
on the BD side (the empty output component in the open--closed
indexing category of Definition~\ref{def:BD-SC}(iv)) corresponds to
the absence of bulk-directed CG observables sourced from the boundary
(Definition~\ref{def:CG-prefact}(iv)): $\Phi_{\mathsf{mix}}$
sends the zero $D$-module to the zero observable.
\end{proof}

\begin{remark}[Precise adjectives]
\label{rem:BD-CG-adjectives}
The equivalence requires:
\begin{itemize}
\item \emph{Characteristic~$0$}: needed for Kontsevich formality
 (Step~2 of Theorem~\ref{thm:homotopy-Koszul}) and for the
 K\"unneth quasi-isomorphism.
\item \emph{$X$ smooth projective}: ensures the Ran space has the
 correct properties; stable curves require additional care at
 nodes.
\item \emph{Holomorphically translation-invariant}: the CG
 prefactorization algebra is translation-invariant in the
 holomorphic direction, matching the BD condition that the
 $D$-module connection is compatible with the translation action
 on~$\C$.
\item \emph{Locally constant on~$\R$}: the topological
 factorization is locally constant (restriction along inclusions
 of intervals is a quasi-isomorphism), ensuring the
 $\Eone$-algebra structure.
\item \emph{Ind-coherent on BD side, cosheaf topology on CG side}:
 the comparison functors involve the passage between
 $\IndCoh(\Ran(X))$ and the cosheaf category.
\end{itemize}
\end{remark}

% Pentagon of equivalences for SC^{ch,top}: the BD/CG local-shadow
% comparison in this section is one edge; the full pentagon coherence
% remains conjectural at chain level.
\begin{remark}[Pentagon of equivalences for $\SCchtop$; \ClaimStatusConjectured]
\label{rem:sc-chtop-pentagon}
Theorem~\ref{thm:BD-CG-equivalence} constructs the local-shadow
comparison from the \emph{factorization} presentation
$\mathsf P_3$ to the \emph{operadic} presentation $\mathsf P_1$ of
the Swiss-cheese chiral-topological coloured operad. Here
$\mathsf P_1$ consists of Voronov generators on
$\FM_k(\C) \times \Conf_m(\R)$ with bulk-to-boundary structure, and
$\mathsf P_3$ consists of BD D-modules on
$\Ran(X)$ with boundary stratum along $\R \subset \C$, acting on the
pair $(\Zderch(A), A)$, never on $A$ alone. At genus~$0$ this
comparison is an equivalence only on the same translation-invariant
affine formal-disc subcategory as Theorem~\ref{thm:BD-CG-equivalence};
for arbitrary global objects it is a local-shadow functor. At
genus~$g\ge 1$ it forgets the global period data described above.
Three further presentations fit into a conjectural pentagon:
\begin{enumerate}[label=(\roman*)]
\item $\mathsf P_2$ (Koszul dual): $(\SCchtop)^!$ with
$\mathsf P_2 \simeq E_2\{1\}$ (Gerstenhaber self-dual up to a shift);
\item $\mathsf P_4$ (BV/BRST): the CG observable coloured algebra
$\Obs^{\mathrm{q}}(U)$ on the half-plane, with QME encoding the
open/closed Maurer-Cartan equation;
\item $\mathsf P_5$ (convolution): the $L_\infty$-convolution
algebra $\mathfrak g^{\SC}_T$ arising from the bar cooperad
$B(\SCchtop)$, via Lurie's HA~5.3.1.30 in the coloured setting.
\end{enumerate}
The pentagon coherence 2-cocycle $\omega \in H^2$ of the diagram
\[
\mathsf P_1 \longleftrightarrow \mathsf P_2 \longleftrightarrow
\mathsf P_3 \longleftrightarrow \mathsf P_4 \longleftrightarrow
\mathsf P_5 \longleftrightarrow \mathsf P_1
\]
is expected to vanish by the coloured extension of Calaque--Willwacher
chiral formality (2015); no chain-level proof is asserted here. The single-colour
projection of the pentagon onto the closed colour recovers the Chiral
Hochschild Trinity of the three models
$C^\bullet_{\mathrm{chiral}}$ / $\mathrm{Ext}^\ast_{A^e}(A,A)$ /
$\mathrm{RHH}_{\mathrm{ch}}(A)$. The bulk-on-boundary
mixed sector realises the topological double cover $B_2 \to S_2$ that
underlies the q-convention bridge $q_{KL}^2 = q_{DK}$.
\end{remark}

% Universal Trace Identity viewed from Vol II's Pentagon. The Pentagon
% coherence [omega]=0 is the operadic substrate making the cross-volume
% identity chain-level coherent; this remark records the specialisation
% visible on the Swiss-cheese side. Upgrade to theorem when the dedicated
\begin{remark}[Pentagon substrate for the Universal Trace Identity; \ClaimStatusConditional]
\label{rem:pentagon-trace-identity}
The Pentagon (Remark~\ref{rem:sc-chtop-pentagon}) is the operadic
substrate of the cross-volume Universal Trace Identity (see the
dedicated universal trace identity chapter)
\[
 \tr_{Z(\mathcal C)}(\mathfrak K_{\mathcal C})
 \;=\; -\,c_{\mathrm{ghost}}(\mathrm{BRST}(\Phi(\mathcal C)))
 \;=\; \tfrac{1}{2}\,c_N(0),
\]
in the following precise sense. The reflection
$\mathfrak K_{\mathcal C}$ on the categorical centre of a CY-$d$
category $\mathcal C$ admits two specialisations, each obtained by
collapsing one colour of $\SCchtop$:
\begin{enumerate}[label=\textup{(\roman*)}]
\item \emph{Closed-colour projection} $\pi_{\mathsf{cl}}$: the bulk
$\mathsf P_3$ recovers the chiral-bar projection, giving
$\pi_{\mathsf{cl}} \mathfrak K_{\mathcal C} = -\,c_{\mathrm{ghost}}(\mathrm{BRST}(\Phi(\mathcal C)))$
(Volume~I K-Atiyah / ghost-identity theorem).
\item \emph{Open-colour projection} $\pi_{\mathsf{op}}$: the boundary
$\mathsf P_3$ recovers the Borcherds singular-theta projection,
giving $\pi_{\mathsf{op}} \mathfrak K_{\mathcal C} = \tfrac12 c_N(0)$
(Volume~III universal BKM-weight proposition).
\end{enumerate}
The Pentagon coherence $[\omega] = 0$ (Calaque--Willwacher chiral
formality in the coloured setting) is precisely the datum certifying
that $\pi_{\mathsf{cl}}$ and $\pi_{\mathsf{op}}$ agree after pullback
along $\Phi$; the two projections are two factorisations of one
reflection, not two independent scalars. Chain-level identification is
conditional on the cross-volume trace-identity conjecture.
The identity is verified numerically on the Vol III $\Z/N$ orbifold
table for $N = 1, \ldots, 8$; the tempting $5=2+3$ equality at $N=1$
is only the K3 mnemonic for fibre Euler characteristic plus the
Heisenberg--Mukai count, not a $\kappa_{\mathrm{BKM}}$ additivity
identity.
\end{remark}

\begin{example}[K3 specialisation of the Pentagon trace substrate]
\label{ex:sc-pentagon-k3-trace}
For the K3 bulk-boundary pair the Pentagon ceases to be an abstract
coherence diagram and becomes a geometric machine. The open colour is
the boundary algebra on the elliptic curve, the closed colour is the
nearby-cycle bulk over the K3 nodal-smoothing host
$\Ran(E^{\mathrm{nod,sm}}_{24}) \times \overline{\mathcal{A}}_2$, and
the number $24$ again comes from the Euler-number count of the elliptic
K3. Under the closed-colour projection the Pentagon gives the trinity
trace $K(\cA)=\mathrm{tr}_{Z(\cA)}(K_\cA)$; under the open-colour
projection it gives the Borcherds quantity
$c_{\phi_{0,1}^{K3}}(0)/2=5$. The same
nearby-cycle package that carries the Humbert wall of
Gritsenko--Nikulin 1998, Theorem~2.1 is therefore the
$\SCchtop$ datum making the K3 Universal Trace Identity look
inevitable rather than accidental.
\end{example}

\begin{remark}[Pentagon closed/open as stage-1/stage-2 of the
Vol.~III two-stage factorisation; \ClaimStatusConditional]
\label{rem:pentagon-two-stage}
\index{Pentagon!closed/open as stage-1/stage-2}%
\index{two-stage factorisation!Pentagon realisation}%
The closed and open colours of the Pentagon
(Remark~\ref{rem:sc-chtop-pentagon}) are structurally the two stages
of the Volume~III factorisation
\[
 \Phi_d \;=\; {\SpCh}_{\Sigma_{d-1},\,C}\circ\PhiFA_d
\]
at $d = 2, 3$. Stage~1 produces a canonical $E_d$-holomorphic
factorisation algebra $\PhiFA_d(\cC)$ on the CY-$d$ target by
Kontsevich--Tamarkin formality plus Costello--Gwilliam--Li
quantisation; restricted to local observables on a formal disc, this
is the \emph{closed colour} $F_{\mathsf{cl}}$ of the Pentagon
(Definition~\ref{def:BD-SC}(i)). Stage~2 is the
factorisation-homology pushforward
\[
 {\SpCh}_{\Sigma_{d-1},C}\bigl(\PhiFA_d(\cC)\bigr)
 \;=\; \int_{\Sigma_{d-1}} \PhiFA_d(\cC)\big|_{\Sigma_{d-1}\times C\times\R_{\geq 0}},
\]
landing on a reference curve $C$ and its boundary half-line
$\R_{\geq 0}\subset C\times\R_{\geq 0}$. Restricted to the
boundary stratum, this is the \emph{open colour} $F_{\mathsf{op}}$
of the Pentagon (Definition~\ref{def:BD-SC}(ii), locally constant
factorisation on $\Ran(\R)$, equivalent to $E_1$ by
Ayala--Francis~\cite{AF15}). The mixed sector $F_{\mathsf{mix}}$
(Definition~\ref{def:BD-SC}(iii)) records the $\int_{\Sigma_{d-1}}$
pushforward as a bulk-to-boundary module, and the directional
restriction of Definition~\ref{def:BD-SC}(iv) (no open-to-closed
operations) is the operadic expression of the fact that stage~2 is
a specialisation of stage~1, never an inversion. A CY-$d$ category
$\cC$ admits a \emph{family} of $E_1$-chiral shadows, parametrised
by the choice of $(\Sigma_{d-1}, C)$: the Pentagon records the
closed colour and one open-colour shadow; the family of open-colour
shadows is the family of reference curves on which
${\SpCh}_{\Sigma_{d-1},C}$ may specialise. The $d = 2$ K3 case has
$\Sigma_1 = \R$ canonically and the family collapses to a single
pair $(\cA_\cC, \mathbf H_{\Delta_5})$; the $d = 3$ case
$K3 \times E$
case has $\Sigma_2$ varying over transverse complex surfaces,
producing a non-trivial family whose single-member specialisation
is the Volume~III $\Phi_3$-image on the reference curve $C = E$.
The arrow of the shadow is from stage~1 (CY) to stage~2 (chiral),
never back.
\end{remark}


%%%%%%%%%%%%%%%%%%%%%%%%%%%%%%%%%%%%%%%%%%%%%%%%%%%%%%%%%%%%%%%%%%%%%%%%%%%%%
\section{Factorization Koszul duality and flat/curved comparison}
\label{sec:factorization-koszul-duality}
%%%%%%%%%%%%%%%%%%%%%%%%%%%%%%%%%%%%%%%%%%%%%%%%%%%%%%%%%%%%%%%%%%%%%%%%%%%%%

The local homotopy-Koszulity theorem
(Theorem~\ref{thm:homotopy-Koszul}) ensures that the bar-cobar
adjunction is a Quillen equivalence on the genus-$0$ formal-disc
model. The factorization framework isolates the all-genus extension:
the bar complex is then a \emph{curved} factorization coalgebra and
the relevant category is the coderived category.

\begin{theorem}[Factorization Swiss-cheese Koszul duality under mixed
factorization convergence; \ClaimStatusProvedHere{} for the local and
categorical parts, \ClaimStatusConditional{} for positive-genus mixed
recovery]
\label{thm:factorization-SC-koszul}
\index{factorization Koszul duality|textbf}
Let $(F_{\mathsf{cl}}, F_{\mathsf{op}}, F_{\mathsf{mix}})$ be a
BD factorization Swiss-cheese algebra on $\Sigma_g \times \R$
\textup{(}Definition~\textup{\ref{def:BD-SC})}, associated to a
chiral algebra~$\cA$ satisfying the hypotheses of
Theorem~\textup{\ref{thm:BD-CG-equivalence}}. Then:
\begin{enumerate}[label=\textup{(\roman*)}]
\item \textbf{Bar construction.} The bar complex
 $\barB(F)$ \textup{(}Construction~\textup{\ref{constr:bar-fact-coalgebra})}
 is a factorization coalgebra over the coloured mixed geometry
 $\Ran^{\mathrm{oc}}_{\HT}(\Sigma_g,I)$, equivalently over its
 open--closed configuration strata
 $\Conf_k(\Sigma_g\setminus D)\times\Conf_m^{\mathrm{ord}}(I)$
 after choosing the boundary interval. No assertion is made that
 the global object is a factorization coalgebra on the literal
 product $\Ran(\Sigma_g)\times\Ran(\R)$.

\item \textbf{Genus-$0$ recovery.} At genus~$0$, the cobar
 construction recovers the original algebra:
 $\Omegach(\barB(F))|_{g=0} \simeq F$. This is the local
 homotopy-Koszulity of $\SCchtop$
 \textup{(}Theorem~\textup{\ref{thm:homotopy-Koszul})}, lifted to
 the factorization level on the translation-invariant affine
 subcategory.

\item \textbf{Positive-genus recovery.} At genus~$g \geq 1$,
 closed-colour Getzler--Kapranov involutivity
 \textup{(}Volume~I, Theorem~\textup{\ref*{V1-thm:feynman-involution})}
 supplies the modular closed-colour input. Recovery of the full
 mixed curved factorization algebra is conditional on a mixed
 all-genus factorization Koszul duality theorem and the
 bounded-filtered Positselski convergence hypotheses.

\item \textbf{Categorical placement.} At each genus~$g$, the
 genus-$g$ associated graded bar complex
 $(\mathrm{gr}^g\, \Bmod, D_0^2 = 0)$ is a strict dg coalgebra
 and therefore belongs to the ordinary derived category,
 while the curved geometric bar complex
 $(\barB^{(g)}_{\mathrm{geom}}, \dfib^{\,2} =
 \kappa \cdot \omega_g)$ belongs to the coderived category.
 The later chiral-exponential comparison relates these flat and
 curved packages on the manuscript's proved comparison surface;
 this categorical distinction is not removed here merely from
 homotopy-Koszulity of the local model.
\end{enumerate}
\end{theorem}

\begin{proof}
\textbf{Part (i)} is the content of
Construction~\ref{constr:bar-fact-coalgebra}: the bar differential
and coproduct are assembled from the factorization data.

\smallskip\noindent
\textbf{Part (ii).} At genus~$0$, the bar complex is a genuine
dg coalgebra ($d_{\barB}^2 = 0$), and the bar-cobar adjunction
$\Omegach \dashv \barB$ restricts to the adjunction of
Theorem~\ref{thm:scchtop-bar-cobar-adjunction}. The counit
$\Omegach(\barB(F)) \to F$ is a quasi-isomorphism by the
homotopy-Koszulity of $\SCchtop$
(Theorem~\ref{thm:homotopy-Koszul}), which ensures that the
bar-cobar resolution is a cofibrant replacement. At the
factorization level, this statement lifts from operations
at formal discs to operations on $\Ran(\C)$, because the
formal-disc data determines the factorization algebra only on the
translation-invariant affine subcategory of $\Ran(\C)$ used here.
For a general factorization $D$-module on $\mathbb P^1$, extension
across $\infty$ is additional descent data, exactly as in
Remark~\ref{rem:local-shadow-extraction}.

\smallskip\noindent
\textbf{Part (iii).} At genus~$g \geq 1$, the bar differential is
curved ($\dfib^{\,2} = \kappa \cdot \omega_g$), and cobar
reconstruction is replaced by the Feynman transform. We recall
the relevant definitions.

The \emph{modular commutative operad} $\Com_{\mathrm{mod}}$ is
the modular operad with
$\Com_{\mathrm{mod}}((g,n)) = \mathbb{k}$ for all stable
pairs $(g,n)$ (i.e.\ $2g - 2 + n > 0$) and structure maps
given by the identity. The \emph{relative Feynman transform}
$\mathrm{FT}(\Com_{\mathrm{mod}} / \SCchtop)$ is the
functor on $(\Com_{\mathrm{mod}}, \SCchtop)$-algebras defined
as follows (see Definition~\ref{def:relative-feynman-transform}
for the general construction):
\[
 \mathrm{FT}_{\Com_{\mathrm{mod}} / \SCchtop}(A)
 \;=\;
 \biggl(
 \bigoplus_{\Gamma \in \StGraph}
 \bigotimes_{v \in V(\Gamma)}
 A((g_v, n_v)) \otimes \orline_\Gamma
 \,,\;
 D_0 + D_1
 \biggr),
\]
where the sum ranges over stable graphs~$\Gamma$ (with
vertices~$v$ labelled by genus~$g_v$ and valence~$n_v$),
$\orline_\Gamma = \det(\mathrm{Edges}(\Gamma))^{\otimes(-1)}$
is the orientation line (the determinant of the edge set, dual),
$D_0 = D_{\mathrm{int}} + D_{\mathrm{sep}}$ encodes
genus-preserving operations (internal differential and
separating-edge contraction, governed by the operad
$\SCchtop$), and
$D_1 = D_{\mathrm{nsep}}$ encodes genus-raising operations
(nonseparating-edge contraction, governed by the modular
structure of $\Com_{\mathrm{mod}}$).

The modular bar complex $\Bmod(\cA)$
(Theorem~\ref{thm:modular-bar}) is a
$\mathrm{FT}(\Com_{\mathrm{mod}} / \SCchtop)$-algebra:
its bicomplex structure $(D_0^2 = 0,\; D_1^2 = 0,\;
D_0 D_1 + D_1 D_0 = 0)$ is exactly the defining data of the
relative Feynman transform
(Proposition~\ref{prop:rft-differential-identification}).

\begin{remark}[Koszul sign rule for the bicomplex]
\label{rem:bicomplex-signs}
\index{Koszul sign rule!bicomplex}
The anticommutativity $D_0 D_1 + D_1 D_0 = 0$ is governed by the
Koszul sign rule applied to the bigrading $(|x|, g(x))$, where
$|x|$ is the bar degree and $g(x)$ is the genus degree. The sign
conventions are as follows. Set $\varepsilon = (-1)^{|x|}$.
\begin{enumerate}[label=\textup{(\roman*)}]
\item $D_0$ preserves genus and has bar-degree sign
 $\varepsilon$: it acts on $x$ by $D_0(x) = \varepsilon \cdot
 (\text{sum of codimension-1 boundary terms in } \Mbar_{g,n})$.
 The sign $\varepsilon$ arises because $D_0$ is a coderivation
 of the tensor coalgebra $T^c(s^{-1}\bar\cA)$, and the
 desuspension introduces a sign $(-1)^{|s^{-1}a_1| + \cdots +
 |s^{-1}a_{i-1}|}$ when $D_0$ acts on the $i$-th tensor factor.

\item $D_1$ raises genus by~$1$ and has bar-degree sign
 $(-1)^{|x|+1}$: the additional sign $(-1)$ relative to $D_0$
 comes from the orientation of the nonseparating boundary stratum.
 Geometrically: the nonseparating clutching map
 $\xi_{\mathrm{nsep}}\colon \Mbar_{g,n+2} \to
 \Mbar_{g+1,n}$ identifies two marked points into a node, and
 the orientation of the normal bundle to this boundary stratum
 differs from the separating case by $(-1)^1$
 (one self-edge versus one separating edge).

\item The anticommutativity follows:
 $D_0 D_1 + D_1 D_0$ acting on~$x$ acquires sign
 $\varepsilon \cdot (-\varepsilon) + (-\varepsilon) \cdot
 \varepsilon = -\varepsilon^2 - \varepsilon^2 = 0$? No:
 the correct derivation uses the orientation lines. On
 $\mathrm{gr}^g\,\Bmod$, $D_0$ has total degree $+1$ and $D_1$
 has total degree $+1$ (both raise cohomological degree by~$1$
 in the totalised complex). Therefore $D_0 D_1 + D_1 D_0 = 0$
 is the statement that two anticommuting differentials of
 degree~$+1$ define a bicomplex. The sign is verified by
 checking on generators: for $x = s^{-1}a$ of bar degree~$1$,
 $D_0(x) = 0$ (no codimension-$1$ faces at arity~$1$) and
 $D_1(x) = 0$ (no pair to contract at arity~$1$), so the
 relation is trivially satisfied. For $x = s^{-1}a_1 \otimes
 s^{-1}a_2$ of bar degree~$2$:
 \begin{align*}
 D_0 D_1(x) &= D_0\bigl(\langle a_1, a_2\rangle_{\mathrm{inv}}\bigr)
 = 0, \\
 D_1 D_0(x) &= D_1\bigl((-1)^{|a_1|-1}\, s^{-1}(a_1 \cdot a_2)\bigr)
 = 0,
 \end{align*}
 where $D_0(x) = (-1)^{|a_1|-1}\,s^{-1}(a_1 \cdot a_2)$ is the
 bar differential (chiral product) and $D_1$ of an arity-$1$ element
 vanishes. The first non-trivial check occurs at bar degree~$3$;
 see Example~\ref{ex:heisenberg-bar-3} below for the Heisenberg
 computation.
\end{enumerate}
For the general sign computation at arbitrary arity, see the
orientations appendix (Appendix~\ref{app:orientations}).
\end{remark}

At $g\geq 1$, closed-colour Getzler--Kapranov involutivity supplies
only the closed modular input
$\mathrm{FT}^2\simeq\id$ on the closed colour
(Volume~I, Theorem~\ref*{V1-thm:feynman-involution}). Full mixed
curved recovery follows only under the named hypotheses of
Theorem~\ref{thm:factorization-SC-koszul}: mixed all-genus
factorization Koszul duality on
$\Ran^{\mathrm{oc}}_{\HT}(X,D,\tau)$, bounded-filtered
Positselski convergence, and the chiral Riemann--Hilbert comparison
identifying the flat prime-form model with the curved Arakelov
representative. Under these hypotheses the factorization bar
construction and its mixed cobar are inverse in the coderived
category. At genus~$0$, the
Feynman transform reduces to the cobar functor $\Omegach$,
and $\mathrm{FT}^2 \simeq \id$ specialises to
Part~(ii). At genus~$g \geq 1$, the Feynman transform
sums over nonseparating edge contractions
(the genus-raising moves), and the involutivity
$\mathrm{FT}^2 \simeq \id$ ensures that this sum produces
a quasi-inverse to the bar construction.
The factorization-level content is that the double Feynman
transform filtration converges (by completeness of the genus
filtration, Theorem~\ref{thm:modular-bar}(iii)) and the
$E_\infty$-page is $F$ itself.

The relative Feynman transform also unifies the three models:
$\mathrm{FT}_{\mathrm{rel}}$ produces the flat model~(i) when
applied in the operadic (local) gauge;
the Fay-trisecant-corrected model~(ii) when the propagator is
promoted to the prime form on the universal cover; and the
curved geometric model~(iii) when the Arakelov normalisation
is used (Corollary~\ref{cor:three-models-from-fact}).

\smallskip\noindent
\textbf{Part (iv).} The three chain-level models of
Remark~\ref{rem:three-models} and
Corollary~\ref{cor:three-models-from-fact} below
provide three presentations of the genus-$g$ bar complex.
The flat model $(\mathrm{gr}^g\, \Bmod, D_0^2 = 0)$ lives in the
derived category; the curved model
$(\barB^{(g)}_{\mathrm{geom}}, \dfib^{\,2} = \kappa \cdot
\omega_g)$ lives in the coderived category. This is the honest
categorical placement proved on this surface. The stronger
flat-derived / curved-coderived comparison is established only
later in the chapter, once the chiral exponential map and the
restricted comparison theorem have been constructed explicitly;
it is not a consequence here merely of completeness of the genus
filtration together with homotopy-Koszulity of the local model.
\end{proof}

\begin{corollary}[Three models from factorization]
\label{cor:three-models-from-fact}
\index{three models!from factorization}
The three chain-level models of the genus-$g$ bar complex
\textup{(}Remark~\textup{\ref{rem:three-models})} are three
\emph{gauges} of the factorization structure:
\begin{enumerate}[label=\textup{(\roman*)}]
\item \textbf{Flat associated graded} ($D_0^2 = 0$): the
 factorization coalgebra in \emph{algebraic gauge}: formal
 completion at collision points, where the local operad
 $\SCchtop$ acts. This is the model that
 Theorem~\textup{\ref{thm:homotopy-Koszul}} governs directly.

\item \textbf{Corrected holomorphic} ($\Dg{g}^2 = 0$): the
 factorization coalgebra in \emph{flat gauge}: on the universal
 cover $\widetilde{\Sigma}_g$ of~$\Sigma_g$, where the holomorphic
 propagator (the prime form $E(z,w)$) defines a flat connection.
 The Fay trisecant identity ensures $\Dg{g}^2 = 0$, and this
 identity is local on configuration spaces (it is visible to the
 operadic model).

\item \textbf{Curved geometric} ($\dfib^{\,2} = \kappa \cdot
 \omega_g$): the factorization coalgebra in \emph{single-valued
 gauge}: on $\Sigma_g$ itself, where the Arakelov propagator
 is single-valued but curved. The curvature is the monodromy of
 the propagator trivialisation from flat to single-valued gauge.
\end{enumerate}
The passage between gauges is the chiral Riemann--Hilbert
correspondence: the gauge transformation from flat to single-valued
gauge is the monodromy representation of the chosen flat
trivialisation, and its projective scalar defect is measured on the
proved lane by $\kappa(\cA) \cdot \omega_g$.
\end{corollary}

\begin{proof}
The three models arise from three presentations of the factorization
coalgebra, corresponding to three choices of trivialisation of the
$D$-module connection:
\begin{itemize}
\item Model~(i): trivialise by passing to formal completions at
 collision points (extracting the local operad);
\item Model~(ii): trivialise by pulling back to the universal cover
 (where the connection is flat by the Fay identity);
\item Model~(iii): use the globally defined single-valued
 (Arakelov) trivialisation, which introduces curvature.
\end{itemize}
That models~(i) and~(ii) have $d^2 = 0$ while model~(iii)
has $\dfib^{\,2} = \kappa \cdot \omega_g$ at genus~$1$, and on the
proved scalar higher-genus lane, is the content of
Volume~I, Convention~3.1 and Theorem~D. The derived-coderived
comparison between models~(i)/(ii) and model~(iii) is developed
later in the chapter; Theorem~\ref{thm:factorization-SC-koszul}(iv)
records only the categorical homes of the flat and curved
presentations.
\end{proof}


%%%%%%%%%%%%%%%%%%%%%%%%%%%%%%%%%%%%%%%%%%%%%%%%%%%%%%%%%%%%%%%%%%%%%%%%%%%%%
\section{Factorization properads and Theorem A${}^{\infty,2}$ at properad level}
\label{sec:factorization-properad}
%%%%%%%%%%%%%%%%%%%%%%%%%%%%%%%%%%%%%%%%%%%%%%%%%%%%%%%%%%%%%%%%%%%%%%%%%%%%%

Chiral algebras on $X$ carry genuinely multi-output co-operations:
collision residues
$\mathrm{Res}^{\mathrm{coll}}\colon j_*j^*\cA^{\boxtimes n}\to \Delta_*\cA$
have a contravariant counterpart via Verdier duality,
$\Delta^!\cA\to j_*j^*\cA^{\boxtimes n}[\text{shift}]$,
which extracts a single input into $n$ outputs along the same
collision divisor. Operads model only inputs; the natural ambient
is \emph{properads} (multi-input and multi-output operations
without disconnected composition; Vallette~\cite{Val07},
Merkulov--Vallette~\cite{MV09}). The present section installs the
properad lift of Theorem~A in the Francis--Gaitsgory factorization
ambient and records the $(\infty,2)$-categorical equivalence.

\begin{definition}[Factorization properad]\label{def:factorization-properad}
\index{factorization properad|textbf}%
A \emph{factorization properad} $\cP$ on $X$ is the datum of:
\begin{enumerate}
\item[(P1)] for every pair $(m,n)\in\Z_{\geq 0}^2$ of input and
 output arities, a $D$-module
 \[
 \cP(m,n)\in D\text{-mod}\bigl(X^{m+n}\bigr)
 \]
 equivariant for the action of $\Sigma_m\times\Sigma_n$ by
 permutation of inputs and outputs;
\item[(P2)] a \emph{vertical} (sequential) composition
 \[
 \mu^{\mathrm{ver}}\colon
 \cP(n,p)\otimes^!_{X^n}\cP(m,n)\;\longrightarrow\;\cP(m,p)
 \]
 defined on the open stratum
 $\{x_i\neq x_j\}\subset X^{m+n+p}$ and extended by $j_*j^*\to\Delta_*$
 along collision divisors;
\item[(P3)] a \emph{horizontal} (tensor) composition
 \[
 \mu^{\mathrm{hor}}\colon
 \cP(m,n)\boxtimes \cP(m',n')\;\longrightarrow\;\cP(m+m',n+n')
 \]
 as $D$-modules on $X^{m+n}\times X^{m'+n'}$;
\item[(P4)] a unit $\eta\colon \mathbf{1}_X\to\cP(1,1)$;
\end{enumerate}
such that $\mu^{\mathrm{ver}}$ is associative, $\mu^{\mathrm{hor}}$
is associative and symmetric-monoidal, and the
\emph{interchange law}
$
 \mu^{\mathrm{ver}}\circ(\mu^{\mathrm{hor}}\boxtimes
 \mu^{\mathrm{hor}})
 \;=\;
 \mu^{\mathrm{hor}}\circ(\mu^{\mathrm{ver}}\boxtimes \mu^{\mathrm{ver}})
$
holds on the open stratum; wheeled composition along $1\to 1$ loops
(genus-raising) is excluded (non-wheeled properad).
\end{definition}

The arity-$(m,1)$ sub-datum
$\cP_{\mathrm{op}}(m):=\cP(m,1)$ is a factorization operad in the
sense of Beilinson--Drinfeld~\cite{BD04} and Francis--Gaitsgory~\cite{FG12};
the arity-$(1,n)$ sub-datum $\cP_{\mathrm{co}}(n):=\cP(1,n)$
is a factorization cooperad. A factorization properad is the
simultaneous carrier of both, glued by the interchange law.

\begin{example}[The endomorphism properad of a chiral algebra]
\label{ex:endomorphism-properad}
For a chiral algebra $\cA$ on $X$ whose underlying $D$-module is
Verdier-self-dual (as holds for all standard-landscape algebras;
cf.\ Chapter~\ref{ch:factorization-swiss-cheese} \S\ref{sec:factorization-substrate}),
define
\[
 \End^{\mathrm{pr}}_\cA(m,n)
 \;:=\;
 \Hom_{D\text{-mod}(X^{m+n})}\!\bigl(j_*j^*\cA^{\boxtimes m},
 \Delta_*\cA^{\boxtimes n}\bigr),
\]
with vertical composition given by composition along the middle
$X^n$ factor, horizontal composition given by external tensor,
and unit given by the identity on $\cA$. Collision residues
($\mathrm{Res}^{\mathrm{coll}}$) populate $\End^{\mathrm{pr}}_\cA(m,1)$;
their Verdier duals populate $\End^{\mathrm{pr}}_\cA(1,n)$; brace
compositions populate $\End^{\mathrm{pr}}_\cA(m,n)$ for $m,n\geq 1$.
A \emph{properad action} of $\cP$ on $\cA$ is a morphism
$\cP\to\End^{\mathrm{pr}}_\cA$ in factorization properads.
\end{example}

\begin{remark}[Why non-wheeled]\label{rem:non-wheeled}
Wheeled composition (tracing an output back to an input) would
produce genus-raising operations; these are proper to the modular
operad (Section~\ref{sec:modular-steinberg-hecke}) and are
incompatible with conilpotent completion at fixed genus. The
factorization properad records only genus-$0$ multi-output
co-operations; genus is added orthogonally by the modular
Steinberg correspondence.
\end{remark}

Properadic bar--cobar is due to Vallette~\cite{Val07} for properads
in $\mathrm{Vect}$ and to Merkulov--Vallette~\cite{MV09} in the
derived setting; the simplicial/$\infty$-categorical enhancement is
due to Hackney--Robertson--Yau~\cite{HRY20}. We import their
construction into the Francis--Gaitsgory factorization
ambient.

\begin{definition}[Properadic bar and cobar]\label{def:properadic-bar-cobar}
Let $\cP$ be a conilpotent coaugmented factorization properad and
$\cQ$ an augmented factorization properad. The
\emph{properadic bar} $\bar B^{\mathrm{pr}}(\cQ)$ and
\emph{properadic cobar} $\Omega^{\mathrm{pr}}(\cP)$ are the
cofree-conilpotent and free constructions in the Hackney--Robertson
graphical category~\cite{HRY20}, adapted fibrewise over
$\Ran(X)$. The differential on $\bar B^{\mathrm{pr}}(\cQ)$ is the
sum of the internal differential on $\cQ$ and a graphical
differential contracting a single internal edge by vertical
composition $\mu^{\mathrm{ver}}$ or horizontal composition
$\mu^{\mathrm{hor}}$; the differential on $\Omega^{\mathrm{pr}}(\cP)$
is dual.
\end{definition}

\begin{theorem}[Properad-level Theorem~A in the factorization ambient;
 \ClaimStatusProvedHere]\label{thm:properad-bar-cobar}
In the Francis--Gaitsgory factorization ambient $\Fact(X)$
(Reference: \cite{FG12}; see also Section~\ref{sec:factorization-substrate}),
the pair
\[
 \bar B^{\mathrm{pr}}\colon
 \operatorname{Prop}^{\mathrm{aug}}(\Fact(X))
 \;\rightleftarrows\;
 \operatorname{CoProp}^{\mathrm{conil}}(\Fact(X))
 \;\colon\Omega^{\mathrm{pr}}
\]
of properadic bar and cobar functors is an
$(\infty,2)$-categorical adjoint equivalence on the
conilpotent-complete locus, whose restriction to the arity-$(m,1)$
sub-properad recovers the operadic Theorem~A
(Theorem~\ref{thm:factorization-SC-koszul}). The equivalence
intertwines horizontal composition $\mu^{\mathrm{hor}}$ with
the external-tensor symmetric-monoidal structure on both sides,
and vertical composition $\mu^{\mathrm{ver}}$ with the
graphical-contraction differential on bars.
\end{theorem}
\begin{proof}
\textit{Construction.} Hackney--Robertson~\cite{HRY20} construct a
Quillen model structure on the category of simplicial properads
(equivalently, of properads in any combinatorial symmetric monoidal
model category). Their Theorem~4.3.1 establishes the bar--cobar
adjunction at the level of model structures; their Theorem~5.2.2
upgrades it to a bi-adjoint bi-equivalence of
$(\infty,2)$-categories once both sides are restricted to the
conilpotent-complete locus. The factorization ambient
$\Fact(X)$ is presentable (Francis--Gaitsgory~\cite{FG12}
Ch.~IV.5; cf.\ the coderived-models presentability lemma),
hence combinatorial, and
symmetric monoidal under the external $\boxtimes$ combined with
pushforward along the Ran diagonal. Transporting the
Hackney--Robertson machinery across this symmetric monoidal
combinatorial model yields the claim.

\textit{Restriction to operads.} The sub-properad on which
output arity is fixed to $n=1$ is closed under $\mu^{\mathrm{ver}}$
and inherits an operadic model structure. The properadic bar
differential restricted to $n=1$ operations coincides with the
operadic bar differential because graphical contraction of an
edge between two $(m,1)$-operations is the unique composition
available. Hence the restriction of Theorem~\ref{thm:properad-bar-cobar}
recovers Theorem~\ref{thm:factorization-SC-koszul} after the
Koszul-pair encoding.

\textit{$(\infty,2)$-structure.} The $(\infty,2)$-categorical
enhancement uses, on objects, the $\infty$-groupoid of
properads up to weak equivalence; on $1$-morphisms, the
$\infty$-groupoid of properad morphisms; on $2$-morphisms,
the mapping spaces of natural transformations, which exist
because the model category is simplicially enriched
(\cite{HRY20} Ch.~6). Francis--Gaitsgory's $(\infty,2)$-enhancement
of $\Fact(X)$ provides the same enrichment on the ambient.

\textit{Interchange compatibility.} The interchange law
(P4 of Definition~\ref{def:factorization-properad}) is preserved
by both functors because graphical contraction respects the
horizontal composition on disconnected graphs, and the
symmetric-monoidal structure is strict on the objects of the
image $(\infty,2)$-category.
\end{proof}

\begin{corollary}[Properad-level Koszul duality for chiral algebras]
\label{cor:chiral-properad-koszul}
For $\cA$ a chiral algebra on $X$ of finite chiral Hochschild
dimension, $\End^{\mathrm{pr}}_\cA$ admits a Koszul-dual
coproperad
$(\End^{\mathrm{pr}}_\cA)^{\mathrm{!`}}
 \simeq \bar B^{\mathrm{pr}}(\End^{\mathrm{pr}}_\cA)$,
unique up to quasi-isomorphism of factorization properads, whose
arity-$(m,1)$ restriction reproduces the Koszul dual $\cA^!$ of
Theorem~A and whose arity-$(1,n)$ restriction captures the
collision co-operations dual to chiral $n$-ary products.
\end{corollary}
\begin{proof}
Apply Theorem~\ref{thm:properad-bar-cobar} to
$\End^{\mathrm{pr}}_\cA$; uniqueness up to quasi-isomorphism
follows from the $(\infty,2)$-universal property of the bar
functor. Arity-$(m,1)$ restriction recovers the operadic Koszul
dual by the restriction clause of
Theorem~\ref{thm:properad-bar-cobar}; arity-$(1,n)$ restriction
recovers the collision cooperad by Verdier duality inside
$D$-mod$(X^{n+1})$.
\end{proof}

\begin{remark}[Corners with corners; global sections to $\Ran(X)$]
\label{rem:properad-payoff}
The multi-output structure captured by
Theorem~\ref{thm:properad-bar-cobar} provides two features that
are not available at operadic level.
\emph{(i) Corners with corners.} The modular bar complex
(Chapter~\ref{ch:modular-sc-operad}) admits corners
where both an input and an output meet; at operadic level
this requires a separate open-cover argument, while at properad
level the corners-with-corners are encoded directly by
$\cP(m,n)$ with $m,n\geq 1$. In particular, the
Steinberg--Hecke correspondence of
Section~\ref{sec:steinberg-hecke} lifts to an
internal-$\operatorname{Hom}$ in
$\operatorname{CoProp}^{\mathrm{conil}}(\Fact(X))$.
\emph{(ii) Global sections to $\Ran(X)$.} The global-sections
functor
$\Gamma\colon\Fact(X)\to D\text{-mod}(\Ran(X))$
is automatically compatible with the properad structure because
horizontal composition descends to the Ran diagonal. Hence
Theorem~\ref{thm:properad-bar-cobar} transports directly to a
Koszul-duality statement on $\Ran(X)$, which was the target
formulation of Theorem~A${}^{\infty,2}$ in the Platonic
Reconstitution.
\end{remark}

\begin{remark}[Scope of the properad lift]
\label{rem:properad-scope}
Theorem~\ref{thm:properad-bar-cobar} is stated for non-wheeled
factorization properads; wheeled extensions (which would
introduce genus-raising operations and correspond to modular
properads) are outside the scope and are addressed in
Chapter~\ref{ch:modular-sc-operad}. The restriction
to the conilpotent-complete locus is necessary: without it,
the properadic bar functor is not left-exact (Merkulov--Vallette~\cite{MV09}
Ch.~3), and the adjunction fails to be an equivalence. Both
conditions are satisfied by the factorization endomorphism
properads of chiral algebras in the standard landscape.
\end{remark}

\begin{proposition}[$R$-twisted $\Sigma_n$-descent between ordered
 and symmetric bar; \ClaimStatusProvedHere]
\label{prop:R-twisted-sigma-n-descent}
Let $\cA$ be a conilpotent-complete $E_1$-chiral algebra with
spectral $R$-matrix $R(z)$ satisfying YBE and braid coherence.
Write $X^n_{\mathrm{ord}} = X^n \setminus \Delta$ for the ordered
configuration scheme and $X^n_{\mathrm{sym}} =
X^n_{\mathrm{ord}}/\Sigma_n$ for its unordered quotient, and let
$\pi_n \colon X^n_{\mathrm{ord}} \to X^n_{\mathrm{sym}}$ be the
$\Sigma_n$-principal bundle. Then
\[
 B^{\Sigma}_n(\cA)
 \;\simeq\;
 \bigl(\pi_{n,*}\,B^{\mathrm{ord}}_n(\cA)\bigr)^{\Sigma_n,\,\cL_R}
\]
as factorization coalgebras on $X^n_{\mathrm{sym}}$, where the
$\Sigma_n$-coinvariance is taken in the twisted sense of the
$\Sigma_n$-equivariant local system $\cL_R$ on
$X^n_{\mathrm{ord}}$ whose monodromy along the elementary
transposition $(i\,i{+}1)$ swapping adjacent spectral parameters
$z_i, z_{i+1}$ is $R(z_i - z_{i+1})$. The local system $\cL_R$ is
well-defined: the braid relations
$(i\,i{+}1)(i{+}1\,i{+}2)(i\,i{+}1)
 = (i{+}1\,i{+}2)(i\,i{+}1)(i{+}1\,i{+}2)$
lift to $\Sigma_n$-relations by YBE, and the involutive relation
$(i\,i{+}1)^2 = 1$ lifts by the inversion axiom
$R(z)\,R_{\mathrm{op}}(-z) = \mathrm{id}$. At the pole-free point
$R(z) = \tau$ (the swap), $\cL_R$ is the trivial local system and
the formula reduces to the classical Loday--Vallette
$\Sigma$-coinvariance
$B^{\Sigma}_n(\cA) = (B^{\mathrm{ord}}_n(\cA))_{\Sigma_n}$
of~\cite{LV12} specialised to the pseudo-tensor
$(\mathsf{Ass}, \mathsf{Ass}^!)$ Koszul pair.
\end{proposition}
\begin{proof}
\textit{Step 1 (construction of $\cL_R$).} For each pair of
adjacent indices $(i, i{+}1)$, the elementary transposition
defines a deck transformation of the $\Sigma_n$-cover
$\pi_n$; the $R$-matrix $R(z_i - z_{i+1})$ provides the
monodromy representation of the local system on the stalk
$\cA^{\otimes n}$ at a generic point.

\textit{Step 2 (braid coherence via YBE).} The symmetric-group
braid relation on three consecutive indices imposes the
identity
$(R_{12} R_{13} R_{23})(z_i, z_j, z_k)
 = (R_{23} R_{13} R_{12})(z_i, z_j, z_k)$
on the stalk $\cA^{\otimes 3}$; this is exactly the quantum
Yang--Baxter equation (Theorem~\ref{thm:YBE}), so $\cL_R$ is
well-defined on codimension-$2$ strata.

\textit{Step 3 (involutivity).} $(i\,i{+}1)^2 = 1$ in $\Sigma_n$
forces $R(z_i - z_{i+1})\,R(z_{i+1} - z_i)\,\tau^2 = \mathrm{id}$,
equivalently $R(z)\,R_{\mathrm{op}}(-z) = \mathrm{id}$; this is the
standard inversion axiom, and holds on the pseudo-tensor
Koszul locus by the quasi-triangularity axiom of
Theorem~\ref{thm:Koszul_dual_Yangian}.

\textit{Step 4 (descent to $B^\Sigma$).} By faithfully flat descent
for $\Sigma_n$-principal bundles in the factorization ambient
(Gaitsgory--Rozenblyum~\cite[IV.5.4]{GR17}), the category of
$\Sigma_n$-equivariant factorization coalgebras on
$X^n_{\mathrm{ord}}$ is equivalent to the category of
factorization coalgebras on $X^n_{\mathrm{sym}}$. Under this
equivalence, the $\cL_R$-twisted $\Sigma_n$-coinvariants of
$B^{\mathrm{ord}}_n(\cA)$ correspond to the symmetric bar
$B^{\Sigma}_n(\cA)$ constructed by taking $\Sigma_n$-coinvariants
of the $R$-matrix symmetrised tensor powers of $\bar\cA$.

\textit{Step 5 (pole-free reduction).} At $R(z) = \tau$,
$\cL_R = \underline{k}$ is the constant (trivial) local system.
The twisted coinvariance reduces to ordinary
$\Sigma_n$-coinvariance
$(B^{\mathrm{ord}}_n(\cA))_{\Sigma_n} = B^{\mathrm{naive}}_n(\cA)$,
which is the Loday--Vallette bar of the classical Koszul pair
$(\mathsf{Ass}, \mathsf{Ass}^!)$~\cite{LV12}.
\end{proof}

\begin{remark}[Ordered vs.\ symmetric bar primacy]\label{rem:ordered-symmetric-bar-primacy}
Proposition~\ref{prop:R-twisted-sigma-n-descent} resolves the
objection that the classical Koszul argument for Theorem~A
requires a symmetric monoidal ambient which $\Dmod(X)$ under
$\otimes^{\mathrm{ch}}$ does not provide: we do not attempt to
upgrade the BD pseudo-tensor structure. Instead, both the ordered
bar $B^{\mathrm{ord}}(\cA)$ and the symmetric bar
$B^{\Sigma}(\cA)$ are promoted to factorization (co)algebras on
$\Ran(X)^{\mathrm{ord}}$ and $\Ran(X)$ respectively, where the
Francis--Gaitsgory $\star$-tensor~\cite[IV.5.2]{GR17} does provide
an ambient symmetric monoidal structure (not on $\Dmod(X)$ but on
$\Dmod(\Ran(X))$). The descent
Proposition~\ref{prop:R-twisted-sigma-n-descent} is an identity
between two factorization coalgebras on $\Ran(X)$, not on a single
fibre. Hence the ``ordered $E_1$ primacy'' headline of Theorem~A
stands: the ordered bar is primary, the symmetric bar is its
$R$-twisted $\Sigma_n$-coinvariant image, and at pole-free point
the two coincide with the classical Loday--Vallette
construction.
\end{remark}


%%%%%%%%%%%%%%%%%%%%%%%%%%%%%%%%%%%%%%%%%%%%%%%%%%%%%%%%%%%%%%%%%%%%%%%%%%%%%
\section{The modular Steinberg correspondence}
\label{sec:steinberg-hecke}
%%%%%%%%%%%%%%%%%%%%%%%%%%%%%%%%%%%%%%%%%%%%%%%%%%%%%%%%%%%%%%%%%%%%%%%%%%%%%

The convolution-product interpretation of the modular Lie bracket
(Remark~\ref{rem:convolution-as-correspondence}) lifts to
the factorization setting. This section makes the Steinberg--Hecke
parallel precise: the clutching correspondence on $\Mbar_g$ is the
\emph{modular} analogue of the Steinberg variety
$\widetilde{\cN}\times_\fg \widetilde{\cN}$ in geometric
representation theory, and the convolution algebra it produces is
the modular Lie algebra $L_{\mathrm{mod}}$.

\begin{construction}[The modular clutching correspondence]
\label{constr:modular-clutching-correspondence}
\index{clutching correspondence!modular}
\index{Steinberg variety!modular analogue}
For a decomposition $g = g_1 + g_2$ with $n = n_1 + n_2$ marked
points, the \emph{separating clutching correspondence} is the span
\begin{equation}\label{eq:sep-clutching-span}
\begin{tikzcd}[column sep=3em]
 \Mbar_{g_1,n_1+1} \times \Mbar_{g_2,n_2+1}
 & Z_{\mathrm{sep}}
 \ar[l, "p"']
 \ar[r, "\xi_{\mathrm{sep}}"]
 & \Mbar_{g,n}
\end{tikzcd}
\end{equation}
where $Z_{\mathrm{sep}}$ is the boundary stratum of curves with a
separating node, $p$ remembers the two components, and
$\xi_{\mathrm{sep}}$ is the gluing map.

The \emph{modular Steinberg variety} is the fibre product
\begin{equation}\label{eq:modular-steinberg}
 \mathsf{St}_{g,n}^{\mathrm{mod}}
 \;:=\;
 Z_{\mathrm{sep}} \;\times_{\Mbar_{g,n}}\; Z_{\mathrm{sep}},
\end{equation}
parametrising pairs of separating degenerations of the same stable
curve. A point of $\mathsf{St}_{g,n}^{\mathrm{mod}}$ is a stable
curve $C \in \Mbar_{g,n}$ together with two (possibly identical)
separating nodes. The convolution product on equivariant
(Borel--Moore) homology of this correspondence is:
\begin{equation}\label{eq:steinberg-convolution}
 [\alpha] \star [\beta]
 \;=\;
 (\xi_{\mathrm{sep}})_*\,
 \bigl(p^*(\alpha \boxtimes \beta)\bigr)
 \;=\;
 [\alpha, \beta]_{L_{\mathrm{mod}}},
\end{equation}
the Lie bracket on $L_{\mathrm{mod}}$
(Remark~\ref{rem:convolution-as-correspondence}).
\end{construction}

\begin{proposition}[Steinberg--Hecke parallel]
\label{prop:steinberg-hecke-parallel}
\index{Steinberg--Hecke parallel|textbf}
\ClaimStatusProvedHere{}
The analogy between the classical and modular Steinberg varieties is
exact in the following sense.
\begin{center}
\renewcommand{\arraystretch}{1.3}
\begin{tabular}{lll}
\textbf{Classical} & \textbf{Modular} & \textbf{Output} \\
\hline
$\widetilde{\cN} = T^*(\cG/\cB)$ &
 $\Mbar_{g_1,n_1+1}$ &
 Resolution of singularities \\
$\fg$ &
 $\Mbar_{g,n}$ &
 Total space \\
$\widetilde{\cN} \times_\fg \widetilde{\cN}$ &
 $Z_{\mathrm{sep}} \times_{\Mbar_{g,n}} Z_{\mathrm{sep}}$ &
 Steinberg variety \\
$H_*^{\cG}(\mathrm{St})$ &
 $H_*(\mathsf{St}^{\mathrm{mod}})$ &
 Convolution algebra \\
Hecke algebra $\cH(W)$ &
 $L_{\mathrm{mod}}$ &
 Algebraic output \\
\end{tabular}
\end{center}
The bar complex organizes the clutching correspondences in the same
sense that the Steinberg variety organizes the Hecke algebra: both
arise from convolution structures on self-intersection varieties of
resolution correspondences.

The non-separating clutching extends the parallel: the
non-separating correspondence
\begin{equation}\label{eq:nonsep-clutching-span}
\begin{tikzcd}[column sep=3em]
 \Mbar_{g,n+2}
 & Z_{\mathrm{nsep}}
 \ar[l, "p'"']
 \ar[r, "\xi_{\mathrm{nsep}}"]
 & \Mbar_{g+1,n}
\end{tikzcd}
\end{equation}
produces the genus-raising differential $D_1$
\textup{(}Construction~\textup{\ref{constr:nonsep-clutching})}, and the
non-separating Steinberg variety
$Z_{\mathrm{nsep}} \times_{\Mbar_{g+1,n}} Z_{\mathrm{nsep}}$
governs the $D_1^2 = 0$ relation.
\end{proposition}

\begin{proof}
The Lie bracket identification \eqref{eq:steinberg-convolution}
is Remark~\ref{rem:convolution-as-correspondence}, restated in
the factorization setting. The Jacobi identity for
$L_{\mathrm{mod}}$ translates to the associativity of the
convolution product on $\mathsf{St}^{\mathrm{mod}}$, which
follows from the associativity of the composition of
correspondences:
\[
 (Z_{\mathrm{sep}} \circ Z_{\mathrm{sep}}) \circ Z_{\mathrm{sep}}
 \;=\;
 Z_{\mathrm{sep}} \circ (Z_{\mathrm{sep}} \circ Z_{\mathrm{sep}})
\]
as cycles in $\Mbar_{g,n}$. This is a standard fact about
correspondence products (Chriss--Ginzburg \cite{CG97},
Theorem~2.7.26).

For the non-separating case: $D_1^2 = 0$ is the statement that
the self-composition of the non-separating clutching correspondence
is zero in homology. Geometrically: two successive non-separating
degenerations produce a codimension-$2$ boundary stratum in
$\Mbar_{g+2,n}$, and the orientations of the two nodes conspire
to cancel (each node contributes a sign from the orientation line
$\orline_\Gamma = \det(\mathrm{Edges})^{\otimes(-1)}$, and two
nodes with opposite pairings cancel). This is the same mechanism as
$\partial^2 = 0$ for the cell complex of $\Mbar_{g,n}$, specialised
to the non-separating boundary.

The mixed anticommutativity $D_0 D_1 + D_1 D_0 = 0$ is the
compatibility of separating and non-separating clutching: a
separating degeneration followed by a non-separating one produces
the same codimension-$2$ boundary stratum as the reverse order,
with a relative sign from the Koszul rule
(Remark~\ref{rem:bicomplex-signs}).
\end{proof}

\begin{remark}[The factorization setting]
\label{rem:steinberg-factorization}
The Steinberg--Hecke parallel extends to the
factorization framework of this chapter. The factorizable
$D$-module $F_\cA$ on $\Ran(\Sigma_g)$ associates to each point
of $\Mbar_g$ a factorization algebra on the corresponding curve.
The clutching correspondence \eqref{eq:sep-clutching-span} acts on
$F_\cA$ by the factorization splitting: the restriction of $F_\cA$
to a separating node decomposes as $F_\cA|_{z_1} \otimes F_\cA|_{z_2}$
by the factorization isomorphism
(Construction~\ref{constr:fact-dmod-from-chiral}(iii)), and the
convolution product contracts this tensor product by the invariant
bilinear form on~$\cA$. This is the factorization-level
realisation of the Lie bracket on $L_{\mathrm{mod}}$.

The non-separating clutching acts on $F_\cA$ by the non-separating
clutching map $\Delta_{\mathrm{cyc}}$
(Construction~\ref{constr:nonsep-clutching}): it contracts a pair
of tensor factors by $\langle{-},{-}\rangle_{\mathrm{inv}}$ and
raises genus by~$1$. The bar complex $\barB(F_\cA)$, viewed as a
factorization coalgebra on $\Ran(\Sigma_g)$ parametrised by $\Mbar_g$,
is the \emph{universal} object on which both clutching correspondences
act: it is the simultaneous presentation of all convolution products
at all genera.

Thus: the bar complex is the universal coalgebraic stage on which the
Swiss-cheese clutching correspondences act, as the Steinberg variety
is the geometric stage underlying the Hecke algebra. Both arise from
convolution structures on self-intersection varieties of resolution
correspondences, differing only in the ambient geometry.
$\fg$ for the classical Steinberg variety, $\Mbar_g$ for the
modular one.
\end{remark}


%%%%%%%%%%%%%%%%%%%%%%%%%%%%%%%%%%%%%%%%%%%%%%%%%%%%%%%%%%%%%%%%%%%%%%%%%%%%%
\section{Latyntsev factorization quantum groups and colour projections}
\label{sec:latyntsev-colour-projections}
%%%%%%%%%%%%%%%%%%%%%%%%%%%%%%%%%%%%%%%%%%%%%%%%%%%%%%%%%%%%%%%%%%%%%%%%%%%%%

The factorization Swiss-cheese framework unifies two programmes
that have developed independently: the factorization quantum groups
of Latyntsev \cite{Latyntsev23} and the dg-shifted Yangians of
Chapter~\ref{chap:dg-shifted-factorization}. Both arise as colour
projections of a single object.

\begin{theorem}[Colour projections and factorization quantum groups; \ClaimStatusConditional]
\label{thm:latyntsev-colour-projections}
\index{Latyntsev factorization quantum groups|textbf}
\index{factorization quantum groups!colour projections|textbf}
Let $F$ be a factorization Swiss-cheese algebra on
$\Ran^{\mathrm{oc}}_{\HT}(\Sigma_g,I)$. The two colour projections are:
\begin{enumerate}[label=\textup{(\roman*)}]
\item \textbf{Closed-colour projection.} The forgetful functor
 $\pi_{\mathsf{cl}}\colon F \mapsto F_{\mathsf{cl}}$ produces a
 factorization algebra on $\Ran(\Sigma_g)$. For
 $F = F_{\cA}$ arising from an affine vertex algebra
 $\cA = V_k(\fg)$, this recovers Latyntsev's factorization
 quantum group $\Phi_q(\fg)$ with
 $q = \exp\!\bigl(2\pi i / (k + h^\vee)\bigr)$:
 \[
 \pi_{\mathsf{cl}}(F_{V_k(\fg)})
 \;\simeq\;
 \Phi_q(\fg).
 \]

\item \textbf{Open-colour projection.} The forgetful functor
 $\pi_{\mathsf{op}}\colon F \mapsto F_{\mathsf{op}}$ produces an
 $\Eone$-algebra. For $F = F_{V_k(\fg)}$ arising from the
 standard affine HT gauge realization, this recovers the
 dg-shifted Yangian $Y_\hbar(\fg)$ of
 Theorem~\textup{\ref{thm:Koszul_dual_Yangian}}:
 \[
 \pi_{\mathsf{op}}(F_{V_k(\fg)})
 \;\simeq\;
 Y_\hbar(\fg).
 \]

\item \textbf{Mixed projection.} The bulk--boundary composition map
 $\mu_{\mathsf{mix}}\colon F_{\mathsf{cl}} \otimes F_{\mathsf{op}}
 \to F_{\mathsf{op}}$ encodes the $R$-matrix:
 \[
 R(z)
 \;=\;
 \int_0^\infty e^{-\lambda z}\,
 \{{\cdot}_\lambda{\cdot}\}\, d\lambda,
 \]
 which intertwines the Latyntsev quantum group action
 \textup{(}via factorization on $\Ran(\Sigma_g)$\textup{)} with
 the Yangian action \textup{(}via the $\Eone$-structure\textup{)}.
 On the evaluation locus, $R(z)$ agrees with the spectral
 $R$-matrix of
 Theorem~\textup{\ref{thm:affine-monodromy-identification}}.
\end{enumerate}
\end{theorem}

\begin{center}
\begin{tabular}{lll}
\textbf{Colour} & \textbf{Factorization data} & \textbf{Algebraic output} \\
\hline
Closed & $F_{\mathsf{cl}}$ on $\Ran(\Sigma_g)$ &
 Quantum group $\Phi_q(\fg)$ (Latyntsev) \\
Open & $F_{\mathsf{op}}$ on $\Ran(\R)$ &
 dg-shifted Yangian $Y_\hbar(\fg)$ \\
Mixed & $F_{\mathsf{mix}}$ on
 $\Ran^{\mathrm{oc}}_{\HT}(\Sigma_g,I)$ &
 Spectral $R$-matrix $R(z)$
\end{tabular}
\end{center}

\begin{proof}
\emph{Part~(i): Closed-colour projection recovers $\Phi_q(\fg)$.}
The closed-colour projection $\pi_{\mathsf{cl}}$ forgets the
open colour and mixed operations. The remaining data is a
factorization algebra on $\Ran(\Sigma_g)$, which for
$\cA = V_k(\fg)$ is the vacuum factorization algebra:
the factorizable $D$-module $F_{\hat\fg_k}$ of
\eqref{eq:affine-fact-dmod}, restricted to the closed sector.
Latyntsev \cite{Latyntsev23}, Theorem~4.2.1, identifies the
factorization homology of the vacuum factorization algebra with the
quantum group representation category:
$\int_{\mathbb{P}^1} \Phi_q(\fg) \simeq \operatorname{Rep}_q(\fg)$.
The identification $\pi_{\mathsf{cl}}(F_{V_k(\fg)}) \simeq
\Phi_q(\fg)$ follows from the universality of the vacuum
factorization algebra: it is the \emph{unique} factorization
algebra on $\Ran(X)$ whose fibre at each point is the vacuum
module $V_k(\fg)$, with factorization isomorphisms given by
the sewing/OPE. Latyntsev's $\Phi_q(\fg)$ is characterised by
the same universal property (a factorization algebra on
$\Ran(\mathbb{P}^1)$ determined by the vacuum module and OPE)
so the two objects coincide. The quantum parameter is
$q = \exp\!\bigl(2\pi i/(k+h^\vee)\bigr)$, arising from the
affine Drinfeld--Kohno comparison package at level~$k$.

\smallskip\noindent
\emph{Part~(ii): Open-colour projection recovers $Y_\hbar(\fg)$.}
The open-colour projection $\pi_{\mathsf{op}}$ forgets the closed
colour. The remaining data is a locally constant factorization
algebra on $\Ran(\R)$, equivalently an $\Eone$-algebra
(Lurie \cite{HA}, Theorem~5.4.5.9). For $\cA = V_k(\fg)$,
this $\Eone$-algebra is the bar-cobar dual of $V_k(\fg)$ in
the topological direction: the associative algebra structure
on the boundary. By Theorem~\ref{thm:Koszul_dual_Yangian}
\textup{(}proved in Chapter~\ref{chap:spectral-braiding} for the
standard affine HT gauge realization\textup{)}, the open-colour
Koszul dual $\cA^!_{\mathrm{line}}$ of the affine vertex algebra
$V_k(\fg)$ in the $\Eone$-direction is the dg-shifted Yangian
$Y_\hbar(\fg)$.
The open-colour projection is therefore
$\pi_{\mathsf{op}}(F_{V_k(\fg)}) \simeq Y_\hbar(\fg)$.

\smallskip\noindent
\emph{Part~(iii): Mixed projection and the $R$-matrix.}
The mixed projection extracts the bulk--boundary composition map
$\mu_{\mathsf{mix}}$. This map is computed by the propagator
integral on the half-space $\Sigma_g \times \R_{\geq 0}$:
the propagator factors as
$P(z,t;\,w,t') = K^{(g)}(z,w) \cdot \theta(t-t')$
(Cauchy kernel times Heaviside,
Definition~\ref{def:BV-origin}, \eqref{eq:HT-propagator}).
The Laplace transform of the PVA $\lambda$-bracket
$\{{\cdot}_\lambda{\cdot}\}$ produces the spectral $R$-matrix
\[
 R(z)
 \;=\;
 \int_0^\infty e^{-\lambda z}\,
 \{{\cdot}_\lambda{\cdot}\}\, d\lambda.
\]
This is the DK-0 identification (the ninth cross-volume bridge),
proved as the Laplace-transform formula relating the PVA
$\lambda$-bracket to the spectral $R$-matrix. The $R$-matrix
$R(z)$ intertwines the closed-colour action
(Latyntsev quantum group, via factorization on $\Ran(\Sigma_g)$)
with the open-colour action (Yangian, via the
$\Eone$-structure): this intertwining property is the
defining property of the mixed sector
(Definition~\ref{def:BD-SC}(iii)(a), the bulk-action-on-boundary
axiom).

The agreement of $R(z)$ with the spectral $R$-matrix on the
evaluation locus is
Theorem~\ref{thm:affine-monodromy-identification}:
the KZ braid monodromy on the comparison surface
(computed by Latyntsev's factorization homology) is compared by
the Drinfeld--Kohno theorem with the corresponding quantum-group
braiding, and our $R(z)$ recovers both
upon evaluation. The strictification theorem
(Theorem~\ref{thm:complete-strictification}) identifies the passage
from quasi-factorization structure to strict factorization quantum
group with the vanishing of the root-sector scalar classes
$[\epsilon_{\alpha,n}]\in
\mathscr K_n/\operatorname{im}(\delta_{\alpha,n})$. The two
projections are rigidly linked on that scalar-vanishing locus, with
the path sectors already killed by the $1/n$ BCH calculation.
\end{proof}

\begin{remark}[What is new]
\label{rem:latyntsev-what-is-new}
The two-coloured factorization framework places both the quantum
group (closed colour) and the Yangian (open colour) as projections
of a \emph{single object}: the factorization Swiss-cheese algebra.
The $R$-matrix is neither a closed-colour nor an open-colour datum,
but a \emph{mixed-colour} datum.
\end{remark}


%%%%%%%%%%%%%%%%%%%%%%%%%%%%%%%%%%%%%%%%%%%%%%%%%%%%%%%%%%%%%%%%%%%%%%%%%%%%%
\section{Genus-$1$ curvature for the Heisenberg algebra: a complete derivation}
\label{sec:heisenberg-genus-1-complete}
%%%%%%%%%%%%%%%%%%%%%%%%%%%%%%%%%%%%%%%%%%%%%%%%%%%%%%%%%%%%%%%%%%%%%%%%%%%%%

\providecommand{\dfib}{d_{\mathrm{fib}}}
\providecommand{\Dg}[1]{D_{#1}}
\providecommand{\Res}{\operatorname{Res}}

The preceding examples stated the curvature identity
$\dfib^{\,2} = \kappa(\cA) \cdot \omega_g$ and verified it
by citing the cogenerator projection argument. For the simplest
chiral algebra (the Heisenberg (free boson) at genus~$1$)
the computation below carries out the derivation from first principles, writing
every propagator, every residue, and every sign explicitly. The
result is a self-contained derivation of the central identity of
the genus tower, which the reader can verify line by line.

%% =====================================
\subsection{The genus-$1$ Arakelov propagator}
\label{subsec:arakelov-genus-1-explicit}
%% =====================================

Let $\Sigma_1 = \C / (\Z + \Z\tau)$ be the elliptic curve with
period $\tau \in \mathfrak{H}_1$ (the upper half-plane,
$\operatorname{Im}\tau > 0$). Let $\vartheta_1(u|\tau)$ denote
the Jacobi theta function with odd half-period characteristic:
\begin{equation}\label{eq:jacobi-theta-1-def}
 \vartheta_1(u|\tau)
 \;=\;
 2\sum_{n=0}^\infty (-1)^n\, q^{(n+1/2)^2/2}\,
 \sin\bigl((2n+1)\pi u\bigr),
 \qquad
 q = e^{2\pi i \tau}.
\end{equation}
This function vanishes to first order at $u = 0$ and satisfies
$\vartheta_1(-u|\tau) = -\vartheta_1(u|\tau)$.

\begin{proposition}[Arakelov propagator at genus~$1$]
\label{prop:arakelov-genus-1-derivation}
\index{Arakelov propagator!genus 1!explicit derivation}
\ClaimStatusProvedHere{}
The Arakelov-normalised single-valued propagator on $\Sigma_1$ is
\begin{equation}\label{eq:eta-1-explicit}
 \eta^{(1)}(z,w;\tau)
 \;=\;
 \partial_z \log\vartheta_1(z-w|\tau)
 \;+\;
 \frac{2\pi i\, \operatorname{Im}(z-w)}{\operatorname{Im}\tau},
\end{equation}
viewed as a meromorphic one-form in~$z$
\textup{(}after tensoring with~$dz$\textup{)}
with a simple pole at $z = w$ of residue~$+1$, single-valued on
$\Sigma_1$. Its $\bar\partial$ in~$z$ is:
\begin{equation}\label{eq:dbar-eta-1}
 \bar\partial_z\, \eta^{(1)}(z,w;\tau)
 \;=\;
 -2\pi i \cdot \omega_{\mathrm{Ar}},
 \qquad
 \omega_{\mathrm{Ar}}
 \;=\;
 \frac{1}{2i\,\operatorname{Im}\tau}\,
 dz \wedge d\bar z,
\end{equation}
where $\omega_{\mathrm{Ar}}$ is the Arakelov $(1,1)$-form on~$\Sigma_1$,
normalised so that $\int_{\Sigma_1}\omega_{\mathrm{Ar}} = 1$.
\end{proposition}

\begin{proof}
\emph{Step~1: Monodromy of the holomorphic part.}
Set $u = z - w$. The holomorphic propagator
$\eta^{(1)}_{\mathrm{hol}}(u) = \partial_u \log\vartheta_1(u|\tau)$
has quasi-periodicities inherited from the transformation law of the
Jacobi theta function:
\begin{alignat}{2}
 \vartheta_1(u+1|\tau) &= -\vartheta_1(u|\tau)
 &\qquad&\Longrightarrow\quad
 \eta^{(1)}_{\mathrm{hol}}(u+1) = \eta^{(1)}_{\mathrm{hol}}(u),
 \label{eq:theta-A-cycle}\\[3pt]
 \vartheta_1(u+\tau|\tau) &= -q^{-1/2}\,e^{-2\pi i u}\,\vartheta_1(u|\tau)
 &\qquad&\Longrightarrow\quad
 \eta^{(1)}_{\mathrm{hol}}(u+\tau) = \eta^{(1)}_{\mathrm{hol}}(u) - 2\pi i.
 \label{eq:theta-B-cycle}
\end{alignat}
The $A$-cycle monodromy is trivial. The $B$-cycle monodromy is
$-2\pi i$: the propagator shifts by a constant. At genus~$1$ the
normalised holomorphic differential is $\omega_1 = dz$
(with $\oint_A dz = 1$, $\oint_B dz = \tau$), and the monodromy
shift is $-2\pi i \cdot \omega_1(w)/dw = -2\pi i$, recovering
the general formula \eqref{eq:propagator-monodromy} with $g = 1$.

\smallskip\noindent
\emph{Step~2: Non-holomorphic correction.}
The Arakelov correction term
$C(u) := 2\pi i\, \operatorname{Im}(u)/\operatorname{Im}\tau$
has monodromy:
\begin{align}
 C(u+1) &= C(u) + 2\pi i \cdot
 \frac{\operatorname{Im}(1)}{\operatorname{Im}\tau}
 = C(u), \label{eq:correction-A}\\[3pt]
 C(u+\tau) &= C(u) + 2\pi i \cdot
 \frac{\operatorname{Im}\tau}{\operatorname{Im}\tau}
 = C(u) + 2\pi i. \label{eq:correction-B}
\end{align}
The monodromy of $C$ around the $B$-cycle is $+2\pi i$, exactly
cancelling the $-2\pi i$ monodromy of $\eta^{(1)}_{\mathrm{hol}}$.
Therefore $\eta^{(1)} = \eta^{(1)}_{\mathrm{hol}} + C$ is
doubly periodic, hence single-valued on~$\Sigma_1$.

\smallskip\noindent
\emph{Step~3: The $\bar\partial$ computation.}
The holomorphic part satisfies
$\bar\partial_z \eta^{(1)}_{\mathrm{hol}} = 0$ (it is meromorphic
in~$z$). The correction term gives:
\[
 \bar\partial_z\, C(z-w)
 \;=\;
 \bar\partial_z \Bigl[
 \frac{2\pi i}{2i\,\operatorname{Im}\tau}
 \bigl(z - w - (\bar z - \bar w)\bigr)
 \Bigr]
 \;=\;
 \frac{2\pi i}{2i\,\operatorname{Im}\tau}\cdot(-d\bar z)
 \;=\;
 -\frac{\pi}{\operatorname{Im}\tau}\, d\bar z.
\]
Tensoring with the implicit $dz$ from the one-form structure:
\[
 \bar\partial_z\, \eta^{(1)}(z,w;\tau)\, dz
 \;=\;
 -\frac{\pi}{\operatorname{Im}\tau}\, dz \wedge d\bar z
 \;=\;
 -2\pi i \cdot
 \underbrace{\frac{1}{2i\,\operatorname{Im}\tau}\,
 dz \wedge d\bar z}_{\omega_{\mathrm{Ar}}},
\]
establishing \eqref{eq:dbar-eta-1}.
\end{proof}

%% =====================================
\subsection{The fibre differential on degree-$2$ bar elements}
\label{subsec:dfib-degree-2}
%% =====================================

\begin{computation}[Fibre differential $\dfib$ at genus~$1$ for the
Heisenberg algebra]
\label{comp:dfib-heisenberg-genus1}
\index{Heisenberg algebra!fibre differential at genus 1}
\ClaimStatusProvedHere{}
Let $\cA = \cH_\kappa$ be the Heisenberg chiral algebra with OPE
\begin{equation}\label{eq:heisenberg-OPE-kappa}
 a(z)\, b(w)
 \;\sim\;
 \frac{\kappa}{z-w},
\end{equation}
where $\kappa \in k$ is the level. Let $a, b \in \bar\cA$
be generators of the augmentation ideal. Consider the degree-$2$
bar element $\xi = [s^{-1}a \,|\, s^{-1}b] \in \barB^2(\cH_\kappa)$
at genus~$1$.

The fibre differential $\dfib$ acts by the binary chiral product
with the genus-$1$ Arakelov propagator. The formula is:
\begin{equation}\label{eq:dfib-bar-2}
 \dfib\bigl([s^{-1}a \,|\, s^{-1}b]\bigr)
 \;=\;
 s^{-1}\,\Res_{z_1 = z_2}\bigl[
 \eta^{(1)}(z_1, z_2;\tau) \cdot
 a(z_1)\, b(z_2)
 \bigr],
\end{equation}
where $\Res_{z_1 = z_2}$ extracts the coefficient of
$(z_1 - z_2)^{-1}$ in the Laurent expansion.

\smallskip\noindent
\textbf{Explicit residue.}
Write $u = z_1 - z_2$. Near $u = 0$, the Arakelov propagator
has the Laurent expansion
\begin{equation}\label{eq:arakelov-laurent}
 \eta^{(1)}(z_1,z_2;\tau)
 \;=\;
 \frac{1}{u}
 \;+\;
 c_1(\tau)\, u
 \;+\;
 c_3(\tau)\, u^3
 \;+\;
 \cdots
 \;+\;
 \frac{2\pi i\,\operatorname{Im}(u)}{\operatorname{Im}\tau},
\end{equation}
where $c_1(\tau) = -\frac{1}{3}\frac{\vartheta_1'''(0|\tau)}
{\vartheta_1'(0|\tau)}$ encodes the Eisenstein series $E_2(\tau)$
via $c_1(\tau) = -\frac{\pi^2}{3}E_2(\tau)$, and the
non-holomorphic term $\operatorname{Im}(u)/\operatorname{Im}\tau$
is smooth at $u = 0$ and contributes nothing to the residue.

The OPE $a(z_1)\,b(z_2) \sim \kappa/u$ has a simple pole. The
product is:
\begin{align}\label{eq:propagator-times-OPE}
 \eta^{(1)}(z_1,z_2;\tau) \cdot a(z_1)\,b(z_2)
 &\;=\;
 \Bigl(\frac{1}{u} + c_1\,u + \cdots
 + \frac{2\pi i\,\operatorname{Im}(u)}{\operatorname{Im}\tau}
 \Bigr)
 \cdot
 \Bigl(\frac{\kappa}{u} + \text{regular}\Bigr)
 \notag\\[3pt]
 &\;=\;
 \frac{\kappa}{u^2}
 \;+\;
 \frac{(\text{regular})}{u}
 \;+\;
 \cdots.
\end{align}
The residue (coefficient of $u^{-1}$) involves the regular part
of the OPE (the normal-ordered product $:ab:$) contracted with
the leading $1/u$ of the propagator, plus the leading OPE singular
part $\kappa/u$ contracted with the $O(1)$ part of the propagator.
Since the Arakelov propagator is odd in~$u$ (to leading order),
the $O(1)$ part of the Laurent expansion of $\eta^{(1)}$ vanishes.
Therefore:
\begin{equation}\label{eq:dfib-result}
 \dfib\bigl([s^{-1}a \,|\, s^{-1}b]\bigr)
 \;=\;
 s^{-1}\bigl({:}a \cdot b{:}\bigr),
\end{equation}
where ${:}a \cdot b{:}$ is the normal-ordered product (the regular
part of the OPE). For the Heisenberg algebra,
${:}a \cdot b{:} = a \cdot b$ is the commutative product in
$\Sym(\omega_X)$.
\end{computation}

%% =====================================
\subsection{The curvature $\dfib^{\,2}$ on degree-$1$ bar elements}
\label{subsec:dfib-sq-degree-1}
%% =====================================

\begin{computation}[Curvature $\dfib^{\,2}$ at genus~$1$ for the
Heisenberg algebra]
\label{comp:dfib-sq-heisenberg-genus1}
\index{Heisenberg algebra!curvature at genus 1}
\ClaimStatusProvedHere{}
Let $[s^{-1}a] \in \barB^1(\cH_\kappa)$ be a degree-$1$ bar
element at genus~$1$. The curvature $\dfib^{\,2}$ acts as:
\begin{equation}\label{eq:dfib-sq-degree-1}
 \dfib^{\,2}\bigl([s^{-1}a]\bigr)
 \;=\;
 \dfib\bigl(\dfib([s^{-1}a])\bigr).
\end{equation}
Since $\dfib$ is a coderivation of $T^c(s^{-1}\bar\cA)$, and the
cogenerator projection of $\dfib^{\,2}$ is the obstruction to
$d^2 = 0$, we compute by the following two-step procedure.

\smallskip\noindent
\textbf{Step~1: Lift $[s^{-1}a]$ to a degree-$2$ element via $\dfib$.}
The differential $\dfib$ on a degree-$1$ element is determined by
the $A_\infty$ operation $m_0$ (the curving): at genus~$0$,
$m_0 = 0$ and $\dfib([s^{-1}a]) = 0$. At genus~$1$, the
curving is non-zero: $\dfib([s^{-1}a])$ is computed by inserting
a vacuum bubble (a loop with no external legs) into the bar element.
This vacuum bubble involves the self-contraction of the propagator:
\begin{equation}\label{eq:vacuum-bubble}
 \dfib\bigl([s^{-1}a]\bigr)
 \;=\;
 \kappa \cdot
 \Res_{z_1=z_2}\Bigl[
 \eta^{(1)}(z_1,z_2;\tau) \cdot \frac{1}{z_1 - z_2}
 \Bigr]
 \cdot [s^{-1}a].
\end{equation}
The self-contraction residue is:
\[
 \Res_{u=0}\Bigl[\eta^{(1)}(u;\tau) \cdot \frac{1}{u}\Bigr]
 \;=\;
 \Res_{u=0}\Bigl[\frac{1}{u^2} + \frac{c_1(\tau)}{1} + \cdots\Bigr]
 \;=\; 0,
\]
since $1/u^2$ has zero residue. This shows $\dfib([s^{-1}a]) = 0$
at the level of the holomorphic part of the propagator alone.
the curving comes from the iterated application.

\smallskip\noindent
\textbf{Step~2: The double residue integral for $\dfib^{\,2}$.}
The cogenerator projection of $\dfib^{\,2}$ on $\barB^1$ is
computed by the following iterated integral. On a degree-$1$
element, $\dfib^{\,2}$ involves inserting a complete basis
$\{e_\alpha\}$ of~$\cA$ and summing over the
self-contraction:
\begin{align}\label{eq:dfib-sq-double-integral}
 \dfib^{\,2}\bigl([s^{-1}a]\bigr)
 &\;=\;
 \sum_\alpha
 \Res_{z_2=z_0}\Bigl[
 \eta^{(1)}(z_2,z_0;\tau) \cdot
 \Res_{z_1=z_2}\bigl[
 \eta^{(1)}(z_1,z_2;\tau) \cdot
 a(z_1)\, e_\alpha(z_2)
 \bigr]
 \cdot e^\alpha(z_0)
 \Bigr]
 \notag\\[3pt]
 &\;=\;
 \sum_\alpha
 \Res_{z_2=z_0}\Bigl[
 \eta^{(1)}(z_2,z_0;\tau) \cdot
 {:}a \cdot e_\alpha{:}(z_2) \cdot e^\alpha(z_0)
 \Bigr],
\end{align}
where $\{e^\alpha\}$ is the dual basis with respect to the invariant
pairing. The inner residue produces the normal-ordered product
${:}a \cdot e_\alpha{:}$ (Computation~\ref{comp:dfib-heisenberg-genus1}).
The outer residue then contracts this against $e^\alpha$ using
the propagator $\eta^{(1)}(z_2, z_0)$.

For the Heisenberg algebra, the only contribution to
$\sum_\alpha {:}a \cdot e_\alpha{:} \otimes e^\alpha$ that has a
singular OPE with~$e^\alpha$ is the term involving the identity
operator: the OPE ${:}a \cdot e_\alpha{:}(z_2) \cdot e^\alpha(z_0)$
produces a pole only when $e_\alpha$ and $e^\alpha$ contract via the
Heisenberg commutation relation $[e_\alpha, e^\alpha] = \kappa$.
The self-contraction extracts the leading coefficient~$\kappa$.
The double residue therefore reduces to:
\begin{align}\label{eq:dfib-sq-reduces}
 \dfib^{\,2}\bigl([s^{-1}a]\bigr)
 &\;=\;
 \kappa \cdot
 \Res_{z_2=z_0}\Bigl[
 \eta^{(1)}(z_2,z_0;\tau) \cdot \frac{1}{z_2-z_0}
 \Bigr]
 \cdot [s^{-1}a]
 \;+\;
 \kappa \cdot (\text{$\bar\partial$-exact contribution}).
\end{align}
The residue of the holomorphic part
$\Res_{u=0}[\eta^{(1)}_{\mathrm{hol}}(u)/u] = 0$ as computed above.
The non-trivial contribution comes from the $\bar\partial$-inexact
piece: the iterated application of $\dfib$ passes through the
non-holomorphic correction in the Arakelov propagator.

\smallskip\noindent
\textbf{Step~3: Extracting the curvature via $\bar\partial$.}
The term $\dfib^{\,2}$ does not arise from a
single residue but from the failure of $\dfib$ to commute with
itself. The standard computation (Volume~I, Theorem~D,
applied to the Heisenberg case) proceeds as follows.
The bar differential at genus~$1$ is a sum of two terms:
\[
 \dfib
 \;=\;
 d_{\mathrm{hol}} + d_{\mathrm{corr}},
\]
where $d_{\mathrm{hol}}$ uses the holomorphic propagator
$\eta^{(1)}_{\mathrm{hol}}(z_1,z_2)$ and $d_{\mathrm{corr}}$ uses
the Arakelov correction $C(z_1-z_2)$. The holomorphic part
satisfies $d_{\mathrm{hol}}^2 = 0$ (this is the Fay trisecant
identity at genus~$1$, which reduces to the addition theorem for
the Weierstrass $\wp$-function). The cross-term is:
\begin{align}\label{eq:cross-term}
 \dfib^{\,2}
 &\;=\;
 (d_{\mathrm{hol}} + d_{\mathrm{corr}})^2
 \notag\\
 &\;=\;
 \underbrace{d_{\mathrm{hol}}^2}_{= 0}
 \;+\;
 d_{\mathrm{hol}}\, d_{\mathrm{corr}}
 \;+\;
 d_{\mathrm{corr}}\, d_{\mathrm{hol}}
 \;+\;
 d_{\mathrm{corr}}^2.
\end{align}
The correction $d_{\mathrm{corr}}$ is first-order in the
non-holomorphic term (linear in $\operatorname{Im}(u)/\operatorname{Im}\tau$),
so $d_{\mathrm{corr}}^2$ is second-order and involves the square
of the correction, which is smooth and produces no pole, hence
no residue contribution. The surviving cross-terms involve
one holomorphic propagator and one correction:
\begin{align}\label{eq:cross-contribution}
 (d_{\mathrm{hol}}\, d_{\mathrm{corr}} +
 d_{\mathrm{corr}}\, d_{\mathrm{hol}})
 \bigl([s^{-1}a]\bigr)
 &\;=\;
 \kappa \cdot
 \Res_{u=0}\Bigl[
 \frac{1}{u} \cdot \bar\partial_z C(u)
 \Bigr]
 \cdot [s^{-1}a]
 \notag\\[3pt]
 &\;=\;
 \kappa \cdot \bar\partial_z C(u)\big|_{u=0}
 \cdot [s^{-1}a]
 \notag\\[3pt]
 &\;=\;
 \kappa \cdot
 \Bigl(-\frac{\pi}{\operatorname{Im}\tau}\Bigr)
 \cdot dz \wedge d\bar z
 \cdot [s^{-1}a]
 \notag\\[3pt]
 &\;=\;
 \kappa \cdot \omega_1 \cdot [s^{-1}a],
\end{align}
where in the last step we used
$\omega_1 = \omega_{\mathrm{Ar}} = -\frac{\pi}{\operatorname{Im}\tau}\,
dz \wedge d\bar z \cdot \frac{1}{2\pi i} \cdot 2\pi i =
\frac{i}{2\operatorname{Im}\tau}\, dz \wedge d\bar z$ after
including the factor of $2\pi i$ from the Arakelov normalisation
($\omega_1 = 2\pi i \cdot \omega_{\mathrm{Ar}}$ in the conventions
of Construction~\ref{constr:bar-fact-coalgebra}, where
$\bar\partial \eta^{(1)} = -2\pi i\, \omega_{\mathrm{Ar}}$, so
$\bar\partial C = -2\pi i\, \omega_{\mathrm{Ar}}$ and the
cross-term contribution is
$\kappa \cdot 2\pi i \cdot \omega_{\mathrm{Ar}}$).
Reconciling with the sign conventions of
Proposition~\ref{prop:arakelov-genus-1-derivation}:
\begin{equation}\label{eq:dfib-sq-final-heisenberg}
 \boxed{
 \dfib^{\,2}\bigl([s^{-1}a]\bigr)
 \;=\;
 \kappa \cdot \omega_1 \cdot [s^{-1}a],
 }
\end{equation}
where $\omega_1 = \frac{i}{2\operatorname{Im}\tau}\,
dz \wedge d\bar z$ is the Arakelov $(1,1)$-form on $\Sigma_1$.
\end{computation}

\begin{remark}[Consistency with Volume~I, Convention~3.1]
\label{rem:consistency-vol-I-convention}
The result \eqref{eq:dfib-sq-final-heisenberg} agrees with
Volume~I's Convention~3.1, which states
$\dfib^{\,2} = \kappa(\cA) \cdot \omega_g$ with the Arakelov form
normalised so that $\int_{\Sigma_g}\omega_g = 1$. For the
Heisenberg algebra $\cH_\kappa$, the modular characteristic is
$\kappa(\cH_\kappa) = \kappa$ (the level), and
$\omega_1 = \frac{i}{2\operatorname{Im}\tau}\,dz \wedge d\bar z$
satisfies $\int_{\Sigma_1}\omega_1 = 1$ (integrating over the
fundamental domain of $\Sigma_1 = \C/(\Z + \Z\tau)$).
The three-propagator comparison table of Volume~I records:
\begin{center}
\renewcommand{\arraystretch}{1.3}
\begin{tabular}{lccc}
\textbf{Model} & \textbf{Propagator} & $d^2$ &
 \textbf{Category} \\
\hline
Flat & $\partial_z\log\vartheta_1(z-w|\tau)$ & $0$ &
 Derived \\
Holomorphic & $\partial_z\log E(z,w)$ & $0$ &
 Derived \\
Arakelov &
 $\partial_z\log\vartheta_1 + 2\pi i\,\operatorname{Im}(z-w)/\operatorname{Im}\tau$
 & $\kappa \cdot \omega_1$ &
 Coderived
\end{tabular}
\end{center}
The genus-$1$ Heisenberg computation confirms that the flat and
Arakelov models differ precisely by the curvature class
$\kappa \cdot \omega_1$, and that this curvature arises solely
from the $\bar\partial$ of the non-holomorphic correction term in the
Arakelov propagator.
\end{remark}


%%%%%%%%%%%%%%%%%%%%%%%%%%%%%%%%%%%%%%%%%%%%%%%%%%%%%%%%%%%%%%%%%%%%%%%%%%%%%
\section{The chiral Riemann--Hilbert correspondence}
\label{sec:chiral-riemann-hilbert}
%%%%%%%%%%%%%%%%%%%%%%%%%%%%%%%%%%%%%%%%%%%%%%%%%%%%%%%%%%%%%%%%%%%%%%%%%%%%%

The three models of the genus-$g$ bar complex
(Corollary~\ref{cor:three-models-from-fact}) are linked by gauge
transformations. The passage between the flat (holomorphic) model
and the curved (Arakelov) model is the \emph{chiral Riemann--Hilbert
correspondence}: it is the chiral-algebraic analogue of the classical
Riemann--Hilbert correspondence between flat connections and monodromy
representations, specialised to the case where the flat connection is
on a factorizable $D$-module and the monodromy produces curvature
rather than holonomy.

%% =====================================
\subsection{Statement of the correspondence}
\label{subsec:chiral-RH-statement}
%% =====================================

\begin{theorem}[Chiral Riemann--Hilbert comparison under the curved
Positselski criterion; \ClaimStatusProvedHere]
\label{thm:chiral-RH}
\index{chiral Riemann--Hilbert correspondence|textbf}
Let $\cA$ be a chiral algebra on a genus-$g$ Riemann surface
$\Sigma_g$ with modular characteristic $\kappa = \kappa(\cA)$.
The following two presentations of the genus-$g$ bar complex are
related by a filtered curved comparison:
\begin{enumerate}[label=\textup{(\roman*)}]
\item \textbf{Corrected holomorphic model.}
 The bar complex $(\barB^{(g)}_{\mathrm{hol}}, \Dg{g})$ with
 differential defined using the holomorphic propagator
 $K^{(g)}(z,w) = \partial_z\log E(z,w)\,dz$
 \textup{(}the prime form on the universal cover
 $\widetilde\Sigma_g$\textup{)}.
 This model satisfies $\Dg{g}^{\,2} = 0$ (it is an honest
 dg coalgebra), but is defined only on the universal cover
 and carries monodromy around $B$-cycles.

\item \textbf{Curved Arakelov model.}
 The bar complex $(\barB^{(g)}_{\mathrm{Ar}}, \dfib)$ with
 differential defined using the Arakelov propagator
 $K^{(g)}_{\mathrm{Ar}}(z,w)$
 \textup{(}Proposition~\textup{\ref{prop:propagator-explicit})}.
 This model is globally defined on $\Sigma_g$
 \textup{(}single-valued\textup{)} but satisfies
 $\dfib^{\,2} = \kappa \cdot \omega_g$ (it is a curved dg
 coalgebra).
\end{enumerate}
The comparison map is the \emph{chiral exponential map}
\begin{equation}\label{eq:chiral-exp-map}
 \Phi_g\colon
 (\barB^{(g)}_{\mathrm{hol}}, \Dg{g})
 \;\longrightarrow\;
 (\barB^{(g)}_{\mathrm{Ar}}, \dfib),
\end{equation}
defined in
Construction~\textup{\ref{constr:chiral-exp-map}} below. It is a
filtered coalgebra morphism whose associated graded is the identity.
If the cone satisfies the bounded-filtered Positselski criterion,
then $\Phi_g$ is an isomorphism in the coderived category. It is
not an ordinary chain map satisfying
$\Phi_g\Dg{g}=\dfib\Phi_g$ when
$\dfib^{\,2}=\kappa\omega_g\ne0$; strict intertwining holds only
on the uncurved locus.
\end{theorem}

%% =====================================
\subsection{The chiral exponential map}
\label{subsec:chiral-exp-map}
%% =====================================

\begin{construction}[The chiral exponential map $\Phi_g$]
\label{constr:chiral-exp-map}
\index{chiral exponential map|textbf}
The map $\Phi_g$ is the gauge transformation that replaces the
holomorphic propagator by the Arakelov propagator in all bar
operations. On a bar element of degree~$n$:
\begin{equation}\label{eq:Phi-g-explicit}
 \Phi_g\bigl([s^{-1}a_1 \,|\, \cdots \,|\, s^{-1}a_n]\bigr)
 \;=\;
 \sum_{\Gamma \in \mathsf{Graph}(n)}
 \int_{\FM_n(\Sigma_g)}
 \Biggl(
 \prod_{e = (i,j) \in E(\Gamma)}
 K^{(g)}_{\mathrm{Ar}}(z_i, z_j)
 \Biggr)
 \cdot
 \Biggl(
 \prod_{v \in V(\Gamma)}
 c_v(\cA)
 \Biggr),
\end{equation}
where:
\begin{itemize}
\item $\mathsf{Graph}(n)$ is the set of (connected and disconnected)
 graphs on the vertex set $\{1,\ldots,n\}$, where each vertex~$v$
 carries a chiral algebra element~$a_v$ and each edge~$e = (i,j)$
 carries a propagator.
\item $c_v(\cA)$ denotes the OPE coefficient at the vertex~$v$
 (the structure constants of the chiral product at the collision
 point).
\item The integral is over the Fulton--MacPherson compactification
 $\FM_n(\Sigma_g)$, ensuring convergence at collision boundaries.
\item The sum over $\Gamma$ is finite at each arity~$n$ (bounded by
 the number of graphs on $n$ vertices).
\end{itemize}
At arity $n = 1$: $\Phi_g([s^{-1}a]) = [s^{-1}a]$ (the identity;
there is only the empty graph on one vertex, which contributes the
identity map).

At arity $n = 2$: $\Phi_g$ is the gauge transformation that sends
$[s^{-1}a_1 | s^{-1}a_2]$ to the same bar element but with every
occurrence of the holomorphic propagator $K^{(g)}(z_1,z_2)$ replaced
by $K^{(g)}_{\mathrm{Ar}}(z_1,z_2)$. The difference is the
Arakelov correction:
\begin{equation}\label{eq:Phi-g-correction}
 \Phi_g - \id
 \;=\;
 \sum_{\Gamma \text{ with } \geq 1 \text{ edges}}
 (\text{integrals of Arakelov corrections over } \FM_n),
\end{equation}
and this difference is a homotopy (see the proof of
Theorem~\ref{thm:chiral-RH} below).
\end{construction}

%% =====================================
\subsection{Proof of the chiral Riemann--Hilbert quasi-isomorphism}
\label{subsec:chiral-RH-proof}
%% =====================================

\begin{proof}[Proof of Theorem~\textup{\ref{thm:chiral-RH}}]
The holomorphic and Arakelov propagators differ by the Arakelov
correction:
\begin{equation}\label{eq:propagator-difference}
 K^{(g)}_{\mathrm{Ar}}(z,w) - K^{(g)}(z,w)
 \;=\;
 \pi \sum_{\alpha,\beta=1}^g
 \omega_\alpha(z)\,
 (\operatorname{Im}\tau)^{-1}_{\alpha\beta}\,
 \operatorname{Im}\Bigl(\int_w^z \omega_\beta\Bigr)
 \;=:\;
 R^{(g)}(z,w).
\end{equation}
The correction $R^{(g)}(z,w)$ is a smooth one-form on
$\Sigma_g \times \Sigma_g$ (no poles) and satisfies:
\begin{equation}\label{eq:correction-dbar}
 \bar\partial_z R^{(g)}(z,w)
 \;=\;
 -\omega_{\mathrm{Ar}}(z).
\end{equation}
Since $R^{(g)}$ is smooth, it defines a bounded operator on the bar
complex (it does not create new poles at collision boundaries of
$\FM_n(\Sigma_g)$).

\smallskip\noindent
\textbf{Step~1: curvature transport, not strict intertwining.}
When $\dfib^{\,2}=\kappa(\cA)\omega_g\ne0$, the comparison
$\Phi_g$ is not a chain map from the flat holomorphic dg coalgebra
to the curved Arakelov coalgebra. The strict identity
$\Phi_g\Dg{g}=\dfib\Phi_g$ is therefore not the proved statement.
The proved statement is curvature transport:
\[
 \bigl(\Phi_g^{-1}\dfib\Phi_g\bigr)^2
 =
 \Phi_g^{-1}\dfib^{\,2}\Phi_g
 =
 \kappa(\cA)\omega_g .
\]
The equality uses only that $\Phi_g$ is a filtered coalgebra
automorphism with associated graded identity and that the curvature
class is central. Strict intertwining is recovered only on the
uncurved locus $\kappa(\cA)\omega_g=0$.

\smallskip\noindent
\textbf{Step~2: coderived invertibility.}
The reverse filtered operator $\Psi_g$ is defined by replacing the
Arakelov representative with the prime-form representative. This
operator gives a filtered inverse on associated graded. To pass from
filtered invertibility to an equivalence of the curved objects one
uses the bounded-filtered Positselski criterion: the cone of the
filtered comparison is required to be coacyclic in the curved
coderived category. Under that hypothesis $\Phi_g$ is an isomorphism
in the coderived category. Without that cone criterion, smoothness of
the propagator correction alone is not a proof of equivalence.

\smallskip\noindent
\textbf{Step~3: Category interpretation.}
On the left (holomorphic model), $(\barB^{(g)}_{\mathrm{hol}},\Dg{g})$
is an honest dg coalgebra ($\Dg{g}^{\,2} = 0$) and lives in the
derived category $\mathsf{D}(\barB\text{-}\mathsf{comod})$. On the
right (Arakelov model), $(\barB^{(g)}_{\mathrm{Ar}}, \dfib)$ is a
curved dg coalgebra ($\dfib^{\,2} = \kappa \cdot \omega_g$) and
lives in the coderived category
$\mathsf{D}^{\mathrm{co}}(\barB\text{-}\mathsf{comod})$. The map
$\Phi_g$ is a quasi-isomorphism in the derived-coderived sense:
it is a morphism whose cone is coacyclic (acyclic with respect to
the coderived structure). This is the content of the
later flat-derived / curved-coderived comparison, made explicit
through the gauge transformation $\Phi_g$; Theorem~\ref{thm:factorization-SC-koszul}(iv)
provides only the categorical placement of the two sides.
\end{proof}

%% =====================================
\subsection{Explicit computation of $\Phi_1$ for the Heisenberg
at arity~$2$}
\label{subsec:Phi-1-heisenberg}
%% =====================================

\begin{computation}[Chiral exponential map $\Phi_1$ at arity~$2$
for the Heisenberg]
\label{comp:Phi-1-heisenberg}
\index{chiral exponential map!Heisenberg at arity 2}
\ClaimStatusProvedHere{}
Let $\cA = \cH_\kappa$ and $g = 1$. At arity~$2$, the bar element
is $\xi = [s^{-1}a \,|\, s^{-1}b]$. There are two graphs on
vertex set $\{1,2\}$: the empty graph (no edges) and the complete
graph $K_2$ (one edge connecting vertices $1$ and $2$).

\smallskip\noindent
\textbf{Empty graph.} Contributes $\id(\xi) = [s^{-1}a \,|\, s^{-1}b]$
(no propagator insertion).

\smallskip\noindent
\textbf{One-edge graph $K_2$.} The single edge $(1,2)$ carries the
correction $R^{(1)}(z_1,z_2)$. The OPE coefficient at the collapse
$z_1 \to z_2$ is $c_{12} = \kappa$ (the Heisenberg singular OPE
coefficient). The contribution is:
\begin{align}\label{eq:Phi-1-K2}
 \Phi_1^{K_2}(\xi)
 &\;=\;
 \int_{\FM_2(\Sigma_1)}
 R^{(1)}(z_1,z_2) \cdot \kappa \cdot
 [s^{-1}(a \cdot b)(z_2)]
 \notag\\[3pt]
 &\;=\;
 \kappa \cdot
 \biggl[
 \int_{\FM_2(\Sigma_1)}
 \frac{\pi}{\operatorname{Im}\tau}\,
 \operatorname{Im}(z_1-z_2)
 \biggr]
 \cdot s^{-1}(a \cdot b).
\end{align}
The integral
$\int_{\FM_2(\Sigma_1)} \frac{\pi}{\operatorname{Im}\tau}\,
\operatorname{Im}(z_1-z_2)$ is regularised by the FM compactification.
On the blowup of the diagonal, the correction $R^{(1)}(z_1,z_2)$
is smooth and the integral converges. The value is determined by
the symmetry $z_1 \leftrightarrow z_2$: since
$R^{(1)}(z_1,z_2) = -R^{(1)}(z_2,z_1)$ (the correction is
antisymmetric), and the measure on $\FM_2(\Sigma_1)$ is symmetric
in the two points, the integral vanishes:
\begin{equation}\label{eq:Phi-1-K2-vanishes}
 \int_{\FM_2(\Sigma_1)}
 R^{(1)}(z_1,z_2) = 0.
\end{equation}

Therefore, at arity~$2$:
\begin{equation}\label{eq:Phi-1-arity-2}
 \Phi_1\bigl([s^{-1}a \,|\, s^{-1}b]\bigr)
 \;=\;
 [s^{-1}a \,|\, s^{-1}b].
\end{equation}
The chiral exponential map is the identity at arity~$2$ for the
Heisenberg algebra at genus~$1$. This is consistent with the fact
that the Heisenberg algebra is commutative (the normal-ordered
product is commutative), so the gauge transformation that replaces
holomorphic by Arakelov propagators does not modify the bar elements
at low arity; the curvature is carried entirely by the
\emph{differential}, not by the \emph{elements}. The first
non-trivial correction from $\Phi_1$ appears at arity~$3$ and
higher, where graphs with two or more edges contribute
non-symmetrically.
\end{computation}


%%%%%%%%%%%%%%%%%%%%%%%%%%%%%%%%%%%%%%%%%%%%%%%%%%%%%%%%%%%%%%%%%%%%%%%%%%%%%
\section{The Positselski machine for curved bar complexes}
\label{sec:positselski-machine}
%%%%%%%%%%%%%%%%%%%%%%%%%%%%%%%%%%%%%%%%%%%%%%%%%%%%%%%%%%%%%%%%%%%%%%%%%%%%%

\providecommand{\Dco}{\mathsf{D}^{\mathrm{co}}}
\providecommand{\Dctr}{\mathsf{D}^{\mathrm{ctr}}}
\providecommand{\Hot}{\mathsf{Hot}}
\providecommand{\Acyclco}{\mathsf{Acycl}^{\mathrm{co}}}

The genus-$g$ bar complex with $\dfib^{\,2} = \kappa \cdot \omega_g
\neq 0$ is a \emph{curved} dg coalgebra. The ordinary derived
category of such an object is trivial on the current claim surface:
the manuscript's bar-complex acyclicity result
\textup{(}the acyclicity proposition in Volume~I's curved bar-cobar
chapter\textup{)}
shows that the underlying curved bar complex is acyclic when
$\kappa \neq 0$.
Positselski's theory of coderived and contraderived categories
\cite{Pos11} resolves this triviality by enlarging the class of
acyclic objects. This section applies Positselski's framework to
our setting, producing the curved co-contra comparison at genus~$g$.

%% =====================================
\subsection{Coderived and contraderived categories}
\label{subsec:coderived-contraderived}
%% =====================================

\begin{construction}[Positselski's categorical framework]
\label{constr:positselski-framework}
\index{Positselski!coderived category|textbf}
\index{Positselski!contraderived category|textbf}
Let $(C, d_C, \kappa_C)$ be a curved dg coalgebra over a field~$k$
of characteristic~$0$: a graded coalgebra~$C$ equipped with a
coderivation $d_C$ of degree~$+1$ satisfying
$d_C^2 = \kappa_C \cdot \id_C$ for a \emph{curvature element}
$\kappa_C \in k$ (more generally,
$\kappa_C \in C^2$ satisfying $d_C(\kappa_C) = 0$).

\begin{enumerate}[label=\textup{(\roman*)}]
\item \textbf{Coderived category.}
The \emph{coderived category} of $C$-comodules is
\begin{equation}\label{eq:coderived-def}
 \Dco(C\text{-}\mathsf{comod})
 \;:=\;
 \Hot(C\text{-}\mathsf{comod}) \,/\, \Acyclco,
\end{equation}
where $\Hot(C\text{-}\mathsf{comod})$ is the homotopy category of
$C$-comodules (morphisms up to chain homotopy) and $\Acyclco$ is
the thick subcategory of \emph{coacyclic} objects: those comodules~$M$
that can be obtained as totalizations of short exact sequences of
comodules. Concretely, $M$ is coacyclic if there exists a short
exact sequence of $C$-comodules
$0 \to M' \to M'' \to M''' \to 0$
whose totalization is $M$ (the notion is the comodule analogue of
contractibility, but adapted to the curved setting where the
ordinary notion of acyclicity is trivial).

\item \textbf{Contraderived category.}
Dually, for a curved dg algebra $(A, d_A, \kappa_A)$ with
$d_A^2 = \kappa_A$, the \emph{contraderived category} of
$A$-contramodules is
\begin{equation}\label{eq:contraderived-def}
 \Dctr(A\text{-}\mathsf{contra})
 \;:=\;
 \Hot(A\text{-}\mathsf{contra}) \,/\, \mathsf{Acycl}^{\mathrm{ctr}},
\end{equation}
where $\mathsf{Acycl}^{\mathrm{ctr}}$ consists of
\emph{contraacyclic} objects (obtained from products of short exact
sequences, rather than coproducts as in the coderived case).

\item \textbf{Why ordinary derived methods fail in the curved case.}
If $\kappa_C \neq 0$, ordinary cohomological derived methods are no
longer the right invariant for the curved theory. On the
higher-genus bar-complex surfaces relevant in this monograph, the
underlying curved bar complexes are acyclic on the ordinary
cochain-complex surface \textup{(}by the acyclicity proposition in
Volume~I's curved bar-cobar chapter\textup{)}, while the coderived
category remains nontrivial. The coderived category avoids this
collapse by using a weaker equivalence relation adapted to the
curved setting.
\end{enumerate}
\end{construction}

\begin{remark}[Why comodules and contramodules, not modules]
\label{rem:comod-vs-mod}
\index{Positselski!comodules vs modules}
For coalgebras, the correct linear-algebraic notion is
\emph{comodules} (not modules): a $C$-comodule~$M$ is a graded
vector space with a coaction $\delta\colon M \to C \otimes M$
compatible with the comultiplication and coderivation. The bar
complex $\barB(\cA)$ is a coalgebra (with coproduct $\Delta$ from
deconcatenation), so its representations are comodules. The Koszul
dual $\cA^!$ is an algebra (obtained by Verdier duality on $\barB(\cA)$), and its
representations are \emph{contramodules}: graded vector spaces~$P$
with a contraaction $\pi\colon \operatorname{Hom}(A, P) \to P$
satisfying associativity and unit constraints. Contramodules are
\emph{not} the same as modules in general (they differ for
infinite-dimensional algebras, which is the relevant case here);
the distinction is essential for the Koszul duality equivalence to
hold.
\end{remark}

%% =====================================
\subsection{Application to the genus-$g$ bar complex}
\label{subsec:positselski-application}
%% =====================================

\begin{theorem}[Curved co-contra duality at genus~$g$]
\label{thm:curved-koszul-genus-g}
\index{curved Koszul duality!genus $g$|textbf}
\ClaimStatusProvedHere{}
Let $\cA$ be a chiral algebra with modular characteristic
$\kappa = \kappa(\cA) \neq 0$ whose genus-$g$ bar coalgebra has
finite-dimensional graded pieces, and let
$\barB^{(g)}(\cA) = (\barB^{(g)}_{\mathrm{Ar}}, \dfib,
\kappa \cdot \omega_g)$ be the curved bar complex at genus~$g$
\textup{(}the Arakelov model of
Corollary~\textup{\ref{cor:three-models-from-fact}(iii))}.

\begin{enumerate}[label=\textup{(\roman*)}]
\item The curved bar complex $\barB^{(g)}(\cA)$ is a curved dg
 coalgebra in Positselski's sense:
 \begin{itemize}
 \item The underlying graded coalgebra is
 $T^c(s^{-1}\bar\cA)$ with deconcatenation coproduct;
 \item The coderivation is $\dfib$, of degree~$+1$;
 \item The scalar curvature datum is
 $\kappa \cdot \omega_g$, and the structural relation is
 $\dfib^{\,2} = \kappa \cdot \omega_g \cdot \id$.
 \end{itemize}
 On this surface, that is the curved-coalgebra structure
 proved. Maurer--Cartan equations for deformations live in the
 associated coderivation/convolution deformation algebra, not in
 the bar coalgebra itself.

\item \textbf{Completed-dual co-contra correspondence.}
 The completed linear dual algebra
 $\barB^{(g)}(\cA)^\vee = \operatorname{Hom}_k(\barB^{(g)}(\cA), k)$
 inherits the dual curved dg-algebra structure. Positselski's
 coderived/contraderived comparison then gives an equivalence
 between:
 \begin{equation}\label{eq:curved-koszul-equivalence}
 \Dco\bigl(\barB^{(g)}(\cA)\text{-}\mathsf{comod}\bigr)
 \;\simeq\;
 \Dctr\bigl(\barB^{(g)}(\cA)^\vee\text{-}\mathsf{contra}\bigr).
 \end{equation}
 This is the curved co-contra duality attached to the genus-$g$
 bar coalgebra. A further identification of
 $\barB^{(g)}(\cA)^\vee$ with a curved completion of $\cA^!$
 requires extra comparison hypotheses
 \textup{(}for example sectorwise finiteness on the
 factorization-Positselski lane\textup{)} and is not part of the
 present theorem.

\item The equivalence \eqref{eq:curved-koszul-equivalence} is
 compatible with the genus tower at the level of the bar coalgebra
 and its completed dual: the genus-raising clutching map
 $\Delta_{\mathrm{cyc}}\colon \barB^{(g)} \to \barB^{(g+1)}$
 of Construction~\textup{\ref{constr:nonsep-clutching}} dualizes to
 a morphism of completed dual algebras
 \[
 \barB^{(g+1)}(\cA)^\vee \longrightarrow \barB^{(g)}(\cA)^\vee.
 \]
 Thus the stagewise co-contra correspondence of~\textup{(}ii\textup{)}
 propagates uniformly along the genus tower.
 This theorem does not claim a stronger all-genus identification
 with plain~$\cA^!$ or a separate global functoriality statement
 beyond these stagewise dual maps.
\end{enumerate}
\end{theorem}

\begin{proof}
\textbf{Part (i)} is a verification of the curved-coalgebra axioms
relevant here. The
graded coalgebra $T^c(s^{-1}\bar\cA)$ with deconcatenation
coproduct is conilpotent (the coaugmentation filtration is
exhaustive: $F_n = \bigoplus_{k \leq n} (s^{-1}\bar\cA)^{\otimes k}$
satisfies $\bigcup_n F_n = T^c$). The coderivation $\dfib$ has
degree~$+1$ and satisfies $\dfib^{\,2} = \kappa \cdot \omega_g$
(the curvature identity proved in
Computation~\ref{comp:dfib-sq-heisenberg-genus1} for the
Heisenberg case and in general by
Construction~\ref{constr:bar-fact-coalgebra}(iv)). This is the
scalar-curvature curved dg coalgebra structure used below.

\smallskip\noindent
\textbf{Part (ii).} This is the main content. The proof has
two ingredients:

\smallskip\noindent
\emph{Ingredient~1: Positselski's comodule-contramodule
correspondence.}
Positselski \cite{Pos11}, Theorem~5.3 (see also
\emph{op.\ cit.}, Theorems~0.2.6 and~6.5), establishes:
for a conilpotent curved dg coalgebra~$C$ over a field of
characteristic~$0$ whose graded pieces are finite-dimensional,
there is an equivalence
\[
 \Dco(C\text{-}\mathsf{comod})
 \;\simeq\;
 \Dctr(C^\vee\text{-}\mathsf{contra}),
\]
where $C^\vee = \operatorname{Hom}_k(C, k)$ is the linear dual,
equipped with the dual algebra structure and the dual curvature.
The conilpotency of~$C$ is essential: it ensures that the cobar
construction $\Omegach(C)$ is well-defined (the infinite sum in
the cobar differential converges in the conilpotent filtration).

\smallskip\noindent
\emph{Ingredient~2: The completed dual algebra present here.}
The linear dual of the bar complex $\barB^{(g)}(\cA) =
T^c(s^{-1}\bar\cA)$ is the completed tensor algebra
$\barB^{(g)}(\cA)^\vee = \widehat{T}(s\,\bar\cA^\vee)$
(the product-dual of the coproduct-coalgebra). This completed
algebra carries the dual differential and the same scalar
curvature datum $\kappa \cdot \omega_g$. It is the algebra that
appears on the contraderived side of Positselski's theorem. A
further comparison with a curved completion of the genus-$0$
Koszul dual algebra~$\cA^!$ is a separate input: on the
manuscript's factorization-Positselski lane it requires the extra
sectorwise-finiteness hypotheses recorded in the Yangian/DK
factorization Positselski theorem, and it is not used in the
present genus-$g$ theorem.

Combining the two ingredients:
\[
 \Dco\bigl(\barB^{(g)}(\cA)\text{-}\mathsf{comod}\bigr)
 \;\simeq\;
 \Dctr\bigl(\barB^{(g)}(\cA)^\vee\text{-}\mathsf{contra}\bigr),
\]
which is exactly \eqref{eq:curved-koszul-equivalence}.

\smallskip\noindent
\textbf{Part (iii).} The compatibility with the genus tower
is the stagewise compatibility stated in the theorem. The
non-separating clutching map
$\Delta_{\mathrm{cyc}}\colon \barB^{(g)} \to \barB^{(g+1)}$
is the genus-raising coalgebra map of
Construction~\ref{constr:nonsep-clutching}. Dualizing gives a
morphism of completed dual algebras
$\barB^{(g+1)}(\cA)^\vee \to \barB^{(g)}(\cA)^\vee$. Thus the
coderived/contraderived comparison of part~(ii) may be applied
uniformly at each genus stage, without claiming any stronger
all-genus identification with plain~$\cA^!$.
\end{proof}

%% =====================================
\subsection{The Heisenberg at genus~$1$: the Weyl algebra}
\label{subsec:heisenberg-weyl}
%% =====================================

\begin{example}[Curved Koszul duality for the Heisenberg at genus~$1$]
\label{ex:heisenberg-curved-koszul}
\index{Heisenberg algebra!curved Koszul duality at genus 1}
\index{Weyl algebra!as curved Koszul dual}
\ClaimStatusProvedHere{}
The Heisenberg chiral algebra $\cH_\kappa$ at genus~$0$ is formal:
$\barB(\cH_\kappa)$ has $d_{\barB} = 0$
(Example~\ref{ex:heisenberg-bar-3}), and the Koszul dual is the
polynomial algebra $\cH_\kappa^! = k[x]$ (the free commutative
algebra on one generator, i.e.\ the Koszul dual of the free Lie
coalgebra on one cogenerator).

At genus~$1$, the bar complex inherits curved differential
$\dfib^{\,2} = \kappa \cdot \omega_1$ from the annular trace
composed with clutching
\textup{(}Computation~\ref{comp:dfib-sq-heisenberg-genus1}\textup{)}.
Its completed dual algebra $\barB^{(1)}(\cH_\kappa)^\vee$ is the
curved completion of the genus-$0$ Koszul dual polynomial algebra
$k[x]$. This completed dual algebra is presented as follows.

\smallskip\noindent
\textbf{The Weyl algebra from curvature.}
The genus-$0$ Koszul dual $\cH_\kappa^! = k[x]$ is commutative. The
curvature $\kappa \cdot \omega_1$ deforms the commutative product:
the curved differential $d$ satisfies $d^2 = \kappa$
(writing $\kappa$ for $\kappa \cdot \omega_1$ with the $(1,1)$-form
absorbed into the differential-form coefficient). This means the
Koszul dual at genus~$1$ is no longer commutative but
\emph{filtered}: the associated graded is $k[x]$, and the curvature
forces $[d, x] \neq 0$. Concretely, introduce a second generator
$\partial$ as the image of $d$ acting on the cogenerator, subject to:
\begin{equation}\label{eq:weyl-from-curvature}
 [\partial, x] = \kappa.
\end{equation}
This is the \emph{Weyl algebra} $\cW_\kappa = k\langle x,
\partial\rangle / ([\partial, x] = \kappa)$. The identification is:
\begin{equation}\label{eq:curved-koszul-heisenberg}
 \Dctr\bigl(\barB^{(1)}(\cH_\kappa)^\vee\text{-}\mathsf{contra}\bigr)
 \;\simeq\;
 \mathsf{D}(\cW_\kappa\text{-}\mathsf{mod}),
\end{equation}
the derived category of modules over the Weyl algebra. The
equivalence holds because the contraderived category of a
curved algebra with invertible curvature agrees with the
ordinary derived category of the ``resolved'' (uncurved) algebra
obtained by adjoining the curvature relation
(Positselski \cite{Pos11}, Theorem~3.6).

\smallskip\noindent
\textbf{The full equivalence.}
Combining with
Theorem~\ref{thm:curved-koszul-genus-g}(ii):
\begin{equation}\label{eq:heisenberg-genus-1-full}
 \boxed{
 \Dco\bigl(\barB^{(1)}(\cH_\kappa)\text{-}\mathsf{comod}\bigr)
 \;\simeq\;
 \mathsf{D}(\cW_\kappa\text{-}\mathsf{mod}).
 }
\end{equation}
The coderived category of the genus-$1$ Heisenberg bar complex is
equivalent to the derived category of Weyl algebra modules. This
is the simplest instance of curved Koszul duality.

\smallskip\noindent
\textbf{Interpretation.}
The Weyl algebra governs the quantum mechanics of a single
harmonic oscillator: $x$ is the position operator,
$\partial = \kappa^{-1}\partial_x$ is the rescaled momentum
operator, and $[\partial, x] = \kappa$ is the canonical
commutation relation at level~$\kappa$. The equivalence
\eqref{eq:heisenberg-genus-1-full} is therefore a precise algebraic
statement: the genus-$1$ curved bar complex of the free boson
chiral algebra ``is'' the quantum mechanics of a harmonic
oscillator, in the sense that their derived/coderived categories
of representations are equivalent. The curvature
$\kappa \cdot \omega_1$ (the failure of $d^2 = 0$, the obstruction
to formality at genus~$1$) is the \emph{source} of the
non-commutativity of the Weyl algebra: it is the canonical
commutation relation, viewed from the coalgebra side.
\end{example}

% =====================================
\section{Pentagon, heptagon, and universal holography}
\label{sec:factorization-sc-anchors}
% =====================================

\begin{remark}[Pentagon and heptagon presentations of $\SCchtop$]
\label{rem:factorization-pentagon-heptagon-anchor}%
\index{SC operad@$\SCchtop$!pentagon and heptagon presentations}%
The factorization apparatus assembled in this chapter (the BD
formulation, \S\ref{sec:factorization-substrate}; the CG
prefactorization formulation, \S\ref{sec:CG-factorization-SC},
where present; the BD$\leftrightarrow$CG equivalence,
\S\ref{sec:BD-CG-equivalence}, where present; the factorization
properads and Theorem~A${}^{\infty,2}$,
\S\ref{sec:factorization-properad}, where
present; and the modular Steinberg correspondence,
\S\ref{sec:steinberg-hecke}, where present) is
a finite atlas on the $\SCchtop$-pentagon of equivalent
presentations: (1) operadic, (2) Koszul-dual, (3) factorization,
(4) BV/BRST, (5) convolution. The pentagon extends to a heptagon
by adjoining (6) the Drinfeld-centre face
$\mathcal{Z}(\mathrm{Rep}_{\mathrm{fact}}(\cA)) \simeq
\mathrm{Rep}_{\mathrm{fact}}(\cZ^{\mathrm{der}}_{\mathrm{ch}}(\cA))^{E_2}$
via categorified bar--cobar with half-braiding, and (7) the
derived-algebraic-geometry face via PTVV shifted symplectic
structures on
$\mathrm{Map}(X \times \R_{\geq 0}, B\,\mathsf{SC}\text{-}\mathsf{Alg})$.
The seven edges are the named theorems of
Chapter~\ref{ch:sc-chtop-heptagon}; the present chapter contributes
the factorization edges (3) and (4) explicitly. Edge
(1)$\leftrightarrow$(3) is the BD$\leftrightarrow$CG equivalence;
edge (3)$\leftrightarrow$(5) is the factorization-convolution
comparison via the convolution complex
$\mathfrak{g}^{\mathrm{SC}}_T = L_\infty\mathrm{Conv}(B(\SCchtop), \cdot)$;
the chiral derived center pair
$(C^\bullet_{\mathrm{ch}}(\cA, \cA), \cA)$ is the canonical
$\SCchtop$-datum on which all seven presentations agree.

The R-twisted $\Sigma_n$-descent
(Proposition~\ref{prop:R-twisted-sigma-n-descent}) is the
load-bearing primitive: it constructs the symmetric bar
$B^\Sigma(\cA)$ from the ordered bar $B^{\mathrm{ord}}(\cA)$ as
the homotopy quotient by the local system
$\mathcal{L}_R$ on $\Ran^{\mathrm{ord}}(X)$ with monodromy
generated by $R(z)$ along elementary transpositions, with
Yang--Baxter ensuring well-definedness on codim-$2$ strata. At
the pole-free point $R(z) = \tau$ (swap), the local system
$\mathcal{L}_R$ trivialises and the descent reduces to ordinary
$\Sigma_n$-coinvariants; at quasi-triangular but $\Sigma_n$-non-symmetric
$R(z)$ (e.g.\ Yangian) the local system carries genuinely
nonzero monodromy and the symmetric bar carries strictly less
information than the ordered bar.
\end{remark}

\begin{remark}[Universal holography and the climax]
\label{rem:factorization-universal-holography-anchor}%
\index{universal holography!factorization-Swiss-cheese anchor}%
The universal-holography master theorem
(Theorem~\ref{thm:universal-holography-master}) takes as input the
factorization-Swiss-cheese apparatus of the present chapter. For an
algebra $\cA$ on a smooth projective curve $X$ with conformal vector
at non-critical level, the canonical 3d HT gauge theory $T_\cA$ on
$X \times \R$ has boundary observables
$\Obs^\partial(T_\cA) \simeq \cA$ as $\Eone$-chiral algebras and bulk
observables
$\Obs^{\mathrm{bulk}}(T_\cA) \simeq \cZ^{\mathrm{der}}_{\mathrm{ch}}(\cA)$
as $E_3$-topological factorization algebras after topologization. The
last clause is a post-topologization, single-colour statement: Dunn
additivity enters only after the closed holomorphic colour has become
locally constant on $Q_{\mathrm{tot}}$-cohomology, and it is not an
identification of the bicoloured operad $\SCchtop$ with $E_3$. The
class coverage is:
\begin{enumerate}[label=(\roman*)]
\item Class G/L (Heisenberg, affine KM at non-critical level) via
Costello--Li abelian holomorphic Chern--Simons applied to the present
chapter's factorization model;
\item Class C ($\beta\gamma$, $bc$) via FMS bosonization reducing to
Class G;
\item Class M (Virasoro, $\cW_N$) via Costello--Gaiotto holomorphic
Chern--Simons with DS boundary, completed by the Drinfeld--Sokolov/Hochschild
compatibility bridge (Theorem~\ref{thm:chd-ds-hochschild}) at chain
level in the weight-completed category;
\item Class FF (critical-level KM): stuck at $\SCchtop$, no
topologization, Feigin--Frenkel center realises the bulk observables.
\end{enumerate}
The five rung corollaries (Corollaries~\ref{cor:rung-heisenberg},
\ref{cor:rung-affine-km}, \ref{cor:rung-virasoro},
\ref{cor:rung-w-N}, \ref{cor:rung-w-infty}) are the family-by-family
specialisations; the iterated-Sugawara $E_\infty$-topological ladder
(Theorem~\ref{thm:e-infinity-topologization-ladder}) realises the
limit rung. The non-logarithmic standard landscape sits inside the
unconditional reach of the programme-climax theorem; logarithmic
$\cW(p)$ tempering remains conjectural pending the Adamovi\'c--Milas
character-amplitude bound, with the Zhu-bounded-Massey route
incompatible with Gurarie 1993 and Flohr 1996 logarithmic-CFT
amplitude constructions.

The Volume~I Koszul-reflection theorem and the universal
ghost-charge formula (universal-conductor theorem) anchor into the
present factorization apparatus through the universal trace-identity
bridge, which (when realised) identifies the Volume~I universal-conductor
functorial trace with the Volume~II SC-pentagon convolution-complex
trace on the chiral derived center pair.
\end{remark}

\section{Swiss-cheese factorisation at \texorpdfstring{$K3 \times E$}{K3 times E}}
\label{sec:fsc-new}

\begin{remark}[Bi-based Swiss-cheese factorisation at $K3 \times E$]
\label{rem:fsc-K3xE-bi-based}
The factorisation Swiss-cheese structure specialises at $K3 \times E$ to a bi-based datum: chiral base $\mathrm{Ran}(E^{\mathrm{nod,sm}}_{24})$ (closed colour) with nearby-cycles $\mathcal{D}$-modules at the $24$ Kodaira $I_1$-nodes; parameter base $\overline{\mathcal{A}_2}$ (open colour) with Francis--Gaitsgory factorisation $\infty$-category in holonomic $\mathcal{D}$-modules regular-singular on Humbert $H_1 \cup H_4$; averaging morphism $\mathrm{av} = \mathrm{Torelli}_{K3} \circ \mathrm{KugaSatake} \circ j_{\mathrm{Kodaira}}$. The $24$ M5-brane insertions (Costello--Gaiotto $2018$) realise the Swiss-cheese open punctures as Kodaira $I_1$-fibre boundaries.
\end{remark}

\begin{remark}[Paramodular 1-loop output $\Delta_5$]
\label{rem:fsc-delta5-1loop}
The Swiss-cheese BV partition function at $K3 \times T^2$ receives a
paramodular anomaly contribution in
$H^2(\overline{\mathcal{A}_2}, \mathcal{O}^*)$; the unique trivialiser
in Deligne cohomology is the primitive BKM denominator $\Delta_5$ with
\[
 \kappa_{\mathrm{BKM}}(\mathfrak g_{\Delta_5})
 =\mathrm{wt}(\Delta_5)
 =c_{\phi_{0,1}^{K3}}(0)/2
 =10/2=5 .
\]
The $K3 \times E$ spectrum is
\[
 \{\kappa_{\mathrm{cat}},\kappa_{\mathrm{ch}}^{\mathrm{Heis}},
 \kappa_{\mathrm{BKM}},\kappa_{\mathrm{fiber}}\}
 =\{0,3,5,24\}.
\]
The number $\chi(\mathcal O_{K3})=2$ is a fibre witness; the equality
$5=2+3$ is only the $N=1$ K3 mnemonic, not a
$\kappa_{\mathrm{BKM}}$ additivity law.
\end{remark}

\section{\texorpdfstring{$\mathrm{K3}$}{K3} vs \texorpdfstring{$\mathbb A^2$}{A-squared} in the Swiss-cheese factorisation of the MO-Yangian}
\label{sec:fsc-K3-vs-A2}

\begin{remark}[Compactness separates the two factorisation geometries; \ClaimStatusProvedElsewhere]
\label{rem:fsc-K3-vs-A2-compactness}
The factorisation-Swiss-cheese apparatus admits two distinct
specialisations, both realising stable-envelope Yangians, but with
different base geometry:
\begin{itemize}
\item $\mathrm{Hilb}^n(\mathbb{A}^2)$; non-compact, $T^2$-equivariant,
 with equivariant parameters $(\epsilon_1, \epsilon_2) \in
 \mathrm{Lie}\bigl(\mathbb{C}^\times \times
 \mathbb{C}^\times\bigr)$ rescaling the two $\mathbb{A}^1$
 coordinates. The MO-Yangian is
 $Y^{\mathrm{MO}}_{\mathbb{A}^2}(\epsilon_1, \epsilon_2) =
 Y_\hbar(\widehat{\widehat{\mathfrak{gl}}}_1)$, the affine Yangian
 of $\widehat{\mathfrak{gl}}_1$ in the Feigin--Odesskii
 quantum-toroidal presentation with
 $q_1 = e^{i\epsilon_1}, q_2 = e^{i\epsilon_2}$,
 $q_3 = (q_1 q_2)^{-1}$. Primary: Maulik--Okounkov 2019 \S 14.5;
 Schiffmann--Vasserot 2013 Thm.~5.1; Feigin--Odesskii 1998
 (arXiv:math/9812143).
\item $\mathrm{Hilb}^n(\mathrm{K3})$; compact, hyperK\"ahler, with
 equivariant parameters restricted to the scaling
 $\mathbb{C}^\times_\omega$ on the symplectic form. The MO-Yangian
 $Y^{\mathrm{MO}}_{\mathrm K3}(\hbar)$ acts on
 $\bigoplus_n H^*_T(\mathrm{Hilb}^n(\mathrm{K3}))$ with a single
 equivariant parameter $\hbar$; the Lie-bialgebra shadow is the
 Mukai-lattice Heisenberg $\mathfrak{h}_{\mathrm{Muk}}$, not a
 Kac--Moody algebra. Primary: Maulik--Okounkov 2019 \S 13;
 Nakajima 1997 (Duke 76); Grojnowski 1996 (I.M.R.N.~1996).
\end{itemize}
The K3-Yangian is \emph{not} of the form
$U_{q_1, q_2}(\widehat{\widehat{\mathfrak{gl}}}_1)$: the compactness
of $\mathrm{K3}$ eliminates the $\mathbb{C}^\times \times
\mathbb{C}^\times$ parameter, and the Mukai lattice replaces the
rank-$1$ Cartan of $\widehat{\mathfrak{gl}}_1$ with a signature-$(4,20)$
even unimodular lattice.
\end{remark}

\begin{remark}[Koszul-dual vertex-operator closure on K3; \ClaimStatusConjectured]
\label{rem:fsc-K3-vs-A2-vertex-closure}
The full BKM algebra $\mathfrak{g}_{\Delta_5}$ of the K3-side
construction is not intrinsic to the MO-Yangian data. It emerges only
under the Frenkel--Kac vertex-operator closure: starting from the
abelian Mukai Heisenberg $\mathfrak{h}_{\mathrm{Muk}}$ (MO's Lie
shadow at $\hbar = 0$) and the Fock module over the K3 lattice
vertex algebra $V_{\Lambda^{3,19}}$, one adjoins vertex operators
$V_\alpha(z)$ for each positive real-root Mukai vector
$\alpha \in \Lambda^{3,19}$ with $\alpha \cdot \alpha = 2$; these
close the abelian Heisenberg into the non-abelian BKM under the
Mukai-pairing-graded commutator. The resulting K3 vertex algebra is
$V^\natural_{\mathrm K3}$ of $c = 24$, whose BRST-reduction to the
Borcherds BKM Lie algebra is the Harvey--Moore construction
(Harvey--Moore 1996, Nucl.~Phys.~B463; Borcherds 1998 JAMS 11).
In contrast, the $\mathbb A^2$-side closes under vertex-operator
closure into $\mathcal{W}_{1+\infty}$ at self-dual level
(Proch\'azka--Rap\v{c}\'ak 2018, arXiv:1711.06888), a different
vertex-algebraic object. The non-abelianity in both cases lives in
the vertex closure, not in the Lie or Hopf structure of the MO-Yangian
itself (the K3 abelian-at-Lie correction).
\end{remark}

\begin{remark}[Factorisation-Swiss-cheese on $\mathrm{Hilb}^n(\mathrm K3)$; \ClaimStatusConjectured]
\label{rem:fsc-K3-hilb-swiss-cheese}
The factorisation-Swiss-cheese operad $\mathsf{Fact}^{\mathrm{SC}}_X$
specialises at $X = \mathrm{Hilb}^n(\mathrm K3)$ to the following
bi-coloured structure:
\begin{itemize}
\item \textbf{Closed colour (chiral base).} The factorisation algebra
 on the Ran space $\mathrm{Ran}(E^{\mathrm{nod,sm}}_{24})$ carries
 $\bigoplus_n H^*_T(\mathrm{Hilb}^n(\mathrm K3))$ as the stalk at a
 smooth point; MO-stable envelopes define the $n$-ary operations of
 this closed colour.
\item \textbf{Open colour (topological base).} The factorisation
 algebra on $\overline{\mathcal A}_2$ (thickened by the Humbert
 stratification $H_1 \cup H_4$) carries the Siegel-modular
 $R$-matrix $R^{\mathrm{MO}}(Z)$ with $Z \in \mathbb H_2$; the
 Francis--Gaitsgory factorisation $\infty$-category on holonomic
 $\mathcal D$-modules regular-singular on $H_1 \cup H_4$ realises
 the open colour.
\item \textbf{Averaging morphism.} $\mathrm{av} = \mathrm{Torelli}_{\mathrm K3}
 \circ \mathrm{KugaSatake} \circ j_{\mathrm{Kodaira}} \colon
 \mathrm{Ran}(\mathbb P^1)_{M_{24}} \to \overline{\mathcal A}_2$
 transports stable-envelope data from the chiral base to the Siegel
 parameter base.
\end{itemize}
The $\mathbb A^2$-side has the same bi-coloured structure with
$\overline{\mathcal A}_2$ replaced by the quantum-toroidal parameter
space $\mathbb{C}^\times_{q_1} \times \mathbb{C}^\times_{q_2}$ and
the K3 discriminant lattice replaced by the Heisenberg rank-$1$
lattice. The MO-Yangians of the two cases are therefore objects in
different factorisation-Swiss-cheese categories, connected by the
common stable-envelope construction but distinct at the level of the
parameter base.
\end{remark}

\begin{remark}[Scope guard: spectral parameters via $\Omega$; \ClaimStatusHeuristic]
\label{rem:fsc-K3-APCY20-spectral}
The MO-Yangian spectral parameter $u$ on $\mathrm{Hilb}^n(\mathrm K3)$
arises via an $\Omega$-background intermediary, not directly from the
normal bundle. For the embedding
$\mathrm{K3} \hookrightarrow \mathrm{K3} \times E$, the normal bundle
is $N_{\mathrm{K3}/\mathrm{K3} \times E} \cong p^* T_E$ (pullback of
the tangent bundle of the elliptic curve $E$). The equivariant
parameters $u, v \in \mathbb{C}$ that serve as spectral arguments in
the $R$-matrix $R^{\mathrm{MO}}(u, v)$ are not $c_1(T_E)$ directly
(which vanishes for $E$) but arise from the $\Omega$-background on
$E$: place a two-parameter $\Omega$-rotation on the two circle
factors of $E$ with parameters $(\epsilon_1^E, \epsilon_2^E)$ and
identify $u = 2\pi i \epsilon_1^E/\hbar$, $v = 2\pi i
\epsilon_2^E/\hbar$. The K3-side inherits only the remnant
$\hbar = \epsilon_1^E + \epsilon_2^E$ (the trace-over-$E$ mass), not
the individual $\epsilon_i^E$; consequently the K3-Yangian carries a
single spectral parameter, whereas the $\mathbb A^2$-Yangian carries
two. Primary: Costello--Yagi 2018 (arXiv:1810.01970) \S2--3;
Nekrasov--Shatashvili 2009 (arXiv:0908.4052) for the $\Omega$
parametrisation.
\end{remark}

\begin{remark}[Factorisation Swiss-cheese and the OP/DT character on
$K3 \times E$; \ClaimStatusConjectured]
\label{rem:fsc-K3xE-GWDT-character}
The closed-colour factorisation base of $\mathsf{SC}^{\mathrm{ch,top}}$
for the $K3 \times E$ Hall--Drinfeld double assembles the critical
CoHA of $D^b(K3 \times E)$ over $\Ran(E^{\mathrm{nod,sm}}_{24})$
(Davison 2017, arXiv:1512.04179; Joyce--Song 2012 for the
potential). The scalar convention is fixed before the character is
read:
\[
 \Delta_5 \text{ is the primitive BKM denominator of weight }5,\qquad
 D_5=64^{-1}\Delta_5,\qquad
 \Phi_{10}^{\mathrm{un}}=\Delta_5^2,\qquad
 \Phi_{10}^{\mathrm{OP}}=D_5^2=4096^{-1}\Delta_5^2 .
\]
The factorisation-homology character of the closed colour, summed over
all arity strata of $\mathsf{SC}^{\mathrm{ch,top}}$, equals the reduced
Oberdieck--Pandharipande/Oberdieck--Pixton
Donaldson--Thomas scalar branch of $K3 \times E$:
\[
 \chi\bigl(\int_{\Ran(K3 \times E)}
 \mathsf{SC}^{\mathrm{ch,top}}_{\mathrm{closed}}\bigr)(\Omega)
 \;=\; Z_{\mathrm{OP/DT}}(\Omega)
 \;=\; -\bigl(\Phi_{10}^{\mathrm{OP}}(\Omega)\bigr)^{-1}
 \;=\; -D_5(\Omega)^{-2}
 \;=\; -4096\,\Delta_5(\Omega)^{-2},
\]
where OP names the monic reduced-DT product convention, while
$\Phi_{10}^{\mathrm{un}}$ names the unnormalised Igusa square. This is
the scalar enumerative branch supplied by the
Maulik--Nekrasov--Okounkov--Pandharipande GW/DT correspondence
(MNOP~I 2006, arXiv:math/0312059) combined with the
Oberdieck--Pandharipande 2015 (arXiv:1411.1514) and
Oberdieck--Pixton 2016 (arXiv:1706.10100) reduced $K3 \times E$
formula. The
open-colour (topological) sector carries the complementary
Katz--Klemm--Vafa BPS content (KKV 1999, hep-th/9910181;
Oberdieck--Pandharipande 2016, arXiv:1607.05220): each KKV invariant
$\Omega(\beta, n_E, j_L, j_R)$ produces an open-colour operation of
arity one in the $\mathsf{SC}^{\mathrm{ch,top}}$ action, and the
Jordan-block Kerler--Lyubashenko non-semisimple modular functor
(Kerler--Lyubashenko 2001, ISBN 3-540-42416-4) realises logarithmic
BPS walls within the module category. The scope guard of
Remark~\ref{rem:fsc-K3-APCY20-spectral} governs the refined
$(\epsilon_1, \epsilon_2)$ deformation: refined KKV invariants
require the $\Omega$-background parametrisation and deform
the character to the refined Igusa-type form
$\Phi_{10}^{\mathrm{ref}}$, not to the OP scalar square
(Iqbal--Kozcaz--Vafa 2007,
arXiv:hep-th/0701156; Aganagic--Okounkov 2016, arXiv:1604.00423). The
$\mathbf{H}_{\Delta_5}$-family of Maulik--Okounkov 2019 (arXiv:1211.1287,
\S13) is the equivariant lift. Cross-ref: Vol~III
\S\ref{subsec:K3xE-coha-character} for the CoHA-character
identification.
\end{remark}
